%-----------------------------------------------------------------------------------
%   PACKAGES
%-----------------------------------------------------------------------------------

\documentclass[a4paper,11pt,abstract=true]{scrartcl}

\usepackage[margin=.9in]{geometry}


\usepackage{mlmodern}


\usepackage{amsmath,amsthm,amssymb}
\usepackage{mathtools,mathrsfs}
\usepackage{centernot}
\usepackage{wrapfig}

\usepackage[dvipsnames]{xcolor} % moved up: must be loaded before tikz and tcolorbox

\usepackage{tikz,tikz-cd,asymptote}
\usetikzlibrary{positioning, fit, calc}

\usepackage[most]{tcolorbox}
\usepackage{enumerate}
\usepackage[a]{esvect}
\usepackage{hyperref} % loaded last among the packages to avoid clashes


\makeatletter
\def\@seccntformat#1{\textcolor{Red}{\hspace*{0.2em}$\textbf{\S}$\hspace{-0.3em} \csname the#1\endcsname}\hspace{0.5em}}
\makeatother

\usepackage[headsepline]{scrlayer-scrpage}
\clearpairofpagestyles{}
\ihead{Linear Algebra}
\ohead{\pagemark}
\setkomafont{pageheadfoot}{\normalfont}
\setkomafont{pagenumber}{\normalfont}
\pagestyle{scrheadings}

\setlength{\parindent}{0cm}

%-----------------------------------------------------------------------------------
%   COMMANDS
%-----------------------------------------------------------------------------------

\DeclarePairedDelimiter{\floor}{\lfloor}{\rfloor}

\newcommand{\nonadj}{\,\neg\textsf{adj}\,}
\newcommand{\adj}{\,\textsf{adj}\,}
\newcommand{\RR}{\mathbb{R}}
\newcommand{\ZZ}{\mathbb{Z}}
\newcommand{\NN}{\mathbb{N}}
\newcommand{\CC}{\mathbb{C}}
\newcommand{\QQ}{\mathbb{Q}}
\newcommand{\sgn}{\text{sgn\hspace{0.5mm}}}
\newcommand{\osc}{\text{osc\hspace{0.5mm}}}
\newcommand{\foralls}{\hspace{2mm}\forall}
\newcommand{\rimply}{\Rightarrow}
\newcommand{\limply}{\Leftarrow}
\newcommand{\underdash}{\rule{0.25cm}{0.1mm}\hspace{0.5mm}}
\newcommand{\undersmalldash}{\rule{0.18cm}{0.1mm}\hspace{0.3mm}}
\newcommand{\s}{\square}
\newcommand{\xr}{\xrightarrow}

\newcommand{\genstirlingI}[3]{%
  \genfrac{\lbrack}{\rbrack}{0pt}{#1}{#2}{#3}%
}
\newcommand{\genstirlingII}[3]{%
  \genfrac{\{}{\}}{0pt}{#1}{#2}{#3}%
}
\newcommand{\genlah}[3]{%
  \genfrac{\lfloor}{\rfloor}{0pt}{#1}{#2}{#3}%
}
\newcommand{\stirlingI}[2]{\genstirlingI{}{#1}{#2}}
\newcommand{\stirlingII}[2]{\genstirlingII{}{#1}{#2}}
\newcommand{\lah}[2]{\genlah{}{#1}{#2}}

%-----------------------------------------------------------------------------------
%   ENVIRONMENTS
%-----------------------------------------------------------------------------------

\newtheoremstyle{dotlessP}{}{}{}{}{\color{RedViolet}\bfseries}{.}{ }{}
\theoremstyle{dotlessP}
\newtheorem{problem}{Problem}
\newtheoremstyle{dotlessS}{}{}{}{}{\color{ForestGreen}\bfseries}{.}{ }{}
\theoremstyle{dotlessS}
\newtheorem*{solution}{Solution}

\theoremstyle{remark}
\newtheorem*{remark}{Remark}
\theoremstyle{definition}
\newtheorem*{theorem}{Theorem}
\newtheorem*{lemma}{Lemma}

\tcbset{
  base/.style={
    arc=0mm,
    bottomtitle=0.5mm,
    boxrule=0mm,
    colbacktitle=black!10!white,
    coltitle=black,
    fonttitle=\bfseries,
    left=2.5mm,
    leftrule=1mm,
    right=3.5mm,
    title={#1},
    toptitle=0.75mm,
  }
}

\definecolor{brandblue}{rgb}{0.34, 0.7, 1}

% Key theorems (titled, blue).
\newtcolorbox{mainbox}[1]{
  enhanced jigsaw,
  breakable,
  colframe=brandblue,
  base={#1}
}

\newtcolorbox{subbox}[1]{
  colframe=black!30!white,
  base={#1}
}

% Key definitions and notation.
\newtcolorbox{definition_box}{
enhanced jigsaw,
breakable,
pad at break*=1mm,
colback=lightgray!12!white,
colframe=RedViolet!45!black,
colbacktitle=Gray!38!white,
before={\vspace{4mm}},
after={\vspace{2mm}},
sharp corners
}

% Exercises and open-ended problems.
\newtcolorbox{problem_box}{
enhanced jigsaw,
breakable,
pad at break*=1mm,
colback=white!5!white,
colframe=gray!75!black,
colbacktitle=white!28!white,
sharp corners,
before={\vspace{6mm}},
after={\vspace{2mm}}
}

\newtcolorbox{problem_box_1}{
enhanced jigsaw,
breakable,
pad at break*=1mm,
colback=SeaGreen!5!white,
colframe=SeaGreen!75!black,
colbacktitle=Black!28!white,
sharp corners,
before={\vspace{6mm}},
after={\vspace{2mm}}
}

\newtcolorbox{problem_box_title}[1]{
enhanced jigsaw,
breakable,
pad at break*=1mm,
colback=Apricot!5!white,
colframe=Melon!75!black,
sharp corners,
before={\vspace{2mm}},
after={\vspace{2mm}},
title={#1},
fonttitle=\bfseries
}

% Propositions, lemmas, proofs, and general text (unframed).
\newtcolorbox{claim_box}{
enhanced jigsaw,
breakable,
pad at break*=1mm,
colback=white,
colframe=white,
borderline west={2pt}{0pt}{white},
sharp corners,
before={\vspace{1mm}},
after={\vspace{0.8mm}}
}

% Examples.
\newtcolorbox{example_box}{
enhanced jigsaw,
breakable,
pad at break*=1mm,
colback=lightgray!12!white,
colframe=lightgray!15!white,
borderline west={2pt}{0pt}{black},
sharp corners,
before={\vspace{2mm}},
after={\vspace{2mm}}
}

% Cautions and remarks.
\newtcolorbox{remark_box}{
enhanced jigsaw,
breakable,
pad at break*=1mm,
colback=Apricot!5!white,
colframe=Melon!75!black,
sharp corners,
before={\vspace{2mm}},
after={\vspace{2mm}}
}


\begin{document}

%-----------------------------------------------------------------------------------
%   TOP MATTERS
%-----------------------------------------------------------------------------------


\title{Linear Algebra}
\subtitle{Handkerchief}
\author{Pratyay Sen}
\date{\today}

\maketitle

\tableofcontents

\clearpage


%-----------------------------------------------------------------------------------
%   MAIN BODY

\section{Vector Spaces}

\subsection{Fields}

\textbf{Definition:} A \emph{field} is a set $\mathbb{F}$ equipped with two operations,
$+$ and $*$
\[
+:\mathbb{F}\times \mathbb{F}\longrightarrow \mathbb{F}
\]
\[
*:\mathbb{F}\times \mathbb{F}\longrightarrow \mathbb{F}
\]
such that $0\neq1$ and, for all $a,b,c\in \mathbb{F}$, the following axioms hold:
\begin{enumerate}
    \item (Closure under $+$) $a+b\in \mathbb{F}$;
    \item (Closure under $*$) $a*b\in \mathbb{F}$;
    \item (Associativity of $+$) $(a+b)+c=a+(b+c)$;
    \item (Associativity of $*$) $(a*b)*c=a*(b*c)$;
    \item (Commutativity of $+$) $a+b=b+a$;
    \item (Commutativity of $*$) $a*b=b*a$;
    \item (Additive identity) there exists $0\in \mathbb{F}$ such that
    $a+0=a$;
    \item (Multiplicative identity) there exists $1\in \mathbb{F}$ such that
    $a*1=a$;
    \item (Additive inverses) for every $a\in \mathbb{F}$, there exists $-a\in \mathbb{F}$
    such that $a+(-a)=0$;
    \item (Multiplicative inverses) for every $a\in \mathbb{F}$ with $a\neq0$,
    there exists $a^{-1}\in \mathbb{F}$ such that $a*a^{-1}=1$;
    \item (Distributivity) $a*(b+c)=a*b+a*c$.
\end{enumerate}

\begin{remark_box}
\textbf{A field is a substructure of a ring:} A field is a commutative
ring with identity in which every nonzero element has an inverse over the operation $*$.
For example, over addition and multiplication, $\mathbb{Z}$ is a ring but not a field, since $2$ has no
multiplicative inverse in $\mathbb{Z}$, whereas $\mathbb{Q}$, $\mathbb{R}$,
and $\mathbb{C}$ are fields.
\end{remark_box}

\subsection{Vector Spaces}

\textbf{Definition:} Let $\mathbb{F}$ be a field. A \emph{vector space over $\mathbb{F}$} is a
set $V$ equipped with two binary operations $+$ and $*$
\[
+:V\times V\longrightarrow V
\]
\[
*:\mathbb{F}\times V\longrightarrow V
\]
such that, for all $u,v,w\in V$ and $a,b\in \mathbb{F}$:
\begin{enumerate}
    \item (Commutativity of $+$) $u+v=v+u$;
    \item (Associativity of $+$) $(u+v)+w=u+(v+w)$;
    \item (Identity over $+$) there exists $0_V\in V$ such that $v+0_V=v$;
    \item (Inverse over $+$) for each $v\in V$, there exists $-v\in V$ such that $v+(-v)=0_V$;
    \item (Distributivity of scalar $*$ over vector $+$) $a*(u+v)=a*u+a*v$;
    \item (Distributivity of scalar $+$ over scalar $*$ on vectors) $(a+b)*v=a*v+b*v$;
    \item (Associativity of $*$) $a*(b*v)=(a*b)*v$;
    \item (Identity over $*$) $1*v=v$, where $1$ is the identity of $\mathbb{F}$ over $*$.
\end{enumerate}

\textbf{Note:} As one might have noticed by now, I am deliberately avoiding using the term ``addition'' for $+$ and ``multiplication'' for $*$, because there are other functions that satisfy the field axioms and vector space axioms, that are not addition and multiplication. One such example is given below:

\begin{example_box}
\textbf{Example of nonstandard operations:} On the set $\mathbb{R}$, define
\[
x\oplus y=x+y+1
\qquad\text{and}\qquad
x\odot y=xy+x+y,
\]
where the operations on the right-hand sides are ordinary addition and
multiplication. These new operations are not the ordinary ones; for example,
$1\oplus1=3$ and $1\odot1=3$.

The set $\mathbb{R}$ is a field under $\oplus$ and $\odot$. Its identity over $\oplus$ is $-1$, its  identity over $\odot$ is $0$, the inverse over $\oplus$ of
$x$ is $-x-2$, and the inverse over $\odot$ of $x\neq-1$ is
\[
-\frac{x}{x+1}.
\]
To see why all the field axioms hold, consider the bijection
\[
T:\mathbb{R}\longrightarrow\mathbb{R},\qquad T(x)=x+1.
\]
It satisfies
\[
T(x\oplus y)=T(x)+T(y)
\qquad\text{and}\qquad
T(x\odot y)=T(x)T(y).
\]
Thus, $T$ converts the new operations into the ordinary field operations on
$\mathbb{R}$, so the usual field axioms are inherited.

This field is also a vector space over itself: take $V=\mathbb{R}$, use
$\oplus$ for vector addition, and define scalar multiplication by
\[
a*v=a\odot v=av+a+v.
\]
The vector-space axioms then follow from the field axioms. This example shows
that the symbols used for vector addition and scalar multiplication represent
abstract operations and need not be the usual arithmetic operations. However, now that the reader has understood that the ``addition'' and ``multiplication'' being talked about here need not necessarily be the conventional addition and multiplication, I shall use those terms from now on.
\end{example_box}

\textbf{Definition:} If $\mathbb{F}$ is a field and $n$ is a positive integer, then
\[
\mathbb{F}^n=\underbrace{\mathbb{F}\times \mathbb{F}\times\cdots\times \mathbb{F}}_{n\text{ times}}
=\{(x_1,x_2,\ldots,x_n):x_i\in \mathbb{F}\}.
\]
Thus, $\mathbb{F}^n$ is the set of all ordered $n$-tuples of elements of $\mathbb{F}$. $+$
and $*$ are defined componentwise:
\[
(x_1,\ldots,x_n)+(y_1,\ldots,y_n)
=(x_1+y_1,\ldots,x_n+y_n),
\]
\[
a*(x_1,\ldots,x_n)=(a*x_1,\ldots,a*x_n), \qquad a\in \mathbb{F}.
\]
With these operations, $\mathbb{F}^n$ is a vector space over $\mathbb{F}$.

\textbf{Definition:} The notation $\mathbb{F}^\infty$ denotes the set of all infinite
sequences of elements of $\mathbb{F}$:
\[
\mathbb{F}^\infty=\{(x_1,x_2,x_3,\ldots):x_i\in \mathbb{F}\text{ for every }i\in\mathbb{N}\}.
\]
$+$ and $*$ are defined componentwise:
\[
(x_1,x_2,\ldots)+(y_1,y_2,\ldots)
=(x_1+y_1,x_2+y_2,\ldots),
\]
\[
a*(x_1,x_2,\ldots)=(a*x_1,a*x_2,\ldots), \qquad a\in \mathbb{F}.
\]
With these operations, $\mathbb{F}^\infty$ is a vector space over $\mathbb{F}$. Here, sequences
may have infinitely many nonzero entries.

\subsubsection{Subspaces and affine spaces}
\textbf{Definition:} Let $V$ be a vector space over a field $\mathbb{F}$.
A subset $U\subseteq V$ is a \emph{subspace of $V$} if $U$ is itself a vector
space over $\mathbb{F}$ under the addition and scalar multiplication inherited from
$V$. Equivalently, $U$ is a subspace if:
\begin{enumerate}
    \item $0_V\in U$;
    \item $u+v\in U$ for every $u,v\in U$;
    \item $a*u\in U$ for every $a\in \mathbb{F}$ and $u\in U$.
\end{enumerate}
Equivalently, a nonempty subset $U\subseteq V$ is a subspace if
\[
a*u+b*v\in U
\]
for every $u,v\in U$ and $a,b\in \mathbb{F}$.

\begin{definition_box}
\textbf{Definition:} A subset
$W\subseteq \mathbb{F}^n$ is an \emph{$\mathbb{F}$-linear subspace of $\mathbb{F}^n$} if it contains the
zero vector and is closed under vector addition and multiplication by scalars
from $\mathbb{F}$. Explicitly:
\begin{enumerate}
    \item $(0,\ldots,0)\in W$;
    \item if $u,v\in W$, then $u+v\in W$;
    \item if $a\in \mathbb{F}$ and $u\in W$, then $a*u\in W$.
\end{enumerate}
Equivalently, a nonempty subset $W\subseteq \mathbb{F}^n$ is an $\mathbb{F}$-linear subspace if
\[
a*u+b*v\in W
\]
for every $u,v\in W$ and $a,b\in \mathbb{F}$. Thus, $W$ is closed under all linear
combinations of its elements. This is the special case $V=\mathbb{F}^n$ of the general
definition of a subspace of a vector space $V$.
\end{definition_box}

\begin{example_box}
\textbf{Trivial subspace:} The set containing only the zero
vector,
\[
\{0_V\},
\]
is a subspace of every vector space $V$ and is called the \emph{trivial
subspace} or \emph{zero subspace}. It is closed under vector addition and
scalar multiplication because
\[
0_V+0_V=0_V
\qquad\text{and}\qquad
a*0_V=0_V
\]
for every $a\in \mathbb{F}$. The subspaces $\{0_V\}$ and $V$ itself are sometimes
collectively called the trivial subspaces of $V$; all other subspaces are
called nontrivial proper subspaces.
\end{example_box}

\textbf{Examples:}
\begin{enumerate}
    \item Let $\mathbb{F}=\mathbb{R}$, $V=\mathbb{R}^3$, and
    \[
    U=\{(x,y,z)\in\mathbb{R}^3:x+y+z=0\}.
    \]
    Then $V$ is a vector space over $\mathbb{R}$, and $U$ is a subspace of
    $V$.

    \item Let $\mathbb{F}=\mathbb{C}$, $V=\mathbb{C}^2$, and
    \[
    U=\{(z,iz):z\in\mathbb{C}\}.
    \]
    Then $V$ is a vector space over $\mathbb{C}$, and $U$ is a subspace of
    $V$. Indeed, sums of vectors of the form $(z,iz)$ and complex scalar
    multiples of such vectors have the same form.

    \item Let $\mathbb{F}=\mathbb{R}$, let $V=\mathcal{P}_3(\mathbb{R})$ be the vector
    space of real polynomials of degree at most $3$, and let
    \[
    U=\{p\in\mathcal{P}_3(\mathbb{R}):p(0)=0\}.
    \]
    Then $U$ is a subspace of $V$, since the zero polynomial belongs to $U$
    and, whenever $p(0)=q(0)=0$, we have
    $(p+q)(0)=0$ and $(a*p)(0)=0$ for every $a\in\mathbb{R}$.
\end{enumerate}

\begin{definition_box}
\textbf{Definition:} Let $U$ be a subspace of $V$
and let $v\in V$. The \emph{translate} of $U$ by $v$ is
\[
v+U=\{v+u:u\in U\}.
\]
\end{definition_box}

\textbf{Definition:} A nonempty subset $A\subseteq V$ is an
\emph{affine subspace} of $V$ if there exist a linear subspace $U\subseteq V$
and a vector $v\in V$ such that
\[
A=v+U.
\]
The subspace $U$ is called the \emph{direction subspace} of $A$. Although the
vector $v$ is not unique, the direction subspace $U$ is unique. An affine
subspace looks like a shifted copy of $U$, but it need not contain $0_V$ and
therefore need not be a linear subspace. In fact,
\[
v+U\text{ is a subspace of }V
\quad\Longleftrightarrow\quad
v\in U,
\]
in which case $v+U=U$. Basically, 
every translate $A=v+U$ is an affine subspace with direction subspace $U$.  It is a linear
subspace exactly when $v\in U$, in which case $A=U$.

\textbf{Definition:} Let $U$ be a subspace of $V$ and $v\in V$. The
set
\[
v+U=\{v+u:u\in U\}
\]
is called the \emph{coset of $U$ represented by $v$}. Thus, in a vector
space, cosets of linear subspaces are precisely affine subspaces. Two vectors
may represent the same coset; specifically,
\[
v+U=w+U
\quad\Longleftrightarrow\quad
v-w\in U.
\]
Consequently, any two cosets of the same subspace are either equal or
disjoint (prove it yourself!).

\textbf{Definition:} Let $U$ be a subspace of a vector
space $V$ over $\mathbb{F}$. The \emph{quotient space of $V$ by $U$} is
\[
V/U=\{v+U:v\in V\},
\]
the set of all cosets of $U$ in $V$, equipped with the operations
\[
(v+U)+(w+U)=(v+w)+U,
\qquad
a*(v+U)=(a*v)+U,\quad a\in\mathbb{F}.
\]
Thus the elements of $V/U$ are entire cosets, not individual vectors of $V$.
Two vectors represent the same element of $V/U$ exactly when their
difference belongs to $U$.

These operations are \emph{well-defined}: their results do not depend on
which representatives we choose. Indeed, suppose $v+U=v'+U$ and
$w+U=w'+U$. Then $v-v'\in U$ and $w-w'\in U$. Since $U$ is a subspace,
\[
(v+w)-(v'+w')=(v-v')+(w-w')\in U,
\]
and, for every $a\in\mathbb{F}$,
\[
a*v-a*v'=a*(v-v')\in U.
\]
Hence replacing representatives leaves the resulting cosets unchanged.
The vector-space axioms follow from those of $V$. In particular, the zero
vector of $V/U$ is the coset $0_V+U=U$, and the additive inverse of $v+U$
is $(-v)+U$.

\begin{example_box}
\textbf{Example:} Take
\[
V=\mathbb{R}^2,
\qquad U=\{(t,0):t\in\mathbb{R}\}.
\]
For any $(a,b)\in V$,
\[
(a,b)+U=\{(t,b):t\in\mathbb{R}\}.
\]
Thus each element of $V/U$ is a horizontal line, determined entirely by its
height $b$. Adding two cosets adds their heights, and multiplying a coset by
a scalar multiplies its height. For example,
\[
\bigl((0,2)+U\bigr)+\bigl((0,3)+U\bigr)=(0,5)+U.
\]
The zero element of the quotient is the entire horizontal axis $U$.
\end{example_box}

\textbf{Note:} Please don't get discouraged upon seeing the number of definitions and propositions and theorems, the definitions are mostly just different ways of naming different stuff, and the propositions are just the immediate observable results. I am not saying that all propositions are trivial, but most of them are. 

\subsubsection{Sums and intersections}
\begin{definition_box}
\textbf{Sum of Subspaces} Let $U_1,\ldots,U_m$ be subspaces of
a vector space $V$ over $\mathbb{F}$. Their \emph{sum} is
\[
U_1+\cdots+U_m
=\{u_1+\cdots+u_m:u_i\in U_i\text{ for }i=1,\ldots,m\}.
\]
Thus, the sum consists of all vectors that can be written as a sum of one
vector from each subspace. In particular,
\[
U+W=\{u+w:u\in U,\ w\in W\}.
\]
\end{definition_box}

\begin{problem_box}
\textbf{Proposition:} The sum $U_1+\cdots+U_m$ is a subspace of $V$.

\end{problem_box}

\begin{claim_box}
\textit{Proof:} Because each $U_i$ is a subspace, $0_V\in U_i$; hence
\[
0_V=0_V+\cdots+0_V\in U_1+\cdots+U_m.
\]
Now let
\[
x=u_1+\cdots+u_m
\quad\text{and}\quad
y=v_1+\cdots+v_m
\]
belong to $U_1+\cdots+U_m$, where $u_i,v_i\in U_i$. Then
\[
x+y=(u_1+v_1)+\cdots+(u_m+v_m).
\]
Since each $U_i$ is closed under addition, $u_i+v_i\in U_i$, and therefore
$x+y\in U_1+\cdots+U_m$. Finally, for any $a\in \mathbb{F}$,
\[
a*x=(a*u_1)+\cdots+(a*u_m).
\]
Since each $U_i$ is closed under scalar multiplication, $a*u_i\in U_i$, so
$a*x\in U_1+\cdots+U_m$. Thus the sum contains $0_V$ and is closed under
addition and scalar multiplication; hence it is a subspace of $V$.
\hfill$\square$
\end{claim_box}

\begin{problem_box}
\textbf{Proposition:} The sum $U_1+\cdots+U_m$ is the smallest subspace of
$V$ containing $U_1,\ldots,U_m$.

\end{problem_box}

\begin{claim_box}
\textit{Proof:} First, the sum contains each $U_j$. Indeed, if $u\in U_j$,
then
\[
u=0_V+\cdots+0_V+u+0_V+\cdots+0_V\in U_1+\cdots+U_m,
\]
where $u$ occupies the $j$th position and every other term is the zero vector
of the corresponding subspace. Hence
\[
U_j\subseteq U_1+\cdots+U_m
\qquad\text{for every }j=1,\ldots,m.
\]

Now let $W$ be any subspace of $V$ containing every $U_j$. If
$x\in U_1+\cdots+U_m$, then
\[
x=u_1+\cdots+u_m
\]
for some $u_j\in U_j$. Since $U_j\subseteq W$, every $u_j$ belongs to $W$.
Because $W$ is closed under addition, $x\in W$. Therefore
\[
U_1+\cdots+U_m\subseteq W.
\]
Thus, $U_1+\cdots+U_m$ contains all the subspaces $U_j$ and is contained in
every subspace that contains them. It is therefore the smallest such subspace.
\hfill$\square$
\end{claim_box}

\begin{problem_box}
\textbf{Proposition:} If $U$ and $W$ are
subspaces of $V$, then $U\cap W$ is a subspace of $V$.
\end{problem_box}

\begin{claim_box}
\textit{Proof:} Because $U$ and $W$ are subspaces, both contain the zero
vector. Hence
\[
0_V\in U\cap W.
\]
Now let $u,v\in U\cap W$. Then $u,v\in U$ and $u,v\in W$. Since both
subspaces are closed under addition,
\[
u+v\in U
\qquad\text{and}\qquad
u+v\in W,
\]
so $u+v\in U\cap W$. Finally, if $a\in \mathbb{F}$, then closure of $U$ and $W$ under
scalar multiplication gives
\[
a*u\in U
\qquad\text{and}\qquad
a*u\in W.
\]
Therefore $a*u\in U\cap W$. Thus $U\cap W$ contains $0_V$ and is closed
under addition and scalar multiplication, so it is a subspace of $V$.
\hfill$\square$
\end{claim_box}

\begin{problem_box}
\textbf{Proposition:} If $U$ and $W$ are subspaces
of $V$, then $U\cup W$ is a subspace of $V$ if and only if
\[
U\subseteq W
\qquad\text{or}\qquad
W\subseteq U.
\]

\end{problem_box}

\begin{claim_box}
\textit{Proof:} If $U\subseteq W$, then $U\cup W=W$, which is a subspace.
Similarly, if $W\subseteq U$, then $U\cup W=U$, which is a subspace.

Conversely, suppose that $U\cup W$ is a subspace. Assume neither subspace is
contained in the other. Then there exist
\[
u\in U\setminus W
\qquad\text{and}\qquad
w\in W\setminus U.
\]
Because $U\cup W$ is assumed to be a subspace and $u,w\in U\cup W$, it must
contain $u+w$. Thus $u+w$ belongs to either $U$ or $W$. If $u+w\in U$, then
\[
w=(u+w)-u\in U,
\]
contradicting $w\notin U$. If $u+w\in W$, then
\[
u=(u+w)-w\in W,
\]
contradicting $u\notin W$. Hence one of $U$ and $W$ must be contained in the
other. \hfill$\square$
\end{claim_box}

\subsubsection{Spanning and linear independence}
\begin{definition_box}
\textbf{Definition:} Let $v_1,\ldots,v_m\in V$. The \emph{span} of
$v_1,\ldots,v_m$ is the set of all their linear combinations:
\[
\operatorname{span}(v_1,\ldots,v_m)
=\{a_1*v_1+\cdots+a_m*v_m:a_1,\ldots,a_m\in \mathbb{F}\}.
\]
More generally, if $S\subseteq V$, then
\[
\operatorname{span}(S)
=\left\{a_1*v_1+\cdots+a_m*v_m:
m\in\mathbb{N},\ v_i\in S,\ a_i\in \mathbb{F}\right\}.
\]
\end{definition_box}

Thus, $\operatorname{span}(S)$ is the smallest subspace of $V$ containing
$S$. By convention, $\operatorname{span}(\varnothing)=\{0_V\}$.

\begin{definition_box}
\textbf{Definition:} A list of vectors
$v_1,\ldots,v_m\in V$ is \emph{linearly independent} over $\mathbb{F}$ if no vector
in the list belongs to the span of the others; that is,
\[
v_j\notin\operatorname{span}(v_1,\ldots,v_{j-1},v_{j+1},\ldots,v_m)
\]
for every $j=1,\ldots,m$. Equivalently, the list is linearly independent if
\[
a_1*v_1+\cdots+a_m*v_m=0_V
\]
implies
\[
a_1=\cdots=a_m=0.
\]
In other words, the zero vector can be written as a linear combination of
$v_1,\ldots,v_m$ only by taking every coefficient to be zero. Another
equivalent span criterion is
\[
v_j\notin\operatorname{span}(v_1,\ldots,v_{j-1})
\qquad\text{for every }j=1,\ldots,m,
\]
where $\operatorname{span}(\varnothing)=\{0_V\}$.
\end{definition_box}

\begin{definition_box}
\textbf{Definition:} A list of vectors
$v_1,\ldots,v_m\in V$ is \emph{linearly dependent} over $\mathbb{F}$ if it is not
linearly independent; that is, if there exist scalars
$a_1,\ldots,a_m\in \mathbb{F}$, not all zero, such that
\[
a_1*v_1+\cdots+a_m*v_m=0_V.
\]
Such an equation is called a \emph{nontrivial linear relation}. Equivalently,
a finite list is linearly dependent if one of its vectors lies in the span of
the others.
\end{definition_box}


\textbf{Definition:} A subset $S\subseteq V$ is called a
\emph{generating set} or \emph{spanning set} for $V$ if
\[
\operatorname{span}(S)=V.
\]
In this case, we say that $S$ \emph{generates} or \emph{spans} $V$. This
means that every vector $v\in V$ can be written as a finite linear
combination of elements of $S$.

More generally, if $S\subseteq V$, then $S$ generates the subspace
$\operatorname{span}(S)$, even when it does not generate all of $V$. A
generating set need not be linearly independent; a basis is precisely a
linearly independent generating set.



\textbf{Definition:} A vector space $V$
over $\mathbb{F}$ is \emph{finitely generated} if it has a finite generating set. In
other words, there exist vectors $v_1,\ldots,v_n\in V$ such that
\[
V=\operatorname{span}(v_1,\ldots,v_n).
\]
For vector spaces, being finitely generated is equivalent to being
finite-dimensional: from any finite generating set, redundant vectors can be
removed until a finite basis remains. Conversely, every finite basis is a
finite generating set.


\subsubsection{Row and column spaces}
\textbf{Row Space} Let $A\in \mathbb{F}^{m\times n}$ be a matrix whose
rows are $r_1,\ldots,r_m\in \mathbb{F}^n$. The \emph{row space} of $A$ is the subspace
of $\mathbb{F}^n$ spanned by its rows:
\[
\operatorname{Row}(A)
=\operatorname{span}(r_1,\ldots,r_m)
=\{a_1*r_1+\cdots+a_m*r_m:a_1,\ldots,a_m\in \mathbb{F}\}.
\]

\textbf{Column Space} Let the columns of
$A\in \mathbb{F}^{m\times n}$ be $c_1,\ldots,c_n\in \mathbb{F}^m$. The \emph{column space} of
$A$ is the subspace of $\mathbb{F}^m$ spanned by its columns:
\[
\operatorname{Col}(A)
=\operatorname{span}(c_1,\ldots,c_n)
=\{a_1*c_1+\cdots+a_n*c_n:a_1,\ldots,a_n\in \mathbb{F}\}.
\]
Thus, for an $m\times n$ matrix, the row space is contained in $\mathbb{F}^n$, whereas
the column space is contained in $\mathbb{F}^m$.

\subsubsection{Bases and dimension}
\begin{definition_box}
\textbf{Definition:} A list of vectors
$v_1,\ldots,v_m\in V$ is a \emph{basis} of $V$ if:
\begin{enumerate}
    \item $v_1,\ldots,v_m$ are linearly independent; and
    \item $\operatorname{span}(v_1,\ldots,v_m)=V$.
\end{enumerate}
Equivalently, $v_1,\ldots,v_m$ form a basis of $V$ if every vector $v\in V$
can be written uniquely as
\[
v=a_1*v_1+\cdots+a_m*v_m
\]
for scalars $a_1,\ldots,a_m\in \mathbb{F}$. The scalars $a_1,\ldots,a_m$ are called
the \emph{coordinates} of $v$ with respect to this basis.
\end{definition_box}

\begin{problem_box}
\textbf{Theorem:} If a vector space $V$ over $\mathbb{F}$ is generated by $n$
vectors, then $V$ has a basis containing at most $n$ vectors.

\end{problem_box}

\begin{claim_box}
\textit{Proof:} Suppose
\[
V=\operatorname{span}(v_1,\ldots,v_n).
\]
If $(v_1,\ldots,v_n)$ is linearly independent, then it is already a basis of
$V$, and the result follows.

If the list is linearly dependent, then one of its vectors belongs to the
span of the others. After relabeling if necessary, suppose
\[
v_n\in\operatorname{span}(v_1,\ldots,v_{n-1}).
\]
Removing $v_n$ does not change the span, so
\[
V=\operatorname{span}(v_1,\ldots,v_{n-1}).
\]
If this shorter list is still linearly dependent, remove another vector that
lies in the span of the remaining vectors. Each removal preserves the span
and decreases the length of the list by one.

Because the original list is finite, the process eventually stops with a
linearly independent list $(u_1,\ldots,u_m)$ that still spans $V$. Hence this
list is a basis of $V$, and because it was obtained by deleting vectors from
the original list,
\[
m\leq n.
\]
Thus $V$ has a basis of at most $n$ vectors. \hfill$\square$
\end{claim_box}

\textbf{Definition:} Let
$B=(v_1,\ldots,v_n)$ be a basis of $V$. For each $v\in V$, the unique
expression
\[
v=a_1*v_1+\cdots+a_n*v_n,
\qquad a_1,\ldots,a_n\in \mathbb{F},
\]
is called the \emph{basis expansion of $v$ with respect to $B$}. The scalars
$a_1,\ldots,a_n$ are the coordinates of $v$, and the tuple
\[
{[v]}_B=(a_1,\ldots,a_n)\in \mathbb{F}^n
\]
is called the \emph{coordinate vector of $v$ relative to $B$}.

If $B={(v_i)}_{i\in I}$ is an infinite basis, then the basis expansion of $v$
has the form
\[
v=\sum_{i\in I}a_i*v_i,
\]
where only finitely many coefficients $a_i$ are nonzero. Thus every basis
expansion is a finite linear combination, even when the basis is infinite.

\textbf{Definition:} The \emph{dimension} of a vector space $V$
over $\mathbb{F}$, denoted by $\dim_{\mathbb{F}} V$ or simply $\dim V$, is the number of vectors
in any basis of $V$. All bases of a vector space have the same number of
elements, so the dimension is well-defined. If
\[
(v_1,\ldots,v_n)
\]
is a basis of $V$, then $\dim_{\mathbb{F}} V=n$. If no finite basis exists, then $V$ is
called \emph{infinite-dimensional}. The zero vector space has the empty basis,
and hence
\[
\dim_{\mathbb{F}}\{0_V\}=0.
\]

\begin{problem_box}
\textbf{All Finite Bases have the Same Length}
If
$(v_1,\ldots,v_n)$ and $(w_1,\ldots,w_m)$ are two finite bases of $V$, then
$m=n$.
\end{problem_box}

\begin{claim_box}
\textit{Proof:} We first show that $m\leq n$. Because
$(v_1,\ldots,v_n)$ spans $V$, we can write
\[
w_1=a_1*v_1+\cdots+a_n*v_n.
\]
At least one coefficient is nonzero. After relabeling the $v_i$ if necessary,
suppose $a_1\neq0$. We can then solve for $v_1$:
\[
v_1=a_1^{-1}*w_1-a_1^{-1}*a_2*v_2-\cdots-a_1^{-1}*a_n*v_n.
\]
Therefore
\[
(w_1,v_2,\ldots,v_n)
\]
still spans $V$.

Now write $w_2$ as a linear combination of this new spanning list. The
coefficient of at least one of $v_2,\ldots,v_n$ must be nonzero; otherwise
$w_2$ would belong to $\operatorname{span}(w_1)$, contradicting the linear
independence of $(w_1,w_2)$. Relabeling the remaining $v_i$ if necessary, we
can solve for $v_2$ and replace it by $w_2$. Hence
\[
(w_1,w_2,v_3,\ldots,v_n)
\]
still spans $V$.

Continuing in this way, after $k$ steps we obtain a spanning list
\[
(w_1,\ldots,w_k,v_{k+1},\ldots,v_n).
\]
This process cannot continue beyond $k=n$. If $m>n$, then after replacing all
the $v_i$, the list $(w_1,\ldots,w_n)$ would span $V$. Consequently,
$w_{n+1}$ would lie in $\operatorname{span}(w_1,\ldots,w_n)$, contradicting
the linear independence of $(w_1,\ldots,w_m)$. Thus $m\leq n$.

Interchanging the roles of the two bases gives $n\leq m$. Therefore
$m=n$. \hfill$\square$
\end{claim_box}

\begin{problem_box}
\textbf{Size and Extension of Linearly Independent Sets} Let $V$
be a vector space over $\mathbb{F}$ with $\dim_{\mathbb{F}} V=n$. If $L\subseteq V$ is linearly
independent, then
\[
|L|\leq n.
\]
Equality holds if and only if $L$ is a basis of $V$. If $|L|<n$, then $L$
can be extended to a basis of $V$.
\end{problem_box}

\begin{claim_box}
\textit{Proof:} Fix a basis $B=(b_1,\ldots,b_n)$ of $V$. The replacement
argument used above shows that an independent list cannot contain more
vectors than a spanning list. Since $B$ spans $V$, a linearly independent set
$L$ cannot contain $n+1$ elements. Hence $L$ is finite and $|L|\leq n$.

Write $L=\{u_1,\ldots,u_m\}$, where $m\leq n$. Beginning with the spanning
list $(b_1,\ldots,b_n)$, write $u_1$ as a linear combination of its vectors.
At least one coefficient is nonzero, so we can solve for the corresponding
$b_i$ and replace it by $u_1$ without changing the span.

Suppose that $u_1,\ldots,u_{k-1}$ have already replaced $k-1$ vectors of
$B$. Express $u_k$ as a linear combination of the resulting spanning list.
The coefficient of at least one remaining vector from $B$ must be nonzero;
otherwise $u_k$ would lie in
$\operatorname{span}(u_1,\ldots,u_{k-1})$, contradicting the linear
independence of $L$. We may therefore solve for that basis vector and replace
it by $u_k$ while preserving the span. Continuing in this way gives a
spanning list of the form
\[
(u_1,\ldots,u_m,c_{m+1},\ldots,c_n),
\]
where the vectors $c_{m+1},\ldots,c_n$ were chosen from $B$. This list has
$n$ elements. It must be linearly independent: otherwise one of its vectors
could be removed while preserving the span, producing a spanning list of
$n-1$ vectors. This is impossible because the replacement argument shows
that the independent list $B$, which has $n$ vectors, cannot be longer than
a spanning list. Thus the displayed list is a basis extending $L$.

If $m=n$, there are no additional vectors $c_i$, so $L$ itself is a basis.
Conversely, if $L$ is a basis, then the theorem that all finite bases have the
same length gives $|L|=n$. If $m<n$, the vectors
$c_{m+1},\ldots,c_n$ extend $L$ to a basis. \hfill$\square$
\end{claim_box}

\begin{problem_box}
\textbf{Bases of Subspaces}
Every subspace of a
finite-dimensional vector space has a finite basis.
\end{problem_box}

\begin{claim_box}
\textit{Proof:} Let $U$ be a subspace of a finite-dimensional vector space
$V$, and write $\dim_{\mathbb{F}} V=n$. Every linearly independent subset of $U$ is also
a linearly independent subset of $V$, so the preceding theorem shows that it
has at most $n$ elements. Therefore, among all linearly independent subsets
of $U$, choose one, say $L$, having the greatest possible number of elements.

We claim that $L$ spans $U$. If it did not, there would exist
\[
u\in U\setminus\operatorname{span}(L).
\]
Then $L\cup\{u\}$ would be linearly independent: in any relation involving
$u$, a nonzero coefficient of $u$ would imply
$u\in\operatorname{span}(L)$, while a zero coefficient would reduce the
relation to one among the elements of $L$. Thus all coefficients must be
zero. But $L\cup\{u\}$ has more elements than $L$, contradicting the maximal
choice of $L$. Hence $L$ spans $U$. Since $L$ is both linearly independent
and spanning, it is a basis of $U$. \hfill$\square$
\end{claim_box}

\begin{problem_box}
\textbf{Size and Reduction of Spanning Sets}
Let $V$ be a vector
space over $\mathbb{F}$ with $\dim_{\mathbb{F}} V=n$. If $S\subseteq V$ spans $V$, then
\[
|S|\geq n.
\]
Equality holds if and only if $S$ is a basis of $V$. If $|S|>n$, then $S$
contains a subset that is a basis of $V$.
\end{problem_box}

\begin{claim_box}
\textit{Proof:} Let $B=(b_1,\ldots,b_n)$ be a basis of $V$. Because $S$
spans $V$, each $b_i$ is a finite linear combination of elements of $S$.
For each $i$, choose a finite subset $S_i\subseteq S$ such that
\[
b_i\in\operatorname{span}(S_i),
\]
and set
\[
S_0=S_1\cup\cdots\cup S_n.
\]
The set $S_0$ is finite and spans every element of $B$; hence it spans $V$.
By repeatedly deleting any element of $S_0$ that lies in the span of the
others, we obtain a linearly independent spanning subset $B'\subseteq S_0$.
Thus $B'$ is a basis of $V$, and
\[
B'\subseteq S_0\subseteq S.
\]
Because every basis of $V$ has $n$ elements, $|B'|=n$, and therefore
$|S|\geq n$.

If $|S|=n$, then the inclusion $B'\subseteq S$ and the equality
$|B'|=|S|=n$ imply $B'=S$, so $S$ is a basis. Conversely, every basis spans
$V$ and has exactly $n$ elements. Finally, if $|S|>n$, the subset
$B'\subseteq S$ constructed above is a basis of $V$. \hfill$\square$
\end{claim_box}

\begin{definition_box}
\textbf{Definition:} For $j=1,\ldots,n$, let
\[
e_j=(0,\ldots,0,1,0,\ldots,0),
\]
where the $1$ occurs in the $j$th position. The list
\[
(e_1,\ldots,e_n)
\]
is called the \emph{standard basis} of $\mathbb{F}^n$. Every
$x=(x_1,\ldots,x_n)\in \mathbb{F}^n$ has the unique representation
\[
x=x_1*e_1+\cdots+x_n*e_n,
\]
and therefore $\dim_{\mathbb{F}} \mathbb{F}^n=n$.
\end{definition_box}

\begin{example_box}
\textbf{Other standard bases.}
\begin{enumerate}
    \item The standard basis of the polynomial space
    $\mathcal{P}_n(\mathbb{F})$ is
    \[
    (1,x,x^2,\ldots,x^n),
    \]
    so $\dim_{\mathbb{F}}\mathcal{P}_n(\mathbb{F})=n+1$.

    \item For the matrix space $\mathbb{F}^{m\times n}$, let $E_{ij}$ be the matrix
    whose $(i,j)$ entry is $1$ and whose other entries are $0$. The matrices
    \[
    E_{11},E_{12},\ldots,E_{mn}
    \]
    form its standard basis, so $\dim_{\mathbb{F}} \mathbb{F}^{m\times n}=mn$.
\end{enumerate}
\end{example_box}

\begin{problem_box}
\textbf{Proposition:} Any three polynomials in $\mathcal{P}_1(\mathbb{F})$ are
linearly dependent.

\end{problem_box}

\begin{claim_box}
\textit{Proof:} The list $(1,x)$ is a basis of $\mathcal{P}_1(\mathbb{F})$, so
\[
\dim_{\mathbb{F}}\mathcal{P}_1(\mathbb{F})=2.
\]
Suppose $p_1,p_2,p_3\in\mathcal{P}_1(\mathbb{F})$. If $p_1$ and $p_2$ are linearly
dependent, then $(p_1,p_2,p_3)$ is also linearly dependent. Otherwise,
$(p_1,p_2)$ is a linearly independent list of two vectors in a
two-dimensional space and therefore forms a basis. Hence
\[
p_3=a*p_1+b*p_2
\]
for some $a,b\in \mathbb{F}$. Rearranging gives
\[
a*p_1+b*p_2-1*p_3=0.
\]
Because the coefficient of $p_3$ is $-1\neq0$, this is a nontrivial linear
relation. Thus $(p_1,p_2,p_3)$ is linearly dependent. \hfill$\square$
\end{claim_box}

\begin{problem_box}
\textbf{Proposition:} Let $n\geq0$ be an
integer. Any set of $n+2$ polynomials in $\mathcal{P}_n(\mathbb{F})$, the
space of polynomials in $x$ of degree at most $n$ together with the zero
polynomial, is linearly dependent over $\mathbb{F}$.

\end{problem_box}

\begin{claim_box}
\textit{Proof:} The monomials
\[
1,x,x^2,\ldots,x^n
\]
form a basis of $\mathcal{P}_n(\mathbb{F})$, so this space has dimension
$n+1$. By the theorem on the size of linearly independent sets, every
linearly independent subset of this space has at most $n+1$ elements.
Since $n+2>n+1$, a set $\{p_1,\ldots,p_{n+2}\}$ of $n+2$ polynomials
cannot be linearly independent. Thus there exist scalars
$a_1,\ldots,a_{n+2}\in\mathbb{F}$, not all zero, such that
\[
a_1*p_1+\cdots+a_{n+2}*p_{n+2}=0,
\]
where $0$ denotes the zero polynomial. Hence the set is linearly dependent.
\hfill$\square$
\end{claim_box}

\begin{problem_box}
\textbf{Dimension of a Quotient Space}
Let $V$ be a
finite-dimensional vector space over $\mathbb{F}$ and let $U$ be a subspace
of $V$. Then
\[
\dim(V/U)+\dim U=\dim V.
\]
\end{problem_box}

\begin{claim_box}
\textit{Proof:} Write $\dim U=k$ and $\dim V=n$. Choose a basis
$u_1,\ldots,u_k$ of $U$. By the basis-extension theorem,
extend it to a basis
\[
u_1,\ldots,u_k,v_1,\ldots,v_{n-k}
\]
of $V$. We claim that the cosets
\[
v_1+U,\ldots,v_{n-k}+U
\]
form a basis of $V/U$.

First, they span $V/U$. Every $v\in V$ has an expression
\[
v=\sum_{i=1}^{k}a_i u_i+\sum_{j=1}^{n-k}b_j v_j,
\qquad a_i,b_j\in\mathbb{F}.
\]
The first sum belongs to $U$, so in the quotient it contributes the zero
coset. Hence
\[
v+U=\sum_{j=1}^{n-k}b_j*(v_j+U).
\]

Next, suppose a linear combination of these cosets is zero in $V/U$:
\[
\sum_{j=1}^{n-k}c_j*(v_j+U)=U.
\]
This means $\sum_{j=1}^{n-k}c_j v_j\in U$, so there are scalars
$d_1,\ldots,d_k\in\mathbb{F}$ such that
\[
\sum_{j=1}^{n-k}c_j v_j-\sum_{i=1}^{k}d_i u_i=0_V.
\]
Since the extended basis of $V$ is linearly independent, all coefficients
are zero. In particular, every $c_j=0$, proving independence in $V/U$.

Thus $\dim(V/U)=n-k$, and therefore
\[
\dim(V/U)+\dim U=(n-k)+k=n=\dim V.
\]
Empty lists and empty sums are allowed here, so the proof also covers
$U=\{0_V\}$ and $U=V$. \hfill$\square$
\end{claim_box}

\subsubsection{Direct sums and complements}
\textbf{Internal Direct Sum} Let $U_1,\ldots,U_m$ be subspaces
of $V$. Their sum is called an \emph{internal direct sum} if every vector in
$U_1+\cdots+U_m$ has a unique representation of the form
\[
u_1+\cdots+u_m,
\qquad u_i\in U_i.
\]
In this case, we write
\[
U_1\oplus\cdots\oplus U_m
\]
instead of $U_1+\cdots+U_m$. Thus,
\[
V=U_1\oplus\cdots\oplus U_m
\]
means that every $v\in V$ can be written uniquely as
$v=u_1+\cdots+u_m$ with $u_i\in U_i$.

Equivalently, the sum is direct if
\[
u_1+\cdots+u_m=0_V,
\qquad u_i\in U_i,
\]
implies $u_1=\cdots=u_m=0_V$. Another equivalent condition is
\[
U_j\cap\sum_{i\neq j}U_i=\{0_V\}
\qquad\text{for every }j.
\]
For two subspaces, this simplifies to
\[
U+W=U\oplus W
\quad\Longleftrightarrow\quad
U\cap W=\{0_V\}.
\]

\textbf{Complementary Subspace} Let $U$ be a subspace of $V$.
A subspace $W\subseteq V$ is called a \emph{complementary subspace} or
\emph{complement} of $U$ in $V$ if
\[
V=U\oplus W.
\]
Equivalently, $W$ is a complement of $U$ if
\[
U+W=V
\qquad\text{and}\qquad
U\cap W=\{0_V\}.
\]
Thus every $v\in V$ has a unique decomposition
\[
v=u+w,
\qquad u\in U,\quad w\in W.
\]
A complementary subspace need not be unique. If $V$ is finite-dimensional,
then every subspace has a complement, and any complement $W$ satisfies
\[
\dim V=\dim U+\dim W.
\]

\begin{problem_box}
\textbf{Theorem:} Let $V$ be
finite-dimensional and let $U$ be a subspace of $V$. Then there exists a
subspace $W\subseteq V$ such that
\[
V=U\oplus W.
\]

\end{problem_box}

\begin{claim_box}
\textit{Proof:} Choose a basis $B_U$ of $U$. By the basis-extension theorem,
$B_U$ can be extended to a basis of $V$. Thus there is a set $C\subseteq V$
such that
\[
B_U\cup C
\]
is a basis of $V$. Define
\[
W=\operatorname{span}(C).
\]

Because $B_U\cup C$ spans $V$, every $v\in V$ can be written as a finite
linear combination of elements of $B_U\cup C$. Grouping the terms from
$B_U$ and the terms from $C$ gives
\[
v=u+w
\]
for some $u\in U$ and $w\in W$. Hence $U+W=V$.

It remains to show that $U\cap W=\{0_V\}$. Let $x\in U\cap W$. Because
$x\in U$, it is a linear combination of elements of $B_U$; because
$x\in W$, it is a linear combination of elements of $C$. Subtracting these
two expressions gives a linear combination of elements of $B_U\cup C$ equal
to $0_V$. Since $B_U\cup C$ is linearly independent, every coefficient must
be zero. Therefore $x=0_V$, and thus
\[
U\cap W=\{0_V\}.
\]
Consequently, $V=U\oplus W$, so $W$ is a complementary subspace of $U$.
\hfill$\square$
\end{claim_box}

\begin{remark_box}
\textbf{Is every sum a direct sum?} Nope, this statement is not always true. Every direct sum is a sum, but a sum
is direct only when the decomposition of each vector is unique. For two
subspaces, the sum is direct precisely when $U\cap W=\{0_V\}$.
\end{remark_box}

\begin{example_box}
\textbf{Example:} In $V=\mathbb{R}^2$, let
\[
U=\operatorname{span}((1,0))
\qquad\text{and}\qquad
W=\mathbb{R}^2.
\]
Then
\[
U+W=\mathbb{R}^2,
\]
but the sum is not direct because $U\cap W=U\neq\{(0,0)\}$. Equivalently,
the vector $(1,0)$ has two different decompositions:
\[
(1,0)=(1,0)+(0,0)=(0,0)+(1,0),
\]
where the first term in each sum belongs to $U$ and the second belongs to
$W$.
\end{example_box}

\textbf{External Direct Sum} Let ${(V_i)}_{i\in I}$ be a family
of vector spaces over the same field $\mathbb{F}$. Their \emph{external direct sum} is
the vector space
\[
\bigoplus_{i\in I}V_i
=\left\{{(v_i)}_{i\in I}:v_i\in V_i
\text{ and }v_i=0\text{ for all but finitely many }i\right\},
\]
with componentwise addition and scalar multiplication. Its elements are
tuples, not sums taken inside a common ambient vector space.

\textbf{Internal versus External Direct Sum} Let
\[
U=\operatorname{span}((1,0))
\qquad\text{and}\qquad
W=\operatorname{span}((0,1))
\]
be the coordinate axes in $V=\mathbb{R}^2$.

As an internal direct sum,
\[
\mathbb{R}^2=U\oplus W.
\]
Indeed, $U+W=\mathbb{R}^2$ and $U\cap W=\{(0,0)\}$. An element of this
internal direct sum is an ordinary vector $(a,b)\in\mathbb{R}^2$, and its
unique decomposition is
\[
(a,b)=(a,0)+(0,b),
\qquad (a,0)\in U,\quad (0,b)\in W.
\]
Thus, the summands are subspaces of the same ambient vector space, and their
vectors are added inside that space.

As an external direct sum, however, we construct the new vector space
\[
U\oplus_{\mathrm{ext}}W
=\{((a,0),(0,b)):a,b\in\mathbb{R}\}.
\]
Its elements are ordered pairs of vectors. For example,
\[
((2,0),(0,3))\in U\oplus_{\mathrm{ext}}W,
\]
whereas the corresponding element of the internal direct sum is the single
vector
\[
(2,0)+(0,3)=(2,3)\in\mathbb{R}^2.
\]

\textbf{Direct Product} Let ${(V_i)}_{i\in I}$ be a family of
vector spaces over the same field $\mathbb{F}$. Their \emph{direct product} is
\[
\prod_{i\in I}V_i
=\{{(v_i)}_{i\in I}:v_i\in V_i\text{ for every }i\in I\}.
\]
Addition and scalar multiplication are defined componentwise:
\[
{(v_i)}_{i\in I}+{(w_i)}_{i\in I}={(v_i+w_i)}_{i\in I},
\qquad
a*{(v_i)}_{i\in I}={(a*v_i)}_{i\in I}.
\]
With these operations, $\prod_{i\in I}V_i$ is a vector space over $\mathbb{F}$.

If $I=\{1,\ldots,m\}$ is finite, this is written
\[
V_1\times\cdots\times V_m,
\]
and the direct product and external direct sum contain the same tuples. If
$I$ is infinite, the direct product permits infinitely many nonzero
components, whereas the direct sum
\[
\bigoplus_{i\in I}V_i
\]
contains only tuples with finitely many nonzero components. Thus,
\[
\bigoplus_{i\in I}V_i\subseteq\prod_{i\in I}V_i.
\]


\subsection{Bonus: Linear Algebra and Set Theory}
\textit{Read only if you have read the Set Theory portions involved.}

\begin{problem_box}
\textbf{Theorem:} Let $B$ and
$C$ be two bases of an infinite-dimensional vector space $V$. Assuming the
usual axiom of choice, we have
\[
|B|=|C|.
\]

\end{problem_box}

\begin{claim_box}
\textit{Proof:} For each $b\in B$, write $b$ in its unique basis expansion
with respect to $C$. Because a basis expansion is a finite linear combination,
there is a finite subset $C_b\subseteq C$ such that
\[
b\in\operatorname{span}(C_b).
\]
Let
\[
C_0=\bigcup_{b\in B}C_b.
\]
Every element of $B$ lies in $\operatorname{span}(C_0)$. Since $B$ spans
$V$, it follows that $C_0$ also spans $V$.

We claim that $C_0=C$. Certainly $C_0\subseteq C$. If there were some
$c\in C\setminus C_0$, then, because $C_0$ spans $V$, the vector $c$ would
belong to $\operatorname{span}(C_0)$. Thus $c$ would be a linear combination
of other elements of the basis $C$, contradicting the linear independence of
$C$. Hence $C_0=C$.

Therefore $C$ is a union of $|B|$ finite sets. Since $B$ is infinite,
cardinal arithmetic gives
\[
|C|\leq |B|\cdot\aleph_0=|B|.
\]
Interchanging the roles of $B$ and $C$ gives
\[
|B|\leq |C|.
\]
By the Cantor-Schr\"oder-Bernstein theorem, these two inequalities imply
\[
|B|=|C|.
\]
Thus every basis of $V$ has the same cardinality. \hfill$\square$
\end{claim_box}

\begin{example_box}
\textbf{Example:} Regard $\mathbb{R}$ as
a vector space over $\mathbb{Q}$. A basis of this vector space is called a
\emph{Hamel basis of $\mathbb{R}$ over $\mathbb{Q}$}. Thus a Hamel basis is a
subset $B\subseteq\mathbb{R}$, a set of real numbers $\{u_\alpha\}$ such that every real number $\beta$ has a unique representation of the form
\[
\beta = r_1 u_{\alpha_1} + r_2 u_{\alpha_2} + \cdots + r_n u_{\alpha_n},
\]
where all $u_{\alpha_i}$ are unique and $r_i\in\mathbb{Q}\setminus\{0\}$ and $n$ depends on $\beta$.
The existence of such a basis follows from Zorn's lemma, and hence from the
axiom of choice, but the proof does not produce a concrete list or formula
for all its elements. Moreover, every Hamel basis of $\mathbb{R}$ over
$\mathbb{Q}$ is uncountable: if $B$ were countable, then the set of all finite
rational linear combinations of elements of $B$ would be countable, whereas
$\mathbb{R}$ is uncountable. 
\end{example_box}

\begin{problem_box}
\textbf{Theorem:} Assuming the axiom of
choice, every vector space $U$ has a basis.

\end{problem_box}

\begin{claim_box}
\textit{Proof:} Let
\[
\mathcal{I}=\{L\subseteq U:L\text{ is linearly independent}\},
\]
ordered by inclusion. The collection $\mathcal{I}$ is nonempty because the
empty set is linearly independent.

Let $\mathcal{C}$ be a chain in $\mathcal{I}$ and define
\[
L_{\mathcal{C}}=\bigcup_{L\in\mathcal{C}}L.
\]
We claim that $L_{\mathcal{C}}$ is linearly independent. Any linear relation
uses only finitely many vectors from $L_{\mathcal{C}}$. Because
$\mathcal{C}$ is totally ordered by inclusion, all these finitely many
vectors lie together in some member $L\in\mathcal{C}$. Since $L$ is linearly
independent, all coefficients in the relation must be zero. Therefore
$L_{\mathcal{C}}\in\mathcal{I}$, and it is an upper bound for the chain.

Every chain in $\mathcal{I}$ thus has an upper bound. By Zorn's lemma,
$\mathcal{I}$ has a maximal element; call it $B_U$. It remains to show that
$B_U$ spans $U$. If some $u\in U$ did not belong to
$\operatorname{span}(B_U)$, then
\[
B_U\cup\{u\}
\]
would still be linearly independent, contradicting the maximality of $B_U$.
Hence $\operatorname{span}(B_U)=U$. Thus $B_U$ is both linearly independent
and spanning, so it is a basis of $U$. \hfill$\square$
\end{claim_box}

\begin{problem_box}
\textbf{Theorem:} Let
$V$ be an infinite-dimensional vector space over $\mathbb{F}$, and let $U$ be any
subspace of $V$. Assuming the axiom of choice, there exists a subspace
$W\subseteq V$ such that
\[
V=U\oplus W.
\]

\end{problem_box}

\begin{claim_box}
\textit{Proof:} By Zorn's lemma, $U$ has a basis $B_U$. Consider the
collection
\[
\mathcal{L}
=\{L\subseteq V:B_U\subseteq L\text{ and }L\text{ is linearly independent}\},
\]
ordered by inclusion. This collection is nonempty because $B_U\in\mathcal{L}$.

Let $\mathcal{C}$ be a chain in $\mathcal{L}$ and set
\[
L_{\mathcal{C}}=\bigcup_{L\in\mathcal{C}}L.
\]
Every finite subset of $L_{\mathcal{C}}$ is contained in one member of the
chain $\mathcal{C}$ and is therefore linearly independent. Hence
$L_{\mathcal{C}}$ is linearly independent, contains $B_U$, and is an upper
bound for $\mathcal{C}$ in $\mathcal{L}$. By Zorn's lemma, $\mathcal{L}$ has
a maximal element $B$.

The set $B$ spans $V$. Indeed, if some $v\in V$ did not belong to
$\operatorname{span}(B)$, then $B\cup\{v\}$ would be linearly independent,
contradicting the maximality of $B$. Thus $B$ is a basis of $V$ extending
$B_U$.

Let
\[
C=B\setminus B_U
\qquad\text{and}\qquad
W=\operatorname{span}(C).
\]
Because $B=B_U\cup C$ spans $V$, every vector in $V$ is a sum of a vector in
$U=\operatorname{span}(B_U)$ and a vector in $W$. Therefore $U+W=V$.

Now suppose $x\in U\cap W$. Then $x$ has one basis expansion using elements
of $B_U$ and another using elements of $C$. Subtracting these expansions
gives a finite linear combination of elements of the linearly independent set
$B_U\cup C=B$ equal to $0_V$. All coefficients must therefore be zero, so
$x=0_V$. Hence
\[
U\cap W=\{0_V\}.
\]
Consequently, $V=U\oplus W$, and $W$ is a complementary subspace of $U$.
\hfill$\square$
\end{claim_box}


\section{Linear Transformations}

\begin{definition_box}
\textbf{Definition:} Let $V$ and $W$ be vector spaces over the same
field $\mathbb{F}$. A map $T:V\longrightarrow W$ is a \emph{linear map},
or \emph{linear transformation}, if for all $u,v\in V$ and
$a\in\mathbb{F}$:
\begin{enumerate}
    \item (Additivity) $T(u+v)=T(u)+T(v)$;
    \item (Homogeneity) $T(av)=aT(v)$.
\end{enumerate}
Equivalently, for all $u,v\in V$ and $a,b\in\mathbb{F}$,
\[
T(au+bv)=aT(u)+bT(v).
\]
\end{definition_box}

\begin{problem_box}
\textbf{Proposition:} If
$T:V\longrightarrow W$ is a linear map, then
\[
T(0_V)=0_W.
\]

\end{problem_box}

\begin{claim_box}
\textit{Proof:} By additivity and $0_V+0_V=0_V$,
\[
T(0_V)=T(0_V+0_V)=T(0_V)+T(0_V).
\]
Adding the additive inverse of $T(0_V)$ to both sides in $W$ gives
$0_W=T(0_V)$, as required. \hfill$\square$
\end{claim_box}

\begin{problem_box}
\textbf{Proposition:} Let
$V$ be a one-dimensional vector space over $\mathbb{F}$ and let
$T:V\longrightarrow V$ be linear. There exists a unique scalar
$\lambda\in\mathbb{F}$ such that
\[
T(v)=\lambda v\qquad\text{for every }v\in V.
\]

\end{problem_box}

\begin{claim_box}
\textit{Proof:} Choose a basis $\{e\}$ of $V$. Since $T(e)\in V$ and
$e$ spans $V$, there is a scalar $\lambda\in\mathbb{F}$ such that
\[
T(e)=\lambda e.
\]
Every $v\in V$ has the form $v=ae$ for some $a\in\mathbb{F}$.
By linearity and commutativity of scalar multiplication in the field,
\[
T(v)=T(ae)=aT(e)=a(\lambda e)=\lambda(ae)=\lambda v.
\]
Thus the same scalar $\lambda$ works for every vector in $V$.

For uniqueness, suppose also that $T(v)=\mu v$ for every $v\in V$.
Evaluating at $e$ gives $(\lambda-\mu)e=0_V$. Since $e\neq0_V$,
this implies $\lambda-\mu=0$, so $\lambda=\mu$.
\hfill$\square$
\end{claim_box}

\begin{claim_box}
\textbf{Note:} Conversely, multiplication by any fixed scalar is
a linear map. Hence these are exactly the linear maps from a
one-dimensional vector space to itself, including the zero map when
$\lambda=0$.
\end{claim_box}

\begin{definition_box}
\textbf{Notation:} The set of all linear maps from $V$ to $W$ over
$\mathbb{F}$ is denoted by $\mathcal{L}(V,W)$:
\[
\mathcal{L}(V,W)=\{T:V\longrightarrow W:T\text{ is linear}\}.
\]
Thus, to say that $T$ is a linear map from $V$ to $W$, we write
$T\in\mathcal{L}(V,W)$. The symbol $\mathcal{L}(V,W)$ denotes the whole set
of such maps, not a single map. When $W=V$, we abbreviate
$\mathcal{L}(V,V)$ to $\mathcal{L}(V)$; its elements are called
\emph{linear operators on $V$}.
\end{definition_box}

\begin{problem_box}
\textbf{Vector space of linear maps}
Let $V$ and $W$
be vector spaces over the same field $\mathbb{F}$. Then
$\mathcal{L}(V,W)$ is a vector space over $\mathbb{F}$ under pointwise
addition and scalar multiplication, defined by
\[
(S+T)(v)=S(v)+T(v),
\qquad (aT)(v)=aT(v)
\]
for $S,T\in\mathcal{L}(V,W)$, $a\in\mathbb{F}$, and $v\in V$.
Here the operations on the right take place in $W$.
\end{problem_box}

\begin{claim_box}
\textit{Proof:} First we check closure. Let $S,T\in\mathcal{L}(V,W)$,
$u,v\in V$, and $a,b,c\in\mathbb{F}$. By linearity of $S$ and $T$,
\[
\begin{aligned}
(S+T)(au+bv)
&=S(au+bv)+T(au+bv)\\
&=aS(u)+bS(v)+aT(u)+bT(v)\\
&=a(S+T)(u)+b(S+T)(v).
\end{aligned}
\]
Thus $S+T$ is linear. Similarly,
\[
\begin{aligned}
(cT)(au+bv)
&=c\bigl(aT(u)+bT(v)\bigr)\\
&=a(cT)(u)+b(cT)(v),
\end{aligned}
\]
so $cT$ is linear as well.

The zero map $0_{\mathcal{L}}:V\longrightarrow W$, defined by
$0_{\mathcal{L}}(v)=0_W$ for every $v\in V$, is linear and satisfies
$T+0_{\mathcal{L}}=T$. For each $T$, the map $-T=(-1)T$ is linear and
\[
(T+(-T))(v)=T(v)-T(v)=0_W,
\]
so it is the additive inverse of $T$.

The remaining axioms hold pointwise because they hold in $W$. Explicitly,
for $R,S,T\in\mathcal{L}(V,W)$, $a,b\in\mathbb{F}$, and every $v\in V$,
\[
\begin{aligned}
(S+T)(v)&=S(v)+T(v)=T(v)+S(v)=(T+S)(v),\\
((R+S)+T)(v)&=(R(v)+S(v))+T(v)\\
&=R(v)+(S(v)+T(v))=(R+(S+T))(v),\\
(a(S+T))(v)&=a(S(v)+T(v))=(aS+aT)(v),\\
((a+b)T)(v)&=(a+b)T(v)=(aT+bT)(v),\\
(a(bT))(v)&=a(bT(v))=(ab)T(v)=((ab)T)(v),\\
(1T)(v)&=T(v).
\end{aligned}
\]
Two maps are equal when they agree on every vector in their domain.
Consequently, these pointwise identities establish all the corresponding
vector-space axioms for the maps themselves. Hence $\mathcal{L}(V,W)$ is
a vector space over $\mathbb{F}$. \hfill$\square$
\end{claim_box}

\begin{claim_box}
\textbf{Note:} The vectors of this space are linear maps.
Scalar multiplication here means multiplying the output of a map by a
scalar; it is not composition of maps. No finite-dimensional assumption
on $V$ or $W$ is needed.
\end{claim_box}

\begin{definition_box}
\textbf{Definition:} Let $U,V,W$ be vector
spaces over the same field $\mathbb{F}$. If $T\in\mathcal{L}(U,V)$ and
$S\in\mathcal{L}(V,W)$, their \emph{product} $ST$ is the composition
\[
ST=S\circ T:U\longrightarrow W,
\qquad (ST)(u)=S(T(u)).
\]
The map on the right acts first: apply $T$, then $S$. The intermediate
spaces must match so that the output of $T$ can be used as input to $S$.
This product is not pointwise multiplication of the outputs.
\end{definition_box}

\begin{problem_box}
\textbf{Proposition:} The product of two composable linear maps is
linear. In the notation above, $ST\in\mathcal{L}(U,W)$.

\end{problem_box}

\begin{claim_box}
\textit{Proof:} For $u,v\in U$ and $a,b\in\mathbb{F}$,
\[
\begin{aligned}
(ST)(au+bv)&=S(T(au+bv))\\
&=S(aT(u)+bT(v))\\
&=aS(T(u))+bS(T(v))\\
&=a(ST)(u)+b(ST)(v).
\end{aligned}
\]
Thus $ST$ is linear. \hfill$\square$
\end{claim_box}

\begin{claim_box}
\textbf{Properties:} Whenever the domains and codomains match,
composition is associative and distributes over addition:
\[
(RS)T=R(ST),\qquad
S(T_1+T_2)=ST_1+ST_2,\qquad
(S_1+S_2)T=S_1T+S_2T.
\]
Also, $a(ST)=(aS)T=S(aT)$ for $a\in\mathbb{F}$.
The identity map $I_V:V\longrightarrow V$, defined by $I_V(v)=v$, is
linear. For $T\in\mathcal{L}(U,V)$, it satisfies
\[
I_V T=T=T I_U.
\]
In particular, two operators in $\mathcal{L}(V)$ can always be multiplied,
but their product need not commute.
\end{claim_box}

\begin{example_box}
\textbf{Example:} Define operators on $\mathbb{R}^2$ by
\[
T(x,y)=(x,0),\qquad S(x,y)=(y,x).
\]
Then
\[
(ST)(x,y)=(0,x),\qquad (TS)(x,y)=(y,0).
\]
For instance, $(ST)(1,0)=(0,1)$ but $(TS)(1,0)=(0,0)$, so $ST\neq TS$.
For maps between different spaces, the reverse product may not even be
defined.
\end{example_box}

\begin{problem_box}
\textbf{Proposition:} Let
$L:\mathbb{R}^n\longrightarrow\mathbb{R}^m$ be a linear map. The image of
every line in $\mathbb{R}^n$ is either a line or a single point in
$\mathbb{R}^m$.

\end{problem_box}

\begin{claim_box}
\textit{Proof:} Any line $\ell$ in $\mathbb{R}^n$ can be written as
\[
\ell=\{p+tv:t\in\mathbb{R}\},
\]
where $p\in\mathbb{R}^n$ is a point on the line and
$v\in\mathbb{R}^n\setminus\{0\}$ is its direction vector. By linearity,
\[
L(p+tv)=L(p)+tL(v)
\qquad\text{for every }t\in\mathbb{R}.
\]
Therefore the image of the entire line is
\[
L(\ell)=\{L(p)+tL(v):t\in\mathbb{R}\}.
\]
There are two cases:
\begin{enumerate}
    \item If $L(v)\neq0$, this set is the line through $L(p)$ with
    direction vector $L(v)$.
    \item If $L(v)=0$, then $L(p+tv)=L(p)$ for every $t$, so
    $L(\ell)=\{L(p)\}$ is a single point.
\end{enumerate}
Thus a linear map sends every line to a line or a point. The original line
need not pass through the origin, and the image point need not be the
origin. \hfill$\square$
\end{claim_box}

\begin{problem_box}
\textbf{Open-ended problem (a converse for maps of the plane):} Let
$f:\mathbb{R}^2\longrightarrow\mathbb{R}^2$ be a map satisfying:
\begin{enumerate}
    \item $f$ is injective;
    \item $f(0)=0$;
    \item for every line $\ell\subseteq\mathbb{R}^2$, its image
    $f(\ell)=\{f(x):x\in\ell\}$ is a line in $\mathbb{R}^2$.
\end{enumerate}
Prove that $f$ is a linear map over $\mathbb{R}$; that is, prove that
\[
f(au+bv)=af(u)+bf(v)
\]
for all $u,v\in\mathbb{R}^2$ and $a,b\in\mathbb{R}$.
Here the third assumption says that the image is an entire line, not merely
a subset of a line. No continuity assumption is given.
\end{problem_box}

\begin{claim_box}
\textbf{Differentiation:} Let $\mathcal{P}_n(\mathbb{F})$
denote the polynomials of degree at most $n$ over $\mathbb{F}$, including
the zero polynomial. For $n\geq1$, formal differentiation defines a map
\[
D:\mathcal{P}_n(\mathbb{F})\longrightarrow\mathcal{P}_{n-1}(\mathbb{F}),
\qquad D(p)=p'.
\]
Explicitly,
\[
D(a_0+a_1x+\cdots+a_n x^n)
=a_1+2a_2x+\cdots+na_n x^{n-1}.
\]
For any $p,q\in\mathcal{P}_n(\mathbb{F})$ and $a,b\in\mathbb{F}$,
\[
D(ap+bq)=aD(p)+bD(q),
\]
so $D$ is linear. Over $\mathbb{R}$, for example,
$D(3x^2+2x+5)=6x+2$. A nonzero constant polynomial is sent to zero,
showing that a linear map need not be injective.
\end{claim_box}

\begin{claim_box}
\textbf{Integration:} Let $C([0,1],\mathbb{R})$ denote
the vector space of continuous real-valued functions on $[0,1]$, with
pointwise addition and scalar multiplication. Define
\[
J:C([0,1],\mathbb{R})\longrightarrow C([0,1],\mathbb{R}),
\qquad (Jf)(x)=\int_0^x f(t)\,dt.
\]
The function $Jf$ is continuous, so the map has the stated codomain.
For any $f,g\in C([0,1],\mathbb{R})$ and $a,b\in\mathbb{R}$,
\[
\begin{aligned}
(J(af+bg))(x)
&=\int_0^x \bigl(af(t)+bg(t)\bigr)\,dt\\
&=a\int_0^x f(t)\,dt+b\int_0^x g(t)\,dt\\
&=a(Jf)(x)+b(Jg)(x).
\end{aligned}
\]
Thus $J$ is linear. For example, if $f(t)=t^2$, then $(Jf)(x)=x^3/3$.
The fixed lower limit specifies a unique antiderivative with value zero
at $0$; an indefinite integral with an unspecified constant does not
specify a single function-valued map.

Integration over a fixed interval also gives a linear map with scalar
values:
\[
I:C([0,1],\mathbb{R})\longrightarrow\mathbb{R},
\qquad I(f)=\int_0^1 f(t)\,dt.
\]
Indeed, $I(af+bg)=aI(f)+bI(g)$.
\end{claim_box}

\begin{claim_box}
\textbf{Matrix maps:} Fix an $m\times n$ matrix
$A=(a_{ij})$ with entries in $\mathbb{F}$. Regarding vectors as column
vectors, define
\[
T_A:\mathbb{F}^n\longrightarrow\mathbb{F}^m,
\qquad T_A(x)=Ax.
\]
The $i$th coordinate of $T_A(x)$ is $\sum_{j=1}^n a_{ij}x_j$.
For $u,v\in\mathbb{F}^n$ and $a,b\in\mathbb{F}$,
\[
T_A(au+bv)=A(au+bv)=aAu+bAv=aT_A(u)+bT_A(v),
\]
so $T_A$ is linear. For example,
\[
A=\begin{pmatrix}1&2&0\\0&-1&1\end{pmatrix}
\quad\text{gives}\quad
T_A:\mathbb{F}^3\longrightarrow\mathbb{F}^2,
\qquad T_A(x,y,z)=(x+2y,-y+z).
\]
\end{claim_box}

\begin{example_box}
\textbf{Example:} The map
\[
T:\mathbb{R}^2\longrightarrow\mathbb{R}^2,
\qquad T(x,y)=(2x,3y),
\]
is linear: it preserves both vector addition and scalar multiplication.
\end{example_box}

\begin{claim_box}
\textbf{Counterexample (homogeneity without additivity):} Regard
$\mathbb{R}^2$ and $\mathbb{R}$ as vector spaces over $\mathbb{R}$. Define
\[
H:\mathbb{R}^2\longrightarrow\mathbb{R},
\qquad
H(x,y)=
\begin{cases}
\dfrac{x^3}{x^2+y^2}, & (x,y)\neq(0,0),\\[4pt]
0, & (x,y)=(0,0).
\end{cases}
\]
For every nonzero $a\in\mathbb{R}$ and $(x,y)\neq(0,0)$,
\[
H(ax,ay)=\frac{a^3x^3}{a^2(x^2+y^2)}=aH(x,y).
\]
The same identity holds when $a=0$ or $(x,y)=(0,0)$, since both sides
are zero. Thus $H$ is homogeneous for every real scalar, including
negative scalars. However,
\[
H((1,0)+(0,1))=H(1,1)=\frac12,
\qquad H(1,0)+H(0,1)=1+0=1.
\]
So $H$ is not additive and therefore is not linear.
\end{claim_box}

\begin{claim_box}
\textbf{Counterexample (additivity without homogeneity):} Regard
$\mathbb{C}$ as a vector space over $\mathbb{C}$ and define complex
conjugation by
\[
K:\mathbb{C}\longrightarrow\mathbb{C},
\qquad K(z)=\overline{z}.
\]
It is additive because, for all $z,w\in\mathbb{C}$,
\[
K(z+w)=\overline{z+w}=\overline{z}+\overline{w}=K(z)+K(w).
\]
But homogeneity fails for the complex scalar $i$ and the vector $1$:
\[
K(i\cdot1)=\overline{i}=-i,
\qquad iK(1)=i.
\]
Hence $K$ is not linear over $\mathbb{C}$.
\end{claim_box}

\begin{remark_box}
\textbf{Note:} Conjugation does satisfy
$K(az)=aK(z)$ for every real scalar $a$. Thus it is linear when
$\mathbb{C}$ is regarded as a vector space over $\mathbb{R}$, but not when
the scalar field is $\mathbb{C}$. These counterexamples show why both
additivity and homogeneity are required in the definition of a linear map.
\end{remark_box}



\subsection{Isomorphisms}

\begin{problem_box}
\textbf{Proposition:} Let $V$ and
$W$ be vector spaces over $\mathbb{F}$, let $I$ be an index set, and let
$(v_i)_{i\in I}$ be vectors in $V$ and $(w_i)_{i\in I}$ vectors in $W$.
\begin{enumerate}
    \item If $(v_i)_{i\in I}$ is a basis of $V$, there exists a unique
    linear map $T:V\longrightarrow W$ such that $T(v_i)=w_i$ for every
    $i\in I$.
    \item If $(v_i)_{i\in I}$ is linearly independent, there exists at
    least one linear map $T:V\longrightarrow W$ such that
    $T(v_i)=w_i$ for every $i\in I$.
    \item If $(v_i)_{i\in I}$ spans $V$, there exists at most one linear
    map $T:V\longrightarrow W$ such that $T(v_i)=w_i$ for every $i\in I$.
\end{enumerate}
\end{problem_box}

\begin{claim_box}
\textit{Proof:} For (1), every $v\in V$ has a unique expression as a
finite linear combination of basis vectors,
\[
v=\sum_{i\in J}a_i v_i,
\]
where $J\subseteq I$ is finite. Define
\[
T(v)=\sum_{i\in J}a_i w_i.
\]
Uniqueness of the basis expansion makes this definition well-defined,
and the definition directly gives linearity and $T(v_i)=w_i$. If
$S:V\longrightarrow W$ is another linear map with these values, then
\[
S(v)=\sum_{i\in J}a_i S(v_i)
=\sum_{i\in J}a_i w_i=T(v)
\]
for every $v\in V$. Thus $S=T$.

For (2), extend the linearly independent family $(v_i)_{i\in I}$ to a
basis of $V$. Assign each added basis vector the value $0_W$, while
assigning $v_i$ the prescribed value $w_i$. The basis case in (1)
provides a linear map with all these values, hence one with the required
values on the original family.

For (3), suppose $S,T:V\longrightarrow W$ are linear and agree on every
$v_i$. Since the family spans $V$, each $v\in V$ is a finite linear
combination $v=\sum_{i\in J}a_i v_i$. Therefore
\[
S(v)=\sum_{i\in J}a_i S(v_i)
=\sum_{i\in J}a_i T(v_i)=T(v).
\]
Thus $S=T$, proving uniqueness whenever such a map exists.
\hfill$\square$
\end{claim_box}

\begin{definition_box}
\textbf{Definition:} Let $V$ and $W$ be vector spaces
over the same field $\mathbb{F}$. A map $T:V\longrightarrow W$ is a
\emph{vector-space isomorphism} if it is linear and bijective. Thus every
$w\in W$ has exactly one preimage $v\in V$ under $T$.
If an isomorphism exists, we call $V$ and $W$ \emph{isomorphic} and write
\[
V\cong W.
\]
This means that the spaces have the same vector-space structure, even if
their elements are different kinds of objects.
\end{definition_box}

\begin{definition_box}
\textbf{Definition:} A linear map
$T:V\longrightarrow W$ is \emph{invertible} if there is a map
$R:W\longrightarrow V$ such that
\[
RT=I_V,\qquad TR=I_W.
\]
Here $I_V$ and $I_W$ are the identity maps. These equations say that
$R(T(v))=v$ for every $v\in V$ and $T(R(w))=w$ for every $w\in W$:
$R$ undoes $T$ in both directions. The map $R$ is called the
\emph{inverse} of $T$ and is denoted by $T^{-1}$.
\end{definition_box}

\begin{claim_box}
The inverse is unique. Indeed, if $R$ and $Q$ are both inverses of $T$,
then associativity of composition gives
\[
R=R I_W=R(TQ)=(RT)Q=I_V Q=Q.
\]
\end{claim_box}

\begin{problem_box}
\textbf{Lemma:} If
$T:V\longrightarrow W$ is an isomorphism, then its inverse
$T^{-1}:W\longrightarrow V$ is also an isomorphism.

\end{problem_box}

\begin{claim_box}
\textit{Proof:} Since $T$ is bijective, its inverse exists and is bijective.
It remains to prove that $T^{-1}$ is linear. Let $w_1,w_2\in W$ and
$a,b\in\mathbb{F}$. Set
\[
v_1=T^{-1}(w_1),\qquad v_2=T^{-1}(w_2).
\]
Then $T(v_1)=w_1$ and $T(v_2)=w_2$. By linearity of $T$,
\[
T(av_1+bv_2)=aT(v_1)+bT(v_2)=aw_1+bw_2.
\]
Applying $T^{-1}$ to both sides gives
\[
\begin{aligned}
T^{-1}(aw_1+bw_2)
&=av_1+bv_2\\
&=aT^{-1}(w_1)+bT^{-1}(w_2).
\end{aligned}
\]
Thus $T^{-1}$ is linear and bijective, hence an isomorphism.
\hfill$\square$
\end{claim_box}

\begin{problem_box}
\textbf{Proposition:} A linear
map $T:V\longrightarrow W$ is invertible if and only if it is an
isomorphism. In that case $T^{-1}$ is automatically linear.

\end{problem_box}

\begin{claim_box}
\textit{Proof:} Suppose $R$ is an inverse of $T$. If $T(u)=T(v)$,
applying $R$ gives $u=v$, so $T$ is injective. For every $w\in W$,
the identity $T(R(w))=w$ shows that $T$ is surjective. Thus $T$ is
linear and bijective, hence an isomorphism.
Conversely, an isomorphism is bijective, so each $w\in W$ has a unique
preimage. Assigning that preimage to $w$ defines a map $T^{-1}$
satisfying both inverse identities. The preceding lemma proves that
this inverse is linear. \hfill$\square$
\end{claim_box}

\begin{example_box}
\textbf{Example:} On a nonzero vector
space $V$, the map $T(v)=\lambda v$ is invertible exactly when
$\lambda\neq0$. In that case $T^{-1}(v)=\lambda^{-1}v$, since composing
the two maps in either order gives $v$. If $\lambda=0$, all vectors
have the same image, so the map is not injective. In particular,
$T(x,y)=(2x,3y)$ on $\mathbb{R}^2$ is invertible with
$T^{-1}(x,y)=(x/2,y/3)$.
\end{example_box}

\begin{example_box}
\textbf{Example:} Let
$\mathcal{P}_2(\mathbb{F})$ be the space of polynomials of degree at most
$2$, including zero. The map
\[
T:\mathcal{P}_2(\mathbb{F})\longrightarrow\mathbb{F}^3,
\qquad T(a+bx+cx^2)=(a,b,c),
\]
is linear and bijective. Its inverse sends $(a,b,c)$ to $a+bx+cx^2$.
Consequently, $\mathcal{P}_2(\mathbb{F})\cong\mathbb{F}^3$.
\end{example_box}

\begin{example_box}
\textbf{Example:} The map
\[
S:\mathbb{R}^2\longrightarrow\mathbb{R}^2,
\qquad S(x,y)=(2x+y,y)
\]
is linear. Given $(u,v)\in\mathbb{R}^2$, its unique preimage is
$((u-v)/2,v)$, so
\[
S^{-1}(u,v)=\left(\frac{u-v}{2},v\right).
\]
Thus $S$ is an isomorphism. It sends $(1,0)$ to $(2,0)$, showing that an
isomorphism need not preserve Euclidean length. An isomorphism from a
vector space to itself is called an \emph{automorphism}.
\end{example_box}

\begin{example_box}
\textbf{Example:} Regard
$\mathbb{C}$ as a vector space over $\mathbb{R}$. Then
\[
\Phi:\mathbb{C}\longrightarrow\mathbb{R}^2,
\qquad \Phi(a+bi)=(a,b)
\]
is a real-linear isomorphism, with inverse $(a,b)\mapsto a+bi$.
The scalar field is essential: $\mathbb{C}$ has dimension $2$ over
$\mathbb{R}$ but dimension $1$ over $\mathbb{C}$.
\end{example_box}

\begin{problem_box}
\textbf{Proposition:}
Let $T:V_1\longrightarrow V_2$ be a linear map.
\begin{enumerate}
    \item If $W_1$ is a subspace of $V_1$, then
    \[
    W_2=T(W_1)=\{T(w):w\in W_1\}
    \]
    is a subspace of $V_2$.
    \item If $T$ is onto, then the image under $T$ of any spanning set
    of $V_1$ is a spanning set of $V_2$. More generally, if $W_1$ is a
    subspace of $V_1$ and $S$ spans $W_1$, then $T(S)$ spans
    $W_2=T(W_1)$.
    \item If $T$ is injective, then the image under $T$ of every linearly
    independent subset of $V_1$ is linearly independent in $V_2$.
    \item If $T$ is an isomorphism, then $T$ and $T^{-1}$ preserve
    linear independence, spanning, and bases. In particular, a subset
    $B\subseteq V_1$ is a basis of $V_1$ if and only if $T(B)$ is a
    basis of $V_2$.
\end{enumerate}
\end{problem_box}

\begin{claim_box}
\textit{Proof:}
\begin{enumerate}
\item $0_{V_1}\in W_1$, so linearity gives
$0_{V_2}=T(0_{V_1})\in T(W_1)$. If $x,y\in T(W_1)$ and
$a,b\in\mathbb{F}$, write $x=T(u)$ and $y=T(v)$ with $u,v\in W_1$.
Then
\[
ax+by=aT(u)+bT(v)=T(au+bv)\in T(W_1),
\]
because $W_1$ is a subspace. Thus $T(W_1)$ is a subspace of $V_2$.
\hfill$\square$
\item First, let $S$ span $V_1$ and suppose $T$ is onto. For any
$v_2\in V_2$, choose $v_1\in V_1$ with $T(v_1)=v_2$. Since $S$
spans $V_1$, write $v_1$ as a finite linear combination of elements
of $S$. By linearity, $v_2=T(v_1)$ is the same linear combination
of elements of $T(S)$; hence $T(S)$ spans $V_2$. More generally, if
$S$ spans $W_1$, every element of $T(W_1)$ is the image of a linear
combination of elements of $S$, so it is a linear combination of
elements of $T(S)$. Therefore $T(S)$ spans $W_2=T(W_1)$.
\hfill$\square$
\item Let $v_1,\ldots,v_n$ be distinct elements of a linearly
independent subset of $V_1$, and suppose
\[
a_1T(v_1)+\cdots+a_nT(v_n)=0_{V_2}.
\]
By linearity, $T(a_1v_1+\cdots+a_nv_n)=0_{V_2}=T(0_{V_1})$.
Injectivity implies $a_1v_1+\cdots+a_nv_n=0_{V_1}$, and the
independence of the original set gives $a_1=\cdots=a_n=0$.
Thus the image set is linearly independent.
\hfill$\square$
\item If $T$ is an isomorphism, it is both injective and onto, so
(2) and (3) show that it preserves spanning sets and linear independence.
Its inverse is also an isomorphism, so the same holds for $T^{-1}$.
A basis is precisely a linearly independent spanning set; consequently
both $T$ and $T^{-1}$ preserve bases, proving (4).
\hfill$\square$
\end{enumerate} 
\end{claim_box}

\begin{remark_box}
\textbf{Caution:} A map $R:W\longrightarrow V$
with $RT=I_V$ is a \emph{left inverse} of $T$, and this identity alone
implies injectivity. A map with $TR=I_W$ is a \emph{right inverse},
and this identity alone implies surjectivity. In general, one identity
does not imply the other. For example, let
\[
J:\mathbb{R}\longrightarrow\mathbb{R}^2,\quad J(t)=(t,0),
\qquad P:\mathbb{R}^2\longrightarrow\mathbb{R},\quad P(x,y)=x.
\]
Then $PJ=I_{\mathbb{R}}$, but $JP(x,y)=(x,0)$ is not
$I_{\mathbb{R}^2}(x,y)$. Neither $J$ nor $P$ is invertible.
When the domain and codomain have the same finite dimension, however,
the \hyperref[thm:dimension-criteria]{dimension theorem below} shows that
either inverse identity guarantees
invertibility, and the one-sided inverse is then the unique inverse.
\end{remark_box}

\begin{problem_box}
\textbf{Proposition:} If
$T:U\longrightarrow V$ and $S:V\longrightarrow W$ are isomorphisms, then
$ST:U\longrightarrow W$ is an isomorphism, and
\[
{(ST)}^{-1}=T^{-1}S^{-1}.
\]

\end{problem_box}

\begin{claim_box}
\textit{Proof:} Composition preserves linearity. Moreover,
\[
(T^{-1}S^{-1})(ST)=I_U,
\qquad (ST)(T^{-1}S^{-1})=I_W,
\]
so $ST$ is bijective with the stated inverse. \hfill$\square$

Together with the identity map and the inverse lemma, this shows that
being isomorphic is reflexive, symmetric, and transitive: $V\cong V$;
if $V\cong W$, then $W\cong V$; and if $U\cong V$ and $V\cong W$,
then $U\cong W$.
\end{claim_box}

\begin{problem_box}
\textbf{Theorem:}
Let
$V$ and $W$ be finite-dimensional vector spaces over the same field
$\mathbb{F}$. Then
\[
\dim\mathcal{L}(V,W)=(\dim V)(\dim W).
\]
\end{problem_box}

\begin{claim_box}
\textit{Proof:} Put $n=\dim V$ and $m=\dim W$. If either space is
the zero space, the only linear map from $V$ to $W$ is the zero map.
Thus $\dim\mathcal{L}(V,W)=0=mn$ in this case.

Otherwise, fix ordered bases $\mathcal{B}=(v_1,\ldots,v_n)$ of $V$
and $\mathcal{C}=(w_1,\ldots,w_m)$ of $W$. We use the matrix
representation and its linearity, developed in the later subsection
\emph{Matrices}. The matrix representation map
\[
M:\mathcal{L}(V,W)\longrightarrow M_{m\times n}(\mathbb{F})
\]
is linear. It is also bijective: a matrix $A=(a_{ij})$ specifies
exactly one linear map, namely
\[
T_A\left(\sum_{j=1}^n x_j v_j\right)
=\sum_{i=1}^m\left(\sum_{j=1}^n a_{ij}x_j\right)w_i.
\]
Unique coordinates in $\mathcal{B}$ make this well-defined, and the
formula is linear in the $x_j$. Its representing matrix is $A$.
Conversely, any linear map with matrix $A$ must have this formula,
because its values on the basis vectors determine its values on every
vector. Hence $M$ is an isomorphism.

For $1\leq i\leq m$ and $1\leq j\leq n$, let $E_{ij}$ be the
$m\times n$ matrix with a $1$ in row $i$, column $j$, and zeros
elsewhere. Every matrix has the unique expansion
\[
A=\sum_{i=1}^m\sum_{j=1}^n a_{ij}E_{ij}.
\]
Indeed, the entry in row $i$, column $j$ of this sum is $a_{ij}$.
Thus the $mn$ matrices $E_{ij}$ span the matrix space and are linearly
independent, so they form a basis. Since an isomorphism carries a basis
to a basis, their inverse images under $M$ form a basis of
$\mathcal{L}(V,W)$. Consequently,
\[
\dim\mathcal{L}(V,W)=mn=(\dim V)(\dim W).
\]
\hfill$\square$
\end{claim_box}

\subsection{Null spaces and ranges}

\begin{definition_box}
\textbf{Definition:} Let
$T:V\longrightarrow W$ be a linear map over $\mathbb{F}$. The
\emph{null space} or \emph{kernel} of $T$ is
\[
\operatorname{null}T=\ker T=\{v\in V:T(v)=0_W\}.
\]
It consists of all input vectors that $T$ sends to zero. In particular,
the null space is a subset of the \emph{domain} $V$, not the codomain $W$.
It always contains $0_V$, but may also contain nonzero vectors.
\end{definition_box}

\begin{problem_box}
\textbf{Proposition:} The null space of $T$ is a subspace of $V$.

\end{problem_box}

\begin{claim_box}
\textit{Proof:} Since $T(0_V)=0_W$, we have $0_V\in\ker T$.
If $u,v\in\ker T$ and $a,b\in\mathbb{F}$, linearity gives
\[
T(au+bv)=aT(u)+bT(v)=a0_W+b0_W=0_W.
\]
Thus $au+bv\in\ker T$, proving that $\ker T$ is a subspace.
\hfill$\square$
\end{claim_box}

\begin{definition_box}
\textbf{Definition:} The \emph{range} or
\emph{image} of a linear map $T:V\longrightarrow W$ is
\[
\operatorname{range}T=\operatorname{im}T=T(V)
=\{T(v):v\in V\}\subseteq W.
\]
It consists of all output vectors that $T$ actually attains. It need not
equal the whole codomain $W$; equality holds exactly when $T$ is surjective.
\end{definition_box}

\begin{problem_box}
\textbf{Proposition:} The range of $T$ is a subspace of $W$.

\end{problem_box}

\begin{claim_box}
\textit{Proof:} Since $T(0_V)=0_W$, the range contains $0_W$.
Let $w_1,w_2\in\operatorname{range}T$ and $a,b\in\mathbb{F}$.
By the definition of the range, there exist $v_1,v_2\in V$ such that
$w_1=T(v_1)$ and $w_2=T(v_2)$. By linearity,
\[
aw_1+bw_2=aT(v_1)+bT(v_2)=T(av_1+bv_2).
\]
Because $av_1+bv_2\in V$, this output belongs to
$\operatorname{range}T$. Hence the range contains zero and is closed under
linear combinations, so it is a subspace of $W$. \hfill$\square$
\end{claim_box}

\begin{claim_box}
\textbf{Summary:} The two subspaces live in different spaces:
\[
\ker T\text{ is a subspace of }V,
\qquad
\operatorname{range}T\text{ is a subspace of }W.
\]
\end{claim_box}

\begin{definition_box}
\textbf{Definition:} The \emph{rank} of a linear map
$T:V\longrightarrow W$ is the dimension of its range:
\[
\operatorname{rank}T=\dim(\operatorname{range}T).
\]
Thus rank measures the number of linearly independent output directions
attained by $T$. If $V$ and $W$ are finite-dimensional, then
\[
0\leq\operatorname{rank}T\leq\min\{\dim V,\dim W\}.
\]

For a matrix $A\in\mathbb{F}^{m\times n}$, its \emph{rank} is the rank
of the map $T_A(x)=Ax$. Since the range of $T_A$ is the span of the
columns $A_1,\ldots,A_n$,
\[
\operatorname{rank}A
=\dim\operatorname{span}\{A_1,\ldots,A_n\}.
\]
Equivalently, it is the largest number of linearly independent columns
of $A$.
\end{definition_box}

\begin{problem_box}
\textbf{Proposition:} Let $A\in\mathbb{F}^{m\times n}$,
$X\in\mathbb{F}^{n\times n}$, and $Y\in\mathbb{F}^{m\times m}$,
where $X$ and $Y$ are invertible. Then
\[
\operatorname{rank}(AX)=\operatorname{rank}(A)
=\operatorname{rank}(YA).
\]
\end{problem_box}

\begin{claim_box}
\textit{Proof:} Since $X$ is invertible, it maps $\mathbb{F}^n$
onto itself, so
\[
\operatorname{range}T_{AX}
=\{AXx:x\in\mathbb{F}^n\}
=\operatorname{range}T_A.
\]
Hence $\operatorname{rank}(AX)=\operatorname{rank}(A)$.

Since $Y$ is invertible, its restriction to $\operatorname{range}T_A$
is an isomorphism onto $Y(\operatorname{range}T_A)$, and
\[
\operatorname{range}T_{YA}=Y(\operatorname{range}T_A).
\]
Thus the two ranges have the same dimension, giving
$\operatorname{rank}(YA)=\operatorname{rank}(A)$.
\hfill$\square$
\end{claim_box}

\begin{example_box}
\textbf{Example:} The projection
$P:\mathbb{R}^2\longrightarrow\mathbb{R}$, $P(x,y)=x$, has range
$\mathbb{R}$, so $\operatorname{rank}P=1$. The zero map has rank $0$,
and the identity map on $\mathbb{F}^n$ has rank $n$.
\end{example_box}

\begin{example_box}
\textbf{Example:} For
\[
P:\mathbb{R}^2\longrightarrow\mathbb{R},\qquad P(x,y)=x,
\]
the equation $P(x,y)=0$ is equivalent to $x=0$. Hence
\[
\ker P=\{(0,y):y\in\mathbb{R}\}
=\operatorname{span}\{(0,1)\}.
\]
The null space is the vertical axis: the projection discards the entire
vertical component.
\end{example_box}

\begin{example_box}
\textbf{Example:} Over $\mathbb{R}$, consider
\[
D:\mathcal{P}_n(\mathbb{R})\longrightarrow\mathcal{P}_{n-1}(\mathbb{R}),
\qquad D(p)=p',\quad n\geq1.
\]
A polynomial has zero derivative exactly when it is constant. Therefore
\[
\ker D=\{p:p(x)=c,\ c\in\mathbb{R}\}
=\operatorname{span}\{1\}.
\]
Here $1$ denotes the constant polynomial, and the zero polynomial is
included by taking $c=0$.
\end{example_box}

\begin{example_box}
\textbf{Example:} For the matrix map
\[
T_A:\mathbb{R}^3\longrightarrow\mathbb{R}^2,
\qquad T_A(x,y,z)=(x+2y,-y+z),
\]
finding the null space means solving the homogeneous system $Av=0$:
\[
x+2y=0,\qquad -y+z=0.
\]
Taking $y=t$ gives $x=-2t$ and $z=t$, so
\[
\ker T_A=\{(-2t,t,t):t\in\mathbb{R}\}
=\operatorname{span}\{(-2,1,1)\}.
\]
\end{example_box}

\begin{example_box}
\textbf{Examples:} The identity map $I_V$ has
$\ker I_V=\{0_V\}$, whereas the zero map $Z:V\longrightarrow W$,
defined by $Z(v)=0_W$, has $\ker Z=V$.
\end{example_box}

\begin{problem_box}
\textbf{Proposition:} A linear map
$T:V\longrightarrow W$ is injective if and only if $\ker T=\{0_V\}$.

\end{problem_box}

\begin{claim_box}
\textit{Proof:} If $T$ is injective and $v\in\ker T$, then
$T(v)=0_W=T(0_V)$ implies $v=0_V$. Conversely, suppose
$\ker T=\{0_V\}$. If $T(u)=T(v)$, then
\[
T(u-v)=T(u)-T(v)=0_W,
\]
so $u-v\in\ker T$, giving $u=v$. Thus $T$ is injective.
\hfill$\square$
\end{claim_box}

\begin{problem_box}
\textbf{Proposition:} Let $T\in\mathcal{L}(V)$ satisfy $T^2=T$. Then
\[
V=\operatorname{range}T\oplus\ker T,
\]
and $T$ is the projector onto $\operatorname{range}T$ along $\ker T$.
\end{problem_box}

\begin{claim_box}
\textit{Proof:} Set $U=\operatorname{range}T$ and $W=\ker T$.
For each $v\in V$,
\[
v=Tv+(v-Tv).
\]
Here $Tv\in U$, and
\[
T(v-Tv)=Tv-T^2v=0_V,
\]
so $v-Tv\in W$. Thus $V=U+W$.

If $x\in U\cap W$, then $x=Tv$ for some $v$ and $Tx=0$. But
$Tx=T^2v=Tv=x$, hence $x=0_V$. Therefore $U\cap W=\{0_V\}$,
so the sum is direct.

For $u\in U$, write $u=Tv$; then $Tu=T^2v=Tv=u$. Also, $Tw=0_V$
for every $w\in W$. Thus, for the unique decomposition $v=u+w$,
\[
T(v)=T(u)+T(w)=u,
\]
which is precisely the projection onto $U$ along $W$.
\hfill$\square$
\end{claim_box}

\begin{problem_box}
\textbf{Rank-Nullity Theorem}
Let $V$ be a
finite-dimensional vector space over $\mathbb{F}$, let $W$ be a vector
space over the same field, and let $T\in\mathcal{L}(V,W)$. Then
$\operatorname{range}T$ is finite-dimensional and
\[
\dim V=\dim(\ker T)+\dim(\operatorname{range}T).
\]
Note that $W$ need not be
finite-dimensional.
\end{problem_box}

\begin{claim_box}
\textit{Proof:} Since $\ker T$ is a subspace of the finite-dimensional
space $V$, it has a basis $u_1,\ldots,u_k$. Extend this to a basis
\[
u_1,\ldots,u_k,v_1,\ldots,v_r
\]
of $V$, so $\dim V=k+r$. We will show that
\[
T(v_1),\ldots,T(v_r)
\]
is a basis of $\operatorname{range}T$.

\emph{Spanning:} Let $w\in\operatorname{range}T$. There is some $v\in V$
with $T(v)=w$. Expanding $v$ in the chosen basis gives
\[
v=\sum_{i=1}^k a_i u_i+\sum_{j=1}^r b_j v_j.
\]
Because each $u_i$ belongs to $\ker T$, applying $T$ yields
\[
w=T(v)=\sum_{i=1}^k a_i T(u_i)+\sum_{j=1}^r b_j T(v_j)
=\sum_{j=1}^r b_j T(v_j).
\]
Thus $T(v_1),\ldots,T(v_r)$ spans the range.

\emph{Linear independence:} Suppose
\[
\sum_{j=1}^r c_j T(v_j)=0_W.
\]
By linearity, $T(\sum_{j=1}^r c_j v_j)=0_W$, so
$\sum_{j=1}^r c_j v_j\in\ker T$. Since the $u_i$ form a basis of the
kernel, there are scalars $d_1,\ldots,d_k$ such that
\[
\sum_{j=1}^r c_j v_j=\sum_{i=1}^k d_i u_i.
\]
Subtracting the right-hand side gives a linear relation among the basis
vectors of $V$. Hence all its coefficients vanish, and in particular
$c_1=\cdots=c_r=0$. This proves independence.

Consequently, the range has a finite basis of length $r$. Therefore
\[
\dim(\ker T)+\dim(\operatorname{range}T)=k+r=\dim V.
\]
Empty basis lists and empty sums are allowed, so the proof also covers
the zero map and the case $V=\{0_V\}$. \hfill$\square$
\end{claim_box}

\phantomsection%
\label{thm:dimension-criteria}%
\begin{problem_box}
\textbf{Theorem:}
Let $V$ and
$W$ be finite-dimensional vector spaces over the same field $\mathbb{F}$.
Write $n=\dim V$ and $m=\dim W$. Then:
\begin{enumerate}
    \item If $n>m$, no linear map $T:V\longrightarrow W$ is injective.
    \item If $n<m$, no linear map $T:V\longrightarrow W$ is surjective.
    \item If $n=m$, a linear map $T:V\longrightarrow W$ is injective
    if and only if it is surjective, and either condition is equivalent
    to invertibility.
    \item $V$ and $W$ are isomorphic if and only if $n=m$.
\end{enumerate}
\end{problem_box}

\begin{claim_box}

\textit{Proof:} For any $T\in\mathcal{L}(V,W)$, put
$r=\dim(\operatorname{range}T)$ and $k=\dim(\ker T)$.
Rank-nullity and the fact that the range is a subspace of $W$ give
\[
n=k+r,\qquad k\geq0,\qquad r\leq m.
\]
Consequently, $r\leq\min\{n,m\}$. Moreover, the kernel criterion for
injectivity and the dimension criterion for equality of subspaces give
\[
\begin{aligned}
T\text{ is injective}&\ \Longleftrightarrow\ k=0
\ \Longleftrightarrow\ r=n,\\
T\text{ is surjective}&\ \Longleftrightarrow\
\operatorname{range}T=W\ \Longleftrightarrow\ r=m.
\end{aligned}
\]
These two equivalences prove the first three claims at once: if $n>m$,
then $r<n$, so injectivity is impossible; if $n<m$, then $r<m$, so
surjectivity is impossible; and if $n=m$, the two conditions coincide.
In the last case, either condition makes $T$ bijective and hence
invertible. They also show that an isomorphism must have $r=n=m$.

It remains to show that equal dimensions guarantee the existence of
an isomorphism. If $n=m$, choose bases $v_1,\ldots,v_n$ of $V$ and
$w_1,\ldots,w_n$ of $W$, and define
\[
T\left(\sum_{i=1}^n a_i v_i\right)=\sum_{i=1}^n a_i w_i.
\]
Unique basis expansions make $T$ well-defined, and coordinate addition
and scalar multiplication show that it is linear. Every vector of $W$
has the displayed form on the right, so $T$ is surjective. By the
equal-dimension case already proved, $T$ is invertible, hence an
isomorphism. Empty bases cover $n=m=0$. \hfill$\square$
\end{claim_box}

\begin{claim_box}
\textbf{Definition:} Let
$T:V\longrightarrow W$ be linear and put $U=\operatorname{null}T=\ker T$.
Recall that $V/U$ consists of cosets $v+U$, with operations
\[
(v+U)+(w+U)=(v+w)+U,\qquad a(v+U)=av+U.
\]
The \emph{map induced by $T$ on $V/\operatorname{null}T$} is
\[
\widetilde{T}:V/\operatorname{null}T\longrightarrow W,
\qquad \widetilde{T}(v+U)=T(v).
\]
Its input is a coset, and its output is the image under $T$ of any
representative of that coset. No finite-dimensional assumption is needed.
\end{claim_box}

\begin{problem_box}
\textbf{Proposition:} The induced map $\widetilde{T}$ is
well-defined, linear, and injective, and its range equals the range of $T$.

\end{problem_box}

\begin{claim_box}
\textit{Proof:} \emph{Well-definedness.} If $v+U=w+U$, then $v-w\in U=\ker T$.
Hence
\[
T(v)-T(w)=T(v-w)=0_W,
\]
so $T(v)=T(w)$. Thus the value of $\widetilde{T}$ does not depend on
the representative of the coset.

\emph{Linearity:} For $v,w\in V$ and $a,b\in\mathbb{F}$,
\[
\begin{aligned}
\widetilde{T}\bigl(a(v+U)+b(w+U)\bigr)
&=\widetilde{T}\bigl((av+bw)+U\bigr)\\
&=T(av+bw)\\
&=a\widetilde{T}(v+U)+b\widetilde{T}(w+U).
\end{aligned}
\]

\emph{Injectivity:} If $\widetilde{T}(v+U)=\widetilde{T}(w+U)$,
then $T(v)=T(w)$. Therefore $T(v-w)=0_W$, so $v-w\in U$ and
$v+U=w+U$.

\emph{Range:} Every value $\widetilde{T}(v+U)=T(v)$ belongs to
$\operatorname{range}T$. Conversely, every $y\in\operatorname{range}T$
equals $T(v)=\widetilde{T}(v+U)$ for some $v\in V$. Hence
$\operatorname{range}\widetilde{T}=\operatorname{range}T$.
\hfill$\square$
\end{claim_box}

\begin{claim_box}
\textbf{Factorization:} The quotient map
$\pi:V\longrightarrow V/U$, $\pi(v)=v+U$, satisfies
\[
T=\widetilde{T}\circ\pi.
\]
Moreover, $\widetilde{T}$ is the unique map $F:V/U\longrightarrow W$
with $T=F\circ\pi$: for every coset, this identity forces
$F(v+U)=F(\pi(v))=T(v)$.
\end{claim_box}

\begin{problem_box}
\textbf{Corollary:}
Restricting the
\emph{codomain} of $\widetilde{T}$ to $\operatorname{range}T$ gives an
isomorphism
\[
V/\ker T\cong\operatorname{range}T.
\]
\end{problem_box}

\begin{claim_box}
\textit{Proof:} Indeed, the induced map is linear and injective by the
proposition,
and is surjective onto this codomain by the equality of ranges.
With its original codomain $W$, $\widetilde{T}$ is surjective (and hence
an isomorphism) if and only if $T$ is surjective. Passing to the quotient
removes the failure of injectivity, but does not enlarge the range.
\hfill$\square$
\end{claim_box}

\begin{claim_box}
\textbf{Alternative proof of rank-nullity (using quotient spaces):}
Suppose now that $V$ is finite-dimensional. Apply the quotient-space
dimension theorem, using $U=\ker T$:
\[
\dim(V/\ker T)+\dim(\ker T)=\dim V.
\]
Since isomorphic spaces have the same dimension,
$\dim(V/\ker T)=\dim(\operatorname{range}T)$. Substituting yields
\[
\dim(\operatorname{range}T)+\dim(\ker T)=\dim V,
\]
which is the rank-nullity theorem. \hfill$\square$
\end{claim_box}

\begin{example_box}
\textbf{Example:} For the earlier map
$T_A:\mathbb{R}^3\longrightarrow\mathbb{R}^2$ given by
$T_A(x,y,z)=(x+2y,-y+z)$, we found
$\ker T_A=\operatorname{span}\{(-2,1,1)\}$. Its nullity is $1$, so
\[
\dim(\operatorname{range}T_A)=3-1=2.
\]
Since its range is a two-dimensional subspace of $\mathbb{R}^2$, the
range is all of $\mathbb{R}^2$ and $T_A$ is surjective.
\end{example_box}

\begin{example_box}
\textbf{Example:} Every linear map from $\mathbb{R}^3$ to
$\mathbb{R}^2$ has a null space of dimension at least $1$, so it sends
some nonzero vector to zero and cannot be injective.
\end{example_box}

\begin{example_box}
\textbf{Example:} Every linear map from $\mathbb{R}^2$ to
$\mathbb{R}^3$ has a range of dimension at most $2$, and therefore cannot
be surjective. For instance, the inclusion
\[
J:\mathbb{R}^2\longrightarrow\mathbb{R}^3,
\qquad J(x,y)=(x,y,0)
\]
has the $xy$-plane as its range and does not attain $(0,0,1)$.
\end{example_box}

\begin{remark_box}
\textbf{Note:} The dimension obstructions above require linearity;
they are not statements about arbitrary maps between the underlying
sets.
\end{remark_box}

\subsection{Graph of a Linear Transformation}

\begin{definition_box}
\textbf{Definition:} Let
$T:V\longrightarrow W$ be linear over $\mathbb{F}$. The \emph{graph}
of $T$ is
\[
\Gamma(T)=\{(v,T(v)):v\in V\}\subseteq V\times W.
\]
Here $V\times W$ is a vector space with componentwise operations:
\[
(v,w)+(v',w')=(v+v',w+w'),\qquad a(v,w)=(av,aw).
\]
Unlike the range, the graph records both the input and its output.
\end{definition_box}

\begin{problem_box}
\textbf{Proposition:} A map
$f:V\longrightarrow W$ is linear if and only if its graph
$\Gamma(f)$ is a vector subspace of $V\times W$.

\end{problem_box}

\begin{claim_box}
\textit{Proof:} If $f$ is linear, then $(0_V,0_W)$ belongs to its
graph. For $u,v\in V$ and $a,b\in\mathbb{F}$,
\[
a(u,f(u))+b(v,f(v))=(au+bv,af(u)+bf(v))
=(au+bv,f(au+bv)),
\]
which also belongs to the graph. Thus the graph is a subspace.

Conversely, suppose $\Gamma(f)$ is a subspace. The same linear
combination $(au+bv,af(u)+bf(v))$ belongs to $\Gamma(f)$.
The unique pair in the graph with first component $au+bv$ has second
component $f(au+bv)$. Hence
$f(au+bv)=af(u)+bf(v)$, proving linearity. \hfill$\square$
\end{claim_box}

\begin{problem_box}
\textbf{Proposition:}
For a linear map $T:V\longrightarrow W$, the map
\[
G:V\longrightarrow\Gamma(T),\qquad G(v)=(v,T(v)),
\]
is an isomorphism. In particular, if $V$ is finite-dimensional, then
$\dim\Gamma(T)=\dim V$.

\end{problem_box}

\begin{claim_box}
\textit{Proof:} Linearity follows from the componentwise operations
and the linearity of $T$. Equality $G(u)=G(v)$ forces $u=v$ by
comparing first components, so $G$ is injective. Every element of
$\Gamma(T)$ has the form $G(v)$, so $G$ is surjective. Its inverse
is the first-coordinate projection $(v,T(v))\mapsto v$.
The dimension statement follows because isomorphisms preserve
dimension. \hfill$\square$
\end{claim_box}

\begin{problem_box}
\textbf{Proposition:} The
second-coordinate projection restricted to the graph is the linear map
\[
Q:\Gamma(T)\longrightarrow W,\qquad Q(v,T(v))=T(v).
\]
It satisfies
\[
\ker Q=\{(v,0_W):v\in\ker T\},\qquad
\operatorname{range}Q=\operatorname{range}T.
\]
Equivalently,
\[
\Gamma(T)\cap(V\times\{0_W\})=(\ker T)\times\{0_W\}.
\]
Thus $T$ is injective exactly when this intersection is
$\{(0_V,0_W)\}$, and $T$ is surjective exactly when $Q$ is surjective.

\end{problem_box}

\begin{claim_box}
\textit{Proof:} Projection onto the second component is linear.
A pair $(v,T(v))$ is in its kernel precisely when $T(v)=0_W$,
and its possible outputs are precisely the values attained by $T$.
These observations prove both identities and the intersection formula.
The final assertions follow from the kernel criterion for injectivity
and the definition of surjectivity. \hfill$\square$
\end{claim_box}

\begin{problem_box}
\textbf{Proposition:} We have
\[
V\times W=\Gamma(T)\oplus(\{0_V\}\times W).
\]
Here the direct sum means that each pair is uniquely the sum of one
vector in each of these two subspaces.

\end{problem_box}

\begin{claim_box}
\textit{Proof:} Every $(v,w)\in V\times W$ can be written as
\[
(v,w)=(v,T(v))+(0_V,w-T(v)).
\]
The first summand belongs to the graph and the second to
$\{0_V\}\times W$. If a graph element $(v,T(v))$ also lies in
$\{0_V\}\times W$, then $v=0_V$ and $T(v)=0_W$. The intersection
is therefore the zero subspace, making the decomposition unique.
\hfill$\square$
\end{claim_box}

\begin{example_box}
\textbf{Example:} For $T:\mathbb{R}^2\longrightarrow\mathbb{R}$,
$T(x,y)=x+2y$, identify $\mathbb{R}^2\times\mathbb{R}$ with
$\mathbb{R}^3$. Then
\[
\Gamma(T)=\{(x,y,x+2y):x,y\in\mathbb{R}\}
=\operatorname{span}\{(1,0,1),(0,1,2)\}.
\]
The graph is the plane $z=x+2y$, of dimension $2$, while the range is
$\mathbb{R}$, of dimension $1$. Its intersection with $z=0$ is the
line $\{(-2t,t,0):t\in\mathbb{R}\}$, corresponding to
$\ker T=\operatorname{span}\{(-2,1)\}$.
\end{example_box}

\section{Matrices}

\begin{definition_box}
\textbf{Definition:} Let $m,n$ be positive integers. An
\emph{$m\times n$ matrix over $\mathbb{F}$} is a rectangular array of
elements of $\mathbb{F}$ arranged in $m$ rows and $n$ columns:
\[
A=(a_{ij})=
\begin{pmatrix}
a_{11}&a_{12}&\cdots&a_{1n}\\
a_{21}&a_{22}&\cdots&a_{2n}\\
\vdots&\vdots&\ddots&\vdots\\
a_{m1}&a_{m2}&\cdots&a_{mn}
\end{pmatrix}.
\]
The scalar $a_{ij}$ is the \emph{entry} in row $i$ and column $j$.
The size is written with the number of rows first. The set of all such
matrices is denoted by $\mathbb{F}^{m\times n}$ or
$M_{m\times n}(\mathbb{F})$. Two matrices are equal if they have the
same size and all corresponding entries are equal.
\end{definition_box}

\begin{example_box}
\textbf{Example:} The matrix
\[
B=\begin{pmatrix}1&2&0\\3&-1&4\end{pmatrix}
\in\mathbb{R}^{2\times3}
\]
has two rows and three columns. For instance, $b_{21}=3$ and $b_{23}=4$.
Its first row is $(1,2,0)$, and its second column is
$\begin{pmatrix}2\\-1\end{pmatrix}$.
A matrix with $m=n$ is called \emph{square}; a $1\times n$ matrix is a
\emph{row vector}, and an $m\times1$ matrix is a \emph{column vector}.
The \emph{zero matrix} of a given size has every entry equal to zero.
\end{example_box}

\begin{claim_box}
\textbf{Definition:}
A \emph{square matrix of order $n$} has $n$ rows and $n$ columns.
For $A=(a_{ij})\in M_{n\times n}(\mathbb{F})$, its \emph{main diagonal}
consists of $a_{11},a_{22},\ldots,a_{nn}$. We call $A$:
\begin{enumerate}
    \item \emph{Diagonal} if $a_{ij}=0$ whenever $i\neq j$.
    We write $A=\operatorname{diag}(d_1,\ldots,d_n)$ when $a_{ii}=d_i$.
    \item \emph{Upper triangular} if $a_{ij}=0$ whenever $i>j$;
    that is, every entry below the main diagonal is zero.
    \item \emph{Lower triangular} if $a_{ij}=0$ whenever $i<j$;
    that is, every entry above the main diagonal is zero.
    \item A \emph{scalar matrix} if it is diagonal and all its
    diagonal entries are the same scalar $c$. When $c=1$, it is the
    \emph{identity matrix} $I_n$.
\end{enumerate}
Diagonal entries may be zero. A matrix is diagonal if and only if it is
both upper and lower triangular: the two triangular conditions together
say exactly that every off-diagonal entry vanishes. An upper or lower
triangular matrix is called \emph{strictly triangular} of the respective
type if its diagonal entries are also zero.
\end{claim_box}

\begin{example_box}
\textbf{Examples:} Over $\mathbb{R}$, the matrices
\[
D=\begin{pmatrix}2&0&0\\0&0&0\\0&0&-1\end{pmatrix},\qquad
U=\begin{pmatrix}1&2&3\\0&0&4\\0&0&5\end{pmatrix},\qquad
L=\begin{pmatrix}1&0&0\\2&3&0\\4&5&6\end{pmatrix}
\]
are diagonal, upper triangular, and lower triangular, respectively.
The square zero matrix is both diagonal and strictly upper and lower
triangular. In these notes, triangular matrices are square; a rectangular
matrix whose entries satisfy $a_{ij}=0$ for $i>j$ is instead called
\emph{upper trapezoidal}.
\end{example_box}

\begin{claim_box}
\textbf{Definition:} Let
$A=(a_{ij})$ and $B=(b_{ij})$ belong to
$M_{m\times n}(\mathbb{F})$, and let $c\in\mathbb{F}$. Define
\[
A+B=(a_{ij}+b_{ij}),\qquad cA=(ca_{ij}).
\]
Thus addition adds corresponding entries, and scalar multiplication
multiplies every entry by the same scalar. Both operations preserve the
size of the matrix; addition requires the matrices to have the same size.
\end{claim_box}

\begin{example_box}
\textbf{Example:} Over $\mathbb{R}$,
\[
\begin{pmatrix}1&2\\0&-1\end{pmatrix}
+\begin{pmatrix}3&0\\4&2\end{pmatrix}
=\begin{pmatrix}4&2\\4&1\end{pmatrix},
\qquad
3\begin{pmatrix}1&2\\0&-1\end{pmatrix}
=\begin{pmatrix}3&6\\0&-3\end{pmatrix}.
\]
Under these operations, $M_{m\times n}(\mathbb{F})$ is a vector space:
the axioms hold entrywise because they hold in $\mathbb{F}$. Its zero
vector is the zero matrix, and the additive inverse of $A$ is $-A$.
\end{example_box}

\begin{definition_box}
\textbf{Definition:} Let
$A=(a_{ik})\in M_{m\times n}(\mathbb{F})$ and
$B=(b_{kj})\in M_{n\times p}(\mathbb{F})$. Their \emph{product}
$AB$ is the $m\times p$ matrix whose entries are
\[
{(AB)}_{ij}=\sum_{k=1}^n a_{ik}b_{kj},
\qquad 1\leq i\leq m,\quad 1\leq j\leq p.
\]
To find the entry in row $i$, column $j$, multiply corresponding entries
of row $i$ of $A$ and column $j$ of $B$, then add the products.
The number of columns of $A$ must equal the number of rows of $B$:
\[
(m\times n)(n\times p)\longrightarrow m\times p.
\]
This is not entrywise multiplication.
\end{definition_box}

\begin{example_box}
\textbf{Example:} Over $\mathbb{R}$,
let
\[
A=\begin{pmatrix}1&2&0\\0&-1&1\end{pmatrix},
\qquad B=\begin{pmatrix}1&0\\2&1\\3&4\end{pmatrix}.
\]
Since $A$ is $2\times3$ and $B$ is $3\times2$, the product $AB$ is
$2\times2$. Computing each entry gives
\[
\begin{aligned}
AB
&=\begin{pmatrix}
1\cdot1+2\cdot2+0\cdot3 & 1\cdot0+2\cdot1+0\cdot4\\
0\cdot1+(-1)\cdot2+1\cdot3 & 0\cdot0+(-1)\cdot1+1\cdot4
\end{pmatrix}\\[4pt]
&=\begin{pmatrix}5&2\\1&3\end{pmatrix}.
\end{aligned}
\]
\end{example_box}

\begin{remark_box}
\textbf{Note:} The reverse product $BA$ need not
be defined. Even when both products exist, they need not be equal;
in this example $BA$ is $3\times3$, whereas $AB$ is $2\times2$.
\end{remark_box}

\begin{claim_box}
\textbf{Definition:} The \emph{identity matrix}
$I_n\in\mathbb{F}^{n\times n}$ has $1$ in every diagonal position
and $0$ elsewhere. Write $0_{m\times n}$ for the zero matrix of size
$m\times n$.
\end{claim_box}

\begin{claim_box}
\textbf{Properties of matrix addition and scalar multiplication:}
For $A,B,C\in\mathbb{F}^{m\times n}$ and $a,b\in\mathbb{F}$,
\[
\begin{aligned}
A+B&=B+A, & (A+B)+C&=A+(B+C),\\
A+0_{m\times n}&=A, & A+(-A)&=0_{m\times n},\\
a(A+B)&=aA+aB, & (a+b)A&=aA+bA,\\
a(bA)&=(ab)A, & 1A&=A,\qquad 0A=0_{m\times n}.
\end{aligned}
\]
These identities hold entry by entry because the corresponding
identities hold in the field $\mathbb{F}$.
\end{claim_box}

\begin{claim_box}
\textbf{Properties of matrix multiplication:} All products and
sums below require the indicated compatible sizes.
\begin{enumerate}
    \item \emph{Associativity.} If $A$ is $m\times n$, $B$ is
    $n\times p$, and $C$ is $p\times q$, then
    \[
    (AB)C=A(BC)\in\mathbb{F}^{m\times q}.
    \]
    Thus parentheses may be omitted from a product $ABC$, while
    the order of the factors must be retained.
    \item \emph{Distributivity.} If $A$ is $m\times n$ and $B,C$
    are $n\times p$, then
    \[
    A(B+C)=AB+AC.
    \]
    If $A,B$ are $m\times n$ and $C$ is $n\times p$, then
    \[
    (A+B)C=AC+BC.
    \]
    \item \emph{Compatibility with scalars.} If $A$ is $m\times n$,
    $B$ is $n\times p$, and $c\in\mathbb{F}$, then
    \[
    c(AB)=(cA)B=A(cB).
    \]
    \item \emph{Identity and zero.} For $A\in\mathbb{F}^{m\times n}$,
    \[
    I_m A=A=A I_n,\qquad
    A0_{n\times p}=0_{m\times p},\qquad
    0_{q\times m}A=0_{q\times n}.
    \]
\end{enumerate}

\textit{Proof:} For associativity, the entry in row $i$, column $j$ is
\[
\begin{aligned}
{((AB)C)}_{ij}
&=\sum_{k=1}^p\left(\sum_{\ell=1}^n a_{i\ell}b_{\ell k}\right)c_{kj}\\
&=\sum_{\ell=1}^n a_{i\ell}
  \left(\sum_{k=1}^p b_{\ell k}c_{kj}\right)
={\bigl(A(BC)\bigr)}_{ij}.
\end{aligned}
\]
Distributivity follows by expanding each entry, for example
\[
{(A(B+C))}_{ij}
=\sum_{k=1}^n a_{ik}(b_{kj}+c_{kj})
={(AB)}_{ij}+{(AC)}_{ij}.
\]
Expanding $(a_{ik}+b_{ik})c_{kj}$ proves the other distributive law.
A scalar factor can be moved outside each finite sum, proving scalar
compatibility. Finally, multiplication by an identity matrix selects
the corresponding row or column of $A$, while multiplication by a zero
matrix gives only zero summands. These entrywise equalities prove all
the identities. \hfill$\square$
\end{claim_box}

\begin{claim_box}
\textbf{Definition:} Let
$A\in\mathbb{F}^{m\times n}$. A matrix $L\in\mathbb{F}^{n\times m}$
is a \emph{left inverse} of $A$ if $LA=I_n$. A matrix
$R\in\mathbb{F}^{n\times m}$ is a \emph{right inverse} of $A$ if
$AR=I_m$. The side refers to the position of the inverse in the product.
Both kinds of inverse necessarily have size $n\times m$; their having
the same size alone does not imply their existence or equality.
\end{claim_box}

\begin{definition_box}
\textbf{Definition:} A square matrix
$A\in\mathbb{F}^{n\times n}$ is \emph{invertible}, or
\emph{nonsingular}, if there is a matrix $B\in\mathbb{F}^{n\times n}$
such that
\[
AB=BA=I_n.
\]
Such a $B$ is a \emph{two-sided inverse} of $A$. We will prove that it
is unique and denote it by $A^{-1}$. A square matrix with no inverse is
called \emph{singular}. Matrix inversion is not entrywise reciprocation.
\end{definition_box}

\begin{problem_box}
\textbf{Theorem:} If a matrix
$A\in\mathbb{F}^{m\times n}$ admits both a left inverse $L$ and a
right inverse $R$, then $L=R$. Moreover, $m=n$, so $A$ is invertible
and this common inverse is unique.

\end{problem_box}

\begin{claim_box}
\textit{Proof:} Associativity and the identity laws give
\[
L=L I_m=L(AR)=(LA)R=I_n R=R.
\]
To see why $A$ must be square, consider
$T_A:\mathbb{F}^n\longrightarrow\mathbb{F}^m$, $T_A(x)=Ax$.
If $Ax=0$, then $x=LAx=0$, so $T_A$ is injective and $n\leq m$.
For every $y\in\mathbb{F}^m$, we have $y=A(Ry)$, so $T_A$ is
surjective and $m\leq n$. Thus $m=n$.
Finally, any other left inverse $L'$ equals $R$ by the same calculation,
and any other right inverse $R'$ equals $L$. In particular, any two
two-sided inverses coincide. \hfill$\square$
\end{claim_box}

\begin{problem_box}
\textbf{Theorem:}
Let
$A\in\mathbb{F}^{n\times n}$. Then
\[
A\text{ has a left inverse}
\quad\Longleftrightarrow\quad
A\text{ has a right inverse}.
\]
Whenever these equivalent conditions hold, $A$ is invertible and
every one-sided inverse equals the unique two-sided inverse $A^{-1}$.
\end{problem_box}

\begin{claim_box}
\textit{Proof:} $(\Rightarrow)$ Suppose $L$ is a left inverse,
so $LA=I_n$.
If $Ax=0$, then
\[
x=I_n x=(LA)x=L(Ax)=0.
\]
Thus $T_A:\mathbb{F}^n\longrightarrow\mathbb{F}^n$, $T_A(x)=Ax$,
is injective. Since the domain and codomain have the same finite
dimension, rank-nullity implies that $T_A$ is also surjective.
For any $y\in\mathbb{F}^n$, choose $x$ with $Ax=y$. Then
\[
ALy=AL(Ax)=A(LA)x=Ax=y.
\]
Hence $AL$ acts as the identity on every vector, and evaluating on
the standard basis gives $AL=I_n$. Therefore the very same matrix $L$
is also a right inverse, and $A^{-1}=L$.

$(\Leftarrow)$ Suppose $R$ is a right inverse, so $AR=I_n$. Every
$y\in\mathbb{F}^n$ satisfies $y=A(Ry)$, making $T_A$ surjective
and hence injective by equal finite dimensions. For every $x$,
\[
A(RAx-x)=(AR)Ax-Ax=0.
\]
Injectivity implies $RAx-x=0$. Thus $RA=I_n$, so $R$ is also a
left inverse.

This proves both implications. In either case the given inverse is
two-sided, and uniqueness follows from the preceding theorem.
\hfill$\square$
\end{claim_box}

\begin{problem_box}
\textbf{Theorem:} Let
$A\in\mathbb{F}^{n\times n}$, and let $A_1,\ldots,A_n\in\mathbb{F}^n$
be its columns. Then $A$ is invertible if and only if
$A_1,\ldots,A_n$ form a basis of $\mathbb{F}^n$.
\end{problem_box}

\begin{claim_box}
\textit{Proof:} Suppose first that $A$ is invertible, with inverse
$B$. For any $y\in\mathbb{F}^n$,
\[
y=AB y=\sum_{j=1}^n (By)_j A_j,
\]
so the columns of $A$ span $\mathbb{F}^n$. If
$\sum_{j=1}^n c_jA_j=0$, then $Ac=0$, where
$c=(c_1,\ldots,c_n)^{\mathsf T}$. Multiplying by $B$ gives $c=0$,
so the columns are linearly independent. They therefore form a basis.

Conversely, suppose the columns of $A$ form a basis of
$\mathbb{F}^n$. For each standard basis vector $e_j$, choose the unique
$x_j\in\mathbb{F}^n$ such that $Ax_j=e_j$. Let $B$ be the matrix whose
$j$th column is $x_j$. Then $AB=I_n$. Also, if $Ax=0$, the linear
independence of the columns of $A$ implies $x=0$. For any $x\in\mathbb{F}^n$,
\[
A(BAx-x)=(AB)Ax-Ax=Ax-Ax=0.
\]
Thus $BAx-x=0$, so $BA=I_n$. Hence $B$ is a two-sided inverse of $A$,
and $A$ is invertible. \hfill$\square$
\end{claim_box}

\begin{problem_box}
\textbf{Theorem:} Let $M=(m_{ij})\in\mathbb{F}^{n\times n}$ be upper
triangular. Then $M$ is invertible if and only if
\[
m_{ii}\neq 0\qquad\text{for every }1\leq i\leq n.
\]
Moreover, when $M$ is invertible, the diagonal entries of its inverse are
\[
(M^{-1})_{ii}=m_{ii}^{-1},\qquad 1\leq i\leq n.
\]
\end{problem_box}

\begin{claim_box}
\textit{Proof:} Suppose first that some diagonal entry is zero, and let
$k$ be the smallest index with $m_{kk}=0$. Set $x_k=1$ and $x_j=0$
for $j>k$. Since $m_{ii}\neq0$ for $i<k$, define the remaining
coordinates recursively, for $i=k-1,k-2,\ldots,1$, by
\[
x_i=-m_{ii}^{-1}\sum_{j=i+1}^{k}m_{ij}x_j.
\]
The first $k$ rows of $Mx$ are then zero by construction, and rows
$i>k$ are zero because $M$ is upper triangular and $x_j=0$ for $j>k$.
Thus $Mx=0$ with $x_k=1$, so $M$ is not invertible.

Conversely, if every $m_{ii}$ is nonzero, then for every
$b\in\mathbb{F}^n$ the system $Mx=b$ has a unique solution, obtained
by solving the equations from the last row upward:
\[
x_i=m_{ii}^{-1}\left(b_i-\sum_{j=i+1}^{n}m_{ij}x_j\right).
\]
Therefore $M$ is invertible.

Write $N=M^{-1}$. For each $j$, the $j$th column of $N$ is the unique
solution to $Mx=e_j$. Solving from the last row upward shows that its
entries in rows $i>j$ are zero, so $N$ is upper triangular. The $i$th
diagonal entry of $MN=I_n$ is consequently
\[
1=(MN)_{ii}=\sum_{j=i}^{n}m_{ij}n_{ji}=m_{ii}n_{ii}.
\]
Hence $(M^{-1})_{ii}=n_{ii}=m_{ii}^{-1}$ for every $i$.
\hfill$\square$
\end{claim_box}

\begin{remark_box}
\textbf{Caution:}
For rectangular matrices, a left inverse alone or a right inverse alone
need not be unique. For example, for every $t\in\mathbb{F}$,
\[
\begin{pmatrix}1&t\end{pmatrix}\begin{pmatrix}1\\0\end{pmatrix}
=I_1,\qquad
\begin{pmatrix}1&0\end{pmatrix}\begin{pmatrix}1\\t\end{pmatrix}
=I_1.
\]
Thus the column matrix on the left has distinct left inverses, and
the row matrix on the right has distinct right inverses (take $t=0,1$).
Neither rectangular matrix can have both kinds of inverse.
\end{remark_box}

\subsection{Block matrices and Schur complements}

\begin{definition_box}
\textbf{Definition (block matrices and Schur complements):} A matrix
partitioned into compatible submatrices is called a \emph{block matrix}.
For
\[
M=\begin{pmatrix}A&B\\C&D\end{pmatrix},
\]
if $A$ is invertible, define
\[
S_A=D-CA^{-1}B,
\]
called the \emph{Schur complement of $A$ in $M$}. If $D$ is invertible,
define
\[
S_D=A-BD^{-1}C,
\]
called the \emph{Schur complement of $D$ in $M$}.
\end{definition_box}

\begin{problem_box}
\textbf{Theorem:} Let
\[
M=\begin{pmatrix}A&B\\C&D\end{pmatrix},
\]
where $A\in\mathbb{F}^{m\times m}$, $B\in\mathbb{F}^{m\times n}$,
$C\in\mathbb{F}^{n\times m}$, and $D\in\mathbb{F}^{n\times n}$.
Then:
\begin{enumerate}
    \item If $A$ is invertible, then $M$ is invertible if and only if
    $S_A=D-CA^{-1}B$ is invertible.
    \item If $D$ is invertible, then $M$ is invertible if and only if
    $S_D=A-BD^{-1}C$ is invertible.
    \item If $M$ is invertible and
    \[
    M^{-1}=\begin{pmatrix}P&Q\\R&S\end{pmatrix},
    \]
    then $A$ is invertible if and only if $S$ is invertible.
    \item Under the same assumption, $D$ is invertible if and only if
    $P$ is invertible.
    \item If $C=0$, then $M$ is invertible if and only if both $A$ and
    $D$ are invertible. In that case,
    \[
    M^{-1}=\begin{pmatrix}
    A^{-1}&-A^{-1}BD^{-1}\\
    0&D^{-1}
    \end{pmatrix}.
    \]
\end{enumerate}
\end{problem_box}

\begin{claim_box}
\textit{Proof:} For part 1, set $H=D-CA^{-1}B$. We have the factorization
\[
M=
\begin{pmatrix}I_m&0\\CA^{-1}&I_n\end{pmatrix}
\begin{pmatrix}A&B\\0&H\end{pmatrix}.
\]
The first factor is invertible. Since $A$ is invertible, the second
factor is invertible exactly when $H$ is invertible: if the second
factor is invertible and $Hy=0$, applying it to
$\begin{pmatrix}-A^{-1}By\\y\end{pmatrix}$ gives zero, so $y=0$;
thus $H$ is injective and hence invertible. Conversely, if $H$ is
invertible, the second factor has inverse
\[
\begin{pmatrix}A^{-1}&-A^{-1}BH^{-1}\\0&H^{-1}\end{pmatrix}.
\]
In this case multiplication of the factor inverses gives
\[
M^{-1}=
\begin{pmatrix}
A^{-1}+A^{-1}BH^{-1}CA^{-1}&-A^{-1}BH^{-1}\\
-H^{-1}CA^{-1}&H^{-1}
\end{pmatrix}.
\]

For part 2, set $G=A-BD^{-1}C$. Then
\[
M=
\begin{pmatrix}I_m&BD^{-1}\\0&I_n\end{pmatrix}
\begin{pmatrix}G&0\\C&D\end{pmatrix}.
\]
The first factor is invertible. Since $D$ is invertible, the second
factor is invertible exactly when $G$ is invertible: if it is invertible
and $Gx=0$, applying it to
$\begin{pmatrix}x\\-D^{-1}Cx\end{pmatrix}$ gives zero, so $x=0$;
thus $G$ is invertible. Conversely, when $G$ is invertible, the second
factor has inverse
\[
\begin{pmatrix}G^{-1}&0\\-D^{-1}CG^{-1}&D^{-1}\end{pmatrix}.
\]
Multiplying the factor inverses yields
\[
M^{-1}=
\begin{pmatrix}
G^{-1}&-G^{-1}BD^{-1}\\
-D^{-1}CG^{-1}&D^{-1}+D^{-1}CG^{-1}BD^{-1}
\end{pmatrix}.
\]

For part 3, if $A$ is invertible, part 1 and the displayed inverse
formula show that $S=(D-CA^{-1}B)^{-1}$, so $S$ is invertible.
Conversely, suppose $S$ is invertible and $Ax=0$. Since $M$ is
invertible,
\[
\begin{pmatrix}x\\0\end{pmatrix}
=M^{-1}\begin{pmatrix}0\\Cx\end{pmatrix}
=\begin{pmatrix}QCx\\SCx\end{pmatrix}.
\]
Thus $SCx=0$, so $Cx=0$. It follows that
$M\begin{pmatrix}x\\0\end{pmatrix}=0$, and invertibility of $M$
gives $x=0$. Hence $A$ is injective and therefore invertible.

For part 4, if $D$ is invertible, part 2 and its inverse formula show
that $P=(A-BD^{-1}C)^{-1}$, so $P$ is invertible. Conversely,
if $P$ is invertible and $Dy=0$, then
\[
\begin{pmatrix}0\\y\end{pmatrix}
=M^{-1}\begin{pmatrix}By\\0\end{pmatrix}
=\begin{pmatrix}PBy\\RBy\end{pmatrix}.
\]
Therefore $PBy=0$, so $By=0$. Hence
$M\begin{pmatrix}0\\y\end{pmatrix}=0$, and invertibility of $M$
implies $y=0$. Thus $D$ is injective and therefore invertible.

For part 5, if $M$ is invertible, its injectivity implies that $A$ is
injective: $Ax=0$ gives $M\begin{pmatrix}x\\0\end{pmatrix}=0$.
Its surjectivity implies that $D$ is surjective: for any $y\in\mathbb{F}^n$,
write $M\begin{pmatrix}x\\z\end{pmatrix}=\begin{pmatrix}0\\y\end{pmatrix}$;
then $Dz=y$. Since $A$ and $D$ are square, both are invertible.
Conversely, if $A$ and $D$ are invertible and $C=0$, direct
multiplication verifies the stated formula for $M^{-1}$.
\hfill$\square$
\end{claim_box}


\subsection{Further matrix topics}


\begin{example_box}
\textbf{Example:} Unlike
scalars in a field, nonzero matrices can multiply to zero. For example,
\[
A=\begin{pmatrix}1&0\\0&0\end{pmatrix},\qquad
B=\begin{pmatrix}0&0\\1&0\end{pmatrix}
\quad\Longrightarrow\quad
AB=\begin{pmatrix}0&0\\0&0\end{pmatrix}.
\]
Both $A$ and $B$ are nonzero. Moreover, $BA=B\neq0$, so $AB=0$
does not imply $BA=0$, even when both products are defined.
\end{example_box}

\begin{problem_box}
\textbf{Theorem:} Let
$A\in\mathbb{F}^{m\times n}$ and $B\in\mathbb{F}^{n\times p}$.
Write $T_A(x)=Ax$ and $T_B(y)=By$. Then
\[
AB=0_{m\times p}
\quad\Longleftrightarrow\quad
\operatorname{range}T_B\subseteq\ker T_A.
\]
Equivalently, every column of $B$ lies in the null space of $A$.
Consequently, if $AB=0$, then
\[
\operatorname{rank}A+\operatorname{rank}B\leq n.
\]

\end{problem_box}

\begin{claim_box}
\textit{Proof:} If $AB=0$, every vector $z=By$ in the range of $T_B$
satisfies $Az=ABy=0$, so $z\in\ker T_A$. Conversely, the inclusion
implies $ABy=0$ for every $y\in\mathbb{F}^p$. Taking $y$ to be each
standard basis vector shows that every column of $AB$ is zero.
The column formulation follows because column $j$ of $AB$ is $AB_j$,
where $B_j$ is column $j$ of $B$. Finally, the inclusion and
rank-nullity give
\[
\operatorname{rank}B=\dim(\operatorname{range}T_B)
\leq\dim(\ker T_A)=n-\operatorname{rank}A.
\]
Rearranging proves the inequality. \hfill$\square$
\end{claim_box}

\begin{problem_box}
\textbf{Proposition:} Let
$A\in\mathbb{F}^{m\times n}$ and $B\in\mathbb{F}^{n\times p}$.
Then
\[
\operatorname{rank}A+\operatorname{rank}B-n
\leq\operatorname{rank}(AB)
\leq\min(\operatorname{rank}A,\operatorname{rank}B).
\]
\end{problem_box}

\begin{claim_box}
\textit{Proof:} The range of $T_{AB}$ is $A(\operatorname{range}T_B)$,
which is contained in $\operatorname{range}T_A$. Also, applying a linear
map cannot increase dimension, so
\[
\operatorname{rank}(AB)
=\dim\bigl(A(\operatorname{range}T_B)\bigr)
\leq\dim(\operatorname{range}T_B)=\operatorname{rank}B.
\]
Together these give the upper bound.

For the lower bound, restrict $T_A$ to $\operatorname{range}T_B$.
Its image is $\operatorname{range}T_{AB}$, and its kernel is
$\operatorname{range}T_B\cap\ker T_A$. Rank-nullity for this restriction
therefore gives
\[
\operatorname{rank}(AB)
=\operatorname{rank}B-
\dim(\operatorname{range}T_B\cap\ker T_A)
\geq\operatorname{rank}B-\dim(\ker T_A)
=\operatorname{rank}B-(n-\operatorname{rank}A).
\]
Rearranging proves the lower bound. \hfill$\square$
\end{claim_box}

\begin{problem_box}
\textbf{Proposition:} For the same
matrix sizes, suppose $AB=0$.
\begin{enumerate}
    \item If $A$ has a left inverse, then $B=0$.
    In particular, this holds if $A$ is square and invertible.
    \item If $B$ has a right inverse, then $A=0$.
    In particular, this holds if $B$ is square and invertible.
\end{enumerate}

\end{problem_box}

\begin{claim_box}
\textit{Proof:} If $LA=I_n$, then
\[
B=I_n B=(LA)B=L(AB)=0.
\]
If $BR=I_n$, then
\[
A=A I_n=A(BR)=(AB)R=0.
\]
Thus cancellation is justified by an inverse on the appropriate side,
not merely by a factor being nonzero. \hfill$\square$
\end{claim_box}

\begin{problem_box}
\textbf{Corollary:} If $A$ has a left inverse
and $AB=AC$, where $B,C$ have the same compatible size, then $B=C$.
If $A$ has a right inverse and $BA=CA$, then $B=C$.

\end{problem_box}

\begin{claim_box}
\textit{Proof:} Distributivity gives $A(B-C)=0$ in the first case
and $(B-C)A=0$ in the second. Apply the preceding proposition to
obtain $B-C=0$. \hfill$\square$
\end{claim_box}

\begin{problem_box}
\textbf{Proposition:}
Let $A\in\mathbb{F}^{n\times n}$ with $n\geq1$. Then $A$ is singular
if and only if there exists a nonzero $B\in\mathbb{F}^{n\times n}$
such that $AB=0$.

\end{problem_box}

\begin{claim_box}
\textit{Proof:} If $A$ is singular, $T_A$ cannot be injective:
an injective operator on $\mathbb{F}^n$ is bijective, and its linear
inverse would provide an inverse matrix for $A$. Choose a nonzero
$v\in\ker T_A$ and let $B$ have first column $v$ and all other
columns zero. Then $B\neq0$ and every column of $AB$ is zero.
Conversely, if $AB=0$ with $B\neq0$, $A$ cannot be invertible by
the zero-factor proposition. Hence $A$ is singular. \hfill$\square$
\end{claim_box}

\begin{claim_box}
\textbf{Consequence:} If $A,B$ are square matrices of the same
order, $AB=0$, and both factors are nonzero, then both factors are
singular. Indeed, invertibility of either would force the other to be
zero. The converse is false: two singular matrices need not have zero
product. For instance,
$P=\begin{pmatrix}1&0\\0&0\end{pmatrix}$ is singular but $P^2=P\neq0$.
\end{claim_box}

\begin{problem_box}
\textbf{Theorem:}
Let
\[
A=\begin{pmatrix}a&b\\c&d\end{pmatrix},\qquad a,b,c,d\in\mathbb{F}.
\]
The scalar $\det A=ad-bc$ is called the \emph{determinant} of this
$2\times2$ matrix. Then $A$ is invertible if and only if $ad-bc\neq0$,
and in that case
\[
\boxed{A^{-1}=\frac{1}{ad-bc}\begin{pmatrix}d&-b\\-c&a\end{pmatrix}.}
\]
To compute it, interchange the diagonal entries, negate the two
off-diagonal entries, and multiply by ${(ad-bc)}^{-1}$.
\end{problem_box}

\begin{claim_box}
\textit{Proof:} Put $C=\begin{pmatrix}d&-b\\-c&a\end{pmatrix}$.
Direct multiplication gives
\[
AC=\begin{pmatrix}ad-bc&0\\0&ad-bc\end{pmatrix}
=CA=(ad-bc)I_2.
\]
If $ad-bc\neq0$, division by this scalar gives a two-sided inverse,
proving the formula.

If $ad-bc=0$, then $AC=0$. Were $A$ invertible, multiplying by
$A^{-1}$ on the left would give $C=0$. The entries of $C$ then force
$a=b=c=d=0$, so $A=0$, which cannot have an inverse because $0B=0\neq I_2$
for every $B$. This contradiction proves the converse. \hfill$\square$
\end{claim_box}

\begin{example_box}
\textbf{Example:} Over
$\mathbb{R}$, let
\[
A=\begin{pmatrix}2&1\\5&3\end{pmatrix}.
\]
Since $\det A=2\cdot3-1\cdot5=1$, the formula gives
\[
A^{-1}=\begin{pmatrix}3&-1\\-5&2\end{pmatrix}.
\]
Indeed,
\[
\begin{pmatrix}2&1\\5&3\end{pmatrix}
\begin{pmatrix}3&-1\\-5&2\end{pmatrix}
=\begin{pmatrix}1&0\\0&1\end{pmatrix}.
\]
In contrast, $\begin{pmatrix}1&2\\2&4\end{pmatrix}$ has determinant
zero and hence no inverse.
\end{example_box}

\begin{problem_box}
\textbf{Exercise:}
Let $n\geq1$ be an integer, and define $A=(a_{ij})\in\mathbb{R}^{n\times n}$
by
\[
a_{ij}=\frac{1}{i+j-1},\qquad 1\leq i,j\leq n.
\]
Prove that:
\begin{enumerate}
    \item $A$ is invertible for every $n\geq1$.
    \item Every entry of $A^{-1}$ is an integer; that is,
    $A^{-1}\in M_{n\times n}(\mathbb{Z})$.
\end{enumerate}
\end{problem_box}

\begin{claim_box}
\textbf{Solution:} We construct a matrix
$B$ with integer entries such that $AB=I_n$. Since $A$ is square,
this will prove both invertibility and $A^{-1}=B$. The only polynomial
fact we need is that a polynomial of degree at most $n-1$ with $n$
distinct roots is identically zero. Empty products below equal $1$,
and $0!=1$.

\emph{Construct one column of the inverse:} Fix
$j\in\{1,\ldots,n\}$. We want numbers $b_{1j},\ldots,b_{nj}$ such that
\[
\sum_{i=1}^n\frac{b_{ij}}{k+i-1}
=\begin{cases}1,&k=j,\\0,&k\neq j,\end{cases}
\qquad 1\leq k\leq n.
\]
These are exactly the equations for column $j$ of $AB=I_n$.
Introduce the polynomials
\[
Q(x)=\prod_{r=1}^n(x+r-1),\qquad
P_j(x)=Q(j)\prod_{\substack{1\leq k\leq n\\k\neq j}}
\frac{x-k}{j-k}.
\]
The polynomial $P_j$ has degree $n-1$, and its factors give
$P_j(j)=Q(j)$ and $P_j(k)=0$ for all other $k\in\{1,\ldots,n\}$.
Thus $P_j(k)/Q(k)$ has exactly the desired values.

\emph{Write this quotient as a sum:} Set
\[
b_{ij}=\frac{P_j(1-i)}{
\displaystyle\prod_{\substack{1\leq r\leq n\\r\neq i}}(r-i)}.
\]
Every denominator here is nonzero. We claim the polynomial identity
\[
P_j(x)=\sum_{i=1}^n b_{ij}
\prod_{\substack{1\leq r\leq n\\r\neq i}}(x+r-1).
\]
To check it, substitute $x=1-s$, where $1\leq s\leq n$.
All terms on the right vanish except $i=s$, and that term equals
$P_j(1-s)$ by the definition of $b_{sj}$. Both sides have degree
at most $n-1$, so their difference has $n$ distinct roots and is zero.
Dividing by $Q(x)$, whenever $Q(x)\neq0$, gives
\[
\frac{P_j(x)}{Q(x)}=\sum_{i=1}^n\frac{b_{ij}}{x+i-1}.
\]
In particular, substitution of $x=k\in\{1,\ldots,n\}$ proves the
equations in Step 1. Doing this for every $j$ constructs $B$ with
$AB=I_n$, so $A$ is invertible and $B=A^{-1}$.

\emph{Show that the entries are integers:} The products above
can be written using factorials:
\[
\begin{aligned}
Q(j)&=\frac{(n+j-1)!}{(j-1)!},\\
\prod_{\substack{1\leq k\leq n\\k\neq j}}(j-k)
&={(-1)}^{n-j}(j-1)!(n-j)!,\\
\prod_{\substack{1\leq k\leq n\\k\neq j}}(1-i-k)
&={(-1)}^{n-1}\frac{(n+i-1)!}{(i-1)!(i+j-1)},\\
\prod_{\substack{1\leq r\leq n\\r\neq i}}(r-i)
&={(-1)}^{i-1}(i-1)!(n-i)!.
\end{aligned}
\]
Substitution into the formula for $b_{ij}$ yields
\[
b_{ij}={(-1)}^{i+j}
\frac{(n+i-1)!(n+j-1)!}
{(i+j-1){((i-1)!)}^2{((j-1)!)}^2(n-i)!(n-j)!}.
\]
Although this expression is a quotient, it can be regrouped as
\[
\boxed{
b_{ij}={(-1)}^{i+j}(i+j-1)
\binom{n+i-1}{n-j}\binom{n+j-1}{n-i}
\binom{i+j-2}{i-1}^{\!2}.}
\]
Indeed, expanding the three binomial coefficients into factorials
recovers the preceding expression. All their upper and lower arguments
are nonnegative integers with the lower argument no larger than the
upper one. Each binomial coefficient is therefore an integer, so
$b_{ij}\in\mathbb{Z}$. Since $B=A^{-1}$, this completes the proof.
\hfill$\square$
\end{claim_box}

\begin{definition_box}
\textbf{Definition:} The
\emph{leading entry} of a nonzero row is its leftmost nonzero entry.
A matrix is in \emph{row echelon form} (REF) if:
\begin{enumerate}
    \item All zero rows occur below all nonzero rows.
    \item The leading entry of each nonzero row is strictly to the
    right of the leading entry of the preceding nonzero row.
    \item Every entry below a leading entry is zero.
\end{enumerate}
It is in \emph{reduced row echelon form} (RREF) if it also satisfies:
\begin{enumerate}
    \item Every leading entry is $1$.
    \item Each leading $1$ is the only nonzero entry in its column.
\end{enumerate}
Leading entries need not lie on the main diagonal: columns may be
skipped. The zero matrix satisfies these conditions as well.
\end{definition_box}

\begin{problem_box}
\textbf{Proposition:} Every
square matrix in REF is upper triangular. More generally, every
rectangular matrix in REF is upper trapezoidal.

\end{problem_box}

\begin{claim_box}
\textit{Proof:} Let $A$ have $r$ nonzero rows. They are the first $r$
rows because all zero rows are at the bottom. For $1\leq i\leq r$,
let $p_i$ be the column index of the leading entry in row $i$.
The echelon conditions give
\[
1\leq p_1<p_2<\cdots<p_r,
\]
so $p_i\geq i$. If $j<i$, then $j<p_i$, and hence $a_{ij}=0$
because all entries before the leading entry of row $i$ vanish.
For $i>r$, the entire row is zero. Thus $a_{ij}=0$ whenever $i>j$,
which proves both assertions. \hfill$\square$
\end{claim_box}

\begin{problem_box}
\textbf{Theorem:} If
$A\in M_{n\times n}(\mathbb{F})$ is in RREF and has no zero row,
then $A=I_n$.

\end{problem_box}

\begin{claim_box}
\textit{Proof:} Every row has a leading entry, so there are $n$
leading columns with indices
\[
1\leq p_1<p_2<\cdots<p_n\leq n.
\]
The only possibility is $p_i=i$ for every $i$. Thus every diagonal
entry is a leading $1$. Since $A$ is reduced, that $1$ is the only
nonzero entry in its column. Consequently, all off-diagonal entries
are zero, and $A=I_n$. \hfill$\square$
\end{claim_box}

\begin{remark_box}
\textbf{Caution:} An upper triangular
matrix need not be in REF\@: for example,
$\begin{pmatrix}0&1\\0&1\end{pmatrix}$ has its two leading entries
in the same column. Also, the square assumption in the preceding
theorem is essential. The rectangular matrix
\[
\begin{pmatrix}1&0&2\\0&1&3\end{pmatrix}
\]
is in RREF and has no zero row, but is not an identity matrix.
\end{remark_box}

\begin{definition_box}
\textbf{Definition:} The
\emph{transpose} of $A=(a_{ij})\in\mathbb{F}^{m\times n}$ is the
$n\times m$ matrix $A^{\mathsf{T}}$ defined by
\[
{(A^{\mathsf{T}})}_{ji}=a_{ij},
\qquad 1\leq i\leq m,\quad 1\leq j\leq n.
\]
It interchanges rows and columns without conjugating entries, even
over $\mathbb{C}$. The notation $A'$ is also used for the transpose;
here we use $A^{\mathsf{T}}$ to make its meaning explicit. For example,
\[
A=\begin{pmatrix}1&2&3\\4&5&6\end{pmatrix}
\quad\Longrightarrow\quad
A^{\mathsf{T}}=\begin{pmatrix}1&4\\2&5\\3&6\end{pmatrix}.
\]
Transposing twice restores the original matrix, and transposing an
identity matrix leaves it unchanged:
\[
{(A^{\mathsf{T}})}^{\mathsf{T}}=A,\qquad I_n^{\mathsf{T}}=I_n.
\]
Both identities follow directly by interchanging the row and column
indices in the definition.
\end{definition_box}

\begin{problem_box}
\textbf{Proposition:} If $A$ is
$m\times n$ and $B$ is $n\times p$, then
\[
{(AB)}^{\mathsf{T}}=B^{\mathsf{T}}A^{\mathsf{T}}.
\]
In particular, transposition reverses the order of multiplication.

\end{problem_box}

\begin{claim_box}
\textit{Proof:} Both sides are $p\times m$. For each $i,j$,
\[
\begin{aligned}
{\bigl({(AB)}^{\mathsf{T}}\bigr)}_{ij}
&={(AB)}_{ji}=\sum_{k=1}^n a_{jk}b_{ki}\\
&=\sum_{k=1}^n {(B^{\mathsf{T}})}_{ik}{(A^{\mathsf{T}})}_{kj}
={(B^{\mathsf{T}}A^{\mathsf{T}})}_{ij}.
\end{aligned}
\]
Equality of all entries proves the formula. \hfill$\square$
\end{claim_box}

\begin{problem_box}
\textbf{Theorem:} If
$A\in\mathbb{F}^{n\times n}$ is invertible, then $A^{\mathsf{T}}$
is invertible and
\[
\boxed{{(A^{\mathsf{T}})}^{-1}={(A^{-1})}^{\mathsf{T}}.}
\]
In prime notation, this is ${(A')}^{-1}={(A^{-1})}'$.

\end{problem_box}

\begin{claim_box}
\textit{Proof:} Transpose the identities $A^{-1}A=I_n$ and
$AA^{-1}=I_n$, using the product rule:
\[
A^{\mathsf{T}}{(A^{-1})}^{\mathsf{T}}=I_n,\qquad
{(A^{-1})}^{\mathsf{T}}A^{\mathsf{T}}=I_n.
\]
Thus ${(A^{-1})}^{\mathsf{T}}$ is a two-sided inverse of
$A^{\mathsf{T}}$. Uniqueness of the inverse proves the result.
\hfill$\square$
\end{claim_box}

\begin{claim_box}
\textbf{Definition:}
A matrix $A$ is in \emph{column echelon
form} if $A^{\mathsf{T}}$ is in row echelon form, and in
\emph{reduced column echelon form} (RCEF) if $A^{\mathsf{T}}$ is in RREF\@.
Explicitly, RCEF means that zero columns are on the right, the topmost
nonzero entry of each nonzero column is $1$, these leading $1$s move
strictly downward as the columns move right, and each is the only
nonzero entry in its row.
\end{claim_box}

\begin{example_box}
\textbf{Example:} Over $\mathbb{R}$,
\[
R=\begin{pmatrix}1&2&0&3\\0&0&1&-1\\0&0&0&0\end{pmatrix}
\quad\text{is in RREF},\qquad
R^{\mathsf{T}}=
\begin{pmatrix}1&0&0\\2&0&0\\0&1&0\\3&-1&0\end{pmatrix}
\quad\text{is in RCEF}.
\]
The leading entries of $R$ are in positions $(1,1)$ and $(2,3)$;
those of $R^{\mathsf{T}}$ are in positions $(1,1)$ and $(3,2)$.
RREF does not require every column to have a leading $1$, nor does
RCEF require every row to have one.
\end{example_box}

\begin{remark_box}
\textbf{Note:}
Elementary row operations produce RREF\@; the corresponding operations
on columns produce RCEF\@. Equivalently, reduce $A^{\mathsf{T}}$ by rows
and transpose the result to obtain the RCEF of $A$. These are different
operations: to solve $Ax=b$ in the original variables, use row
operations on the entire augmented matrix $[A\mid b]$. Column
operations on $A$ generally change the variables and require a
corresponding change of coordinates; they cannot simply replace row
operations on an augmented system.
\end{remark_box}

\begin{definition_box}
\textbf{Definition:} An \emph{elementary
matrix of order $m$} is a matrix obtained from $I_m$ by one elementary
row operation. There are three types:
\begin{enumerate}
    \item Interchange rows $i$ and $j$, where $i\neq j$.
    \item Multiply row $i$ by a nonzero scalar $c$.
    \item Add $c$ times row $j$ to row $i$, where $i\neq j$.
\end{enumerate}
In the second type, the matrix differs from $I_m$ only in its
$(i,i)$ entry, which is $c$. In the third type, it differs only in
its $(i,j)$ entry, which is $c$ instead of zero. The order of the
elementary matrix matches the number of rows of the matrix on which
the operation will act.
\end{definition_box}

\begin{example_box}
\textbf{Examples:} The elementary matrices
\[
P=\begin{pmatrix}0&1\\1&0\end{pmatrix},\qquad
D=\begin{pmatrix}c&0\\0&1\end{pmatrix}\ (c\neq0),\qquad
H=\begin{pmatrix}1&0\\c&1\end{pmatrix}
\]
are obtained from $I_2$ by $R_1\leftrightarrow R_2$,
$R_1\leftarrow cR_1$, and $R_2\leftarrow R_2+cR_1$, respectively.
For an arbitrary $2\times n$ matrix with rows $u$ and $v$,
\[
P\begin{pmatrix}u\\v\end{pmatrix}=\begin{pmatrix}v\\u\end{pmatrix},
\qquad
D\begin{pmatrix}u\\v\end{pmatrix}=\begin{pmatrix}cu\\v\end{pmatrix},
\qquad
H\begin{pmatrix}u\\v\end{pmatrix}=\begin{pmatrix}u\\v+cu\end{pmatrix}.
\]
\end{example_box}

\begin{problem_box}
\textbf{Theorem:}
Let $A\in\mathbb{F}^{m\times n}$, and let $E$ be obtained from $I_m$
by an elementary row operation. Then $EA$ is exactly the matrix
obtained by performing that same operation on $A$.

\end{problem_box}

\begin{claim_box}
\textit{Proof:} Write $A_{k,*}$ for row $k$ of $A$. Matrix
multiplication says that row $i$ of $EA$ is
\[
{(EA)}_{i,*}=\sum_{k=1}^m e_{ik}A_{k,*}.
\]
Thus each row of $E$ specifies a linear combination of the rows of $A$.
For an unchanged row of $I_m$, this formula returns the corresponding
row of $A$. It remains to check the changed rows:
\begin{enumerate}
    \item If rows $i,j$ of $I_m$ are swapped, row $i$ of $E$ selects
    $A_{j,*}$ and row $j$ selects $A_{i,*}$. Hence $EA$ swaps those rows.
    \item If row $i$ is multiplied by $c$, its only nonzero entry
    is $e_{ii}=c$. The formula gives $cA_{i,*}$.
    \item If $c$ times row $j$ is added to row $i$, then row $i$
    of $E$ has entries $e_{ii}=1$ and $e_{ij}=c$, with all others
    zero. The formula gives $A_{i,*}+cA_{j,*}$.
\end{enumerate}
These are precisely the three stated operations. \hfill$\square$
\end{claim_box}

\begin{problem_box}
\textbf{Corollary:}
The inverse of an elementary matrix is the elementary matrix for the
reverse operation: swap the same rows again, scale by $c^{-1}$,
or add $-c$ times the other row, respectively.

\end{problem_box}

\begin{claim_box}
\textit{Proof:} Let $E$ represent an operation and $F$ its reverse.
By the theorem, $FE$ applies the reverse operation to $E$, giving
$I_m$. Similarly, $EF=I_m$. Thus $F=E^{-1}$. \hfill$\square$
\end{claim_box}

\begin{claim_box}
\textbf{Consequence:} If the
successive operations are represented by $E_1,\ldots,E_k$, then the
final matrix is
\[
E_k\cdots E_2E_1A.
\]
The first operation appears nearest to $A$, on the right of the other
elementary factors. This follows by applying the theorem at each step
and using associativity. The product $Q=E_k\cdots E_1$ is invertible,
with inverse $Q^{-1}=E_1^{-1}\cdots E_k^{-1}$, as direct cancellation
in both products verifies.

For an augmented matrix, the same operation acts on every column:
\[
E[A\mid b]=[EA\mid Eb].
\]
Since $E$ is invertible, $Ax=b$ if and only if $EAx=Eb$.
This gives a matrix explanation of why elementary row operations on
the entire augmented matrix preserve the solution set.
\end{claim_box}

\begin{definition_box}
\textbf{Definition (equivalent matrices):} Two matrices
$A,B\in\mathbb{F}^{m\times n}$ are \emph{equivalent} if there are
invertible matrices $P\in\mathbb{F}^{m\times m}$ and
$Q\in\mathbb{F}^{n\times n}$ such that
\[
B=PAQ.
\]
Left multiplication allows invertible row operations, while right
multiplication allows invertible column operations. The matrices must
have the same dimensions.
\end{definition_box}

\begin{definition_box}
\textbf{Definition (similar matrices):} Two square matrices
$A,B\in\mathbb{F}^{n\times n}$ are \emph{similar} if there is an
invertible matrix $S\in\mathbb{F}^{n\times n}$ such that
\[
B=S^{-1}AS.
\]
Similar matrices represent the same linear operator in different
ordered bases. Similarity is a special case of equivalence, by taking
$P=S^{-1}$ and $Q=S$.
\end{definition_box}

\begin{example_box}
\textbf{Example (equivalent matrices):} Let
\[
A=\begin{pmatrix}1&0&0\\0&0&0\end{pmatrix},\qquad
P=\begin{pmatrix}0&1\\1&0\end{pmatrix},\qquad
Q=\begin{pmatrix}0&0&1\\0&1&0\\1&0&0\end{pmatrix}.
\]
Both $P$ and $Q$ are invertible, and
\[
PAQ=\begin{pmatrix}0&0&0\\0&0&1\end{pmatrix}=B.
\]
Thus $A$ and $B$ are equivalent. They are rectangular, so similarity
(which is defined only for square matrices) does not apply. The linear maps represented by A and B are also said to be equivalent.
\end{example_box}

\begin{example_box}
\textbf{Example (similar matrices):} Take
\[
A=\begin{pmatrix}1&0\\0&2\end{pmatrix},\qquad
S=\begin{pmatrix}1&-1\\0&1\end{pmatrix},\qquad
S^{-1}=\begin{pmatrix}1&1\\0&1\end{pmatrix}.
\]
Then
\[
S^{-1}AS
=\begin{pmatrix}1&1\\0&2\end{pmatrix}=B,
\]
so $A$ and $B$ are similar. The matrices are different, but represent
the same linear operator using different bases.
\end{example_box}


\subsection{Rotation matrices}

The two-dimensional rotation matrix is
\[
R_{\theta}=
\begin{pmatrix}
\cos\theta&-\sin\theta\\
\sin\theta&\cos\theta
\end{pmatrix}.
\]
The subscript $\theta$ indicates the angle of rotation. We regard
vectors as column matrices, so matrix multiplication acts on the left.
For a point at distance $r\geq0$ from the origin and direction angle
$\alpha$, its position vector is
\[
v=\begin{pmatrix}
r\cos(\alpha)\\
r\sin(\alpha)
\end{pmatrix}.
\]
Here $r$ is the radius (or length), and $\alpha$ is measured
counterclockwise from the positive horizontal axis. Multiplying gives
\[
R_{\theta}v=
\begin{pmatrix}
r\cos(\alpha)\cos(\theta)-r\sin(\alpha)\sin(\theta)\\
r\sin(\alpha)\cos(\theta)+r\cos(\alpha)\sin(\theta)
\end{pmatrix}
=
\begin{pmatrix}
r\cos(\alpha+\theta)\\
r\sin(\alpha+\theta)
\end{pmatrix}.
\]
Thus $R_{\theta}$ rotates the vector counterclockwise through angle
$\theta$ about the origin: it keeps the radius $r$ fixed and adds
$\theta$ to the direction angle $\alpha$. A negative $\theta$ gives a
clockwise rotation. For example, $\pi/2$ radians is $90^\circ$ and
$\pi$ radians is $180^\circ$:
\[
R_{\pi/2}=\begin{pmatrix}0&-1\\1&0\end{pmatrix},
\qquad
R_{\pi}=\begin{pmatrix}-1&0\\0&-1\end{pmatrix}.
\]
They send $v$ respectively to
\[
R_{\pi/2}v=
\begin{pmatrix}r\cos(\alpha+\pi/2)\\r\sin(\alpha+\pi/2)\end{pmatrix},
\qquad
R_{\pi}v=
\begin{pmatrix}r\cos(\alpha+\pi)\\r\sin(\alpha+\pi)\end{pmatrix}.
\]

\begin{problem_box}
\textbf{Proposition:}
For $v=\begin{pmatrix}r\cos(\alpha)\\r\sin(\alpha)\end{pmatrix}$,
\[
R_{\theta}v=
\begin{pmatrix}r\cos(\alpha+\theta)\\r\sin(\alpha+\theta)\end{pmatrix}.
\]
Thus $R_{\theta}$ keeps the radius fixed and adds $\theta$ to the
direction angle. In particular, the displayed matrices for
$R_{\pi/2}$ and $R_{\pi}$ have the stated actions on $v$.
\end{problem_box}

\begin{claim_box}
\textit{Proof:} By matrix multiplication, the entries of $R_{\theta}v$
are the two expressions on the left side of the displayed equality.
The angle-addition identities
\[
\cos(\alpha+\theta)=\cos\alpha\cos\theta-\sin\alpha\sin\theta,
\qquad
\sin(\alpha+\theta)=\sin\alpha\cos\theta+\cos\alpha\sin\theta
\]
therefore give the expression on its right. The radius remains $r$,
and the direction angle becomes $\alpha+\theta$, proving the rotation
claim. Substituting $(\cos(\pi/2),\sin(\pi/2))=(0,1)$ and
$(\cos\pi,\sin\pi)=(-1,0)$ into the matrix definition gives the two
special matrices. The same angle-addition identities give their stated
actions on $v$. \hfill$\square$
\end{claim_box}

For column vectors $u,v\in\mathbb{R}^2$, $\|u\|_2$ denotes the
Euclidean length of $u$, and $\langle u,v\rangle$ denotes their usual
dot product. A matrix is called \emph{orthogonal} if its transpose
multiplied by itself is the identity matrix. The rotation matrix satisfies
\[
\|R_{\theta}u\|_2=\|u\|_2,
\qquad
\langle R_{\theta}u,R_{\theta}v\rangle=\langle u,v\rangle,
\qquad
R_{\theta}^{\mathsf{T}}R_{\theta}=I_2,
\qquad \det R_{\theta}=1.
\]
The notation $R_{\theta}^{-1}$ means the inverse matrix, which undoes
the rotation; hence
\[
R_{\theta}^{-1}=R_{-\theta}=R_{\theta}^{\mathsf{T}}.
\]

Successive rotations about the origin satisfy
\[
R_{\alpha}R_{\beta}=R_{\alpha+\beta}.
\]

\begin{problem_box}
\textbf{Proposition:}
For every $\theta\in\mathbb{R}$, the matrix $R_{\theta}$ preserves
Euclidean lengths and dot products, is orthogonal, and has determinant
$1$. Its inverse is $R_{-\theta}=R_{\theta}^{\mathsf{T}}$. Moreover,
rotations about the origin compose according to
$R_{\alpha}R_{\beta}=R_{\alpha+\beta}$.
\end{problem_box}

\begin{claim_box}
\textit{Proof:} Write $c=\cos\theta$ and $s=\sin\theta$. Direct
multiplication gives
\[
R_{\theta}^{\mathsf{T}}R_{\theta}
=\begin{pmatrix}c&s\\-s&c\end{pmatrix}
  \begin{pmatrix}c&-s\\s&c\end{pmatrix}
=\begin{pmatrix}c^2+s^2&0\\0&c^2+s^2\end{pmatrix}
=I_2,
\]
using $c^2+s^2=1$. Also,
$\det R_{\theta}=c^2+s^2=1$. The orthogonality identity implies
\[
\|R_{\theta}u\|_2^2
=u^{\mathsf{T}}R_{\theta}^{\mathsf{T}}R_{\theta}u
=u^{\mathsf{T}}u=\|u\|_2^2,
\]
and
\[
\langle R_{\theta}u,R_{\theta}v\rangle
=u^{\mathsf{T}}R_{\theta}^{\mathsf{T}}R_{\theta}v
=u^{\mathsf{T}}v=\langle u,v\rangle.
\]
Since $R_{-\theta}=R_{\theta}^{\mathsf{T}}$, this transpose is the
inverse and undoes the rotation. Finally, multiplying $R_{\alpha}$ by
$R_{\beta}$ and applying the angle-addition identities to the entries
gives $R_{\alpha+\beta}$. The rightmost factor acts first, so the
formula also agrees with adding the two rotation angles.
\hfill$\square$
\end{claim_box}

\begin{problem_box}
\textbf{Proposition:}
A rotation through angle $\theta$ about a center other than the origin
is given by
\[
x\longmapsto c+R_{\theta}(x-c),
\qquad c=\begin{pmatrix}c_1\\c_2\end{pmatrix}.
\]
If $c\neq0$ and $\theta$ is not an integer multiple of $2\pi$, this
map is not linear and so cannot be represented by multiplication by a
single $2\times2$ matrix alone.
\end{problem_box}

\begin{claim_box}
\textit{Proof:} The displacement from the center to $x$ is $x-c$.
Rotating this displacement gives $R_{\theta}(x-c)$; adding $c$ places
the rotated point back relative to the original center. If $F$ denotes
this map, then
\[
F(x)=R_{\theta}x+(I_2-R_{\theta})c,
\qquad F(0)=(I_2-R_{\theta})c.
\]
For a linear map, $F(0)$ must be zero. A direct calculation gives
\[
\det(I_2-R_{\theta})
=(1-\cos\theta)^2+\sin^2\theta
=2(1-\cos\theta).
\]
This is nonzero when $\theta$ is not an integer multiple of $2\pi$,
so $I_2-R_{\theta}$ is invertible. Since $c\neq0$, it follows that
$F(0)\neq0$, proving that $F$ is not linear. When $\theta$ is an
integer multiple of $2\pi$, the map is the identity and is linear.
\hfill$\square$
\end{claim_box}

\begin{problem_box}
\textbf{Proposition:}
In three dimensions, $R_x(\theta)$, $R_y(\theta)$, and $R_z(\theta)$
are rotations through angle $\theta$ about their respective positive
coordinate axes, with matrices
\[
R_x(\theta)=
\begin{pmatrix}
1&0&0\\
0&\cos\theta&-\sin\theta\\
0&\sin\theta&\cos\theta
\end{pmatrix},\qquad
R_y(\theta)=
\begin{pmatrix}
\cos\theta&0&\sin\theta\\
0&1&0\\
-\sin\theta&0&\cos\theta
\end{pmatrix},
\]
\[
R_z(\theta)=
\begin{pmatrix}
\cos\theta&-\sin\theta&0\\
\sin\theta&\cos\theta&0\\
0&0&1
\end{pmatrix}.
\]
Each matrix fixes vectors along its named axis and rotates the
perpendicular coordinate plane. Rotations about different axes in
three dimensions need not commute: changing their order can change the
result.
\end{problem_box}

\begin{claim_box}
\textit{Proof:} Applying the matrices to a general column gives
\[
R_x(\theta)\begin{pmatrix}a\\b\\d\end{pmatrix}
=\begin{pmatrix}a\\b\cos\theta-d\sin\theta\\b\sin\theta+d\cos\theta\end{pmatrix},
\quad
R_y(\theta)\begin{pmatrix}a\\b\\d\end{pmatrix}
=\begin{pmatrix}a\cos\theta+d\sin\theta\\b\\-a\sin\theta+d\cos\theta\end{pmatrix},
\]
\[
R_z(\theta)\begin{pmatrix}a\\b\\d\end{pmatrix}
=\begin{pmatrix}a\cos\theta-b\sin\theta\\a\sin\theta+b\cos\theta\\d\end{pmatrix}.
\]
These formulas show that each transformation leaves its named-axis
coordinate fixed and acts on the perpendicular pair of coordinates by
a planar rotation. To verify that different-axis rotations need not
commute, take $\theta=\pi/2$:
\[
R_x(\pi/2)R_z(\pi/2)=
\begin{pmatrix}0&-1&0\\0&0&-1\\1&0&0\end{pmatrix},
\qquad
R_z(\pi/2)R_x(\pi/2)=
\begin{pmatrix}0&0&1\\1&0&0\\0&1&0\end{pmatrix}.
\]
The products differ, so these rotations do not commute. \hfill$\square$
\end{claim_box}

\subsection{Permutation matrices}

Let $\sigma$ be a permutation of $\{1,\ldots,n\}$, and let
$e_1,\ldots,e_n$ be the standard basis vectors of $\mathbb{F}^n$.
The associated \emph{permutation matrix} $P_{\sigma}$ is defined by
\[
P_{\sigma}e_j=e_{\sigma(j)},\qquad 1\leq j\leq n.
\]
\begin{problem_box}
\textbf{Proposition:}
Column $j$ of $P_{\sigma}$ is $e_{\sigma(j)}$, so every row and every
column has exactly one entry equal to $1$, and all other entries are
zero. Its action on a vector is
\[
(P_{\sigma}x)_i=x_{\sigma^{-1}(i)}.
\]
Moreover,
\[
P_{\sigma}^{-1}=P_{\sigma^{-1}}=P_{\sigma}^{\mathsf{T}},
\qquad P_{\sigma}^{\mathsf{T}}P_{\sigma}=I_n.
\]
Over $\mathbb{R}$ (and over $\mathbb{C}$ with its standard inner
product), permutation matrices preserve lengths and inner products.
Their determinants are
\[
\det(P_{\sigma})=\operatorname{sgn}(\sigma),
\]
where the sign $\operatorname{sgn}(\sigma)$ is viewed as an element
of $\mathbb{F}$.
\end{problem_box}

\begin{claim_box}
\textit{Proof:} Write $x=\sum_{j=1}^n x_j e_j$. By linearity and the
definition of $P_{\sigma}$,
\[
P_{\sigma}x=\sum_{j=1}^n x_j e_{\sigma(j)}.
\]
The coefficient in position $i$ is therefore $x_{\sigma^{-1}(i)}$.
Because $\sigma$ is a bijection, its values $\sigma(1),\ldots,\sigma(n)$
include every index exactly once. Hence every column and every row has
exactly one $1$.

The columns $e_{\sigma(j)}$ are distinct standard basis vectors, so they
are orthonormal. Thus $P_{\sigma}^{\mathsf{T}}P_{\sigma}=I_n$ and
$P_{\sigma}^{\mathsf{T}}=P_{\sigma}^{-1}$. This proves that the matrix
undoes its coordinate reordering and preserves inner products. It also
preserves lengths by taking the inner product of a vector with itself.

For the determinant, the Leibniz formula can be written as
\[
\det(P_{\sigma})=
\sum_{\pi\in S_n}\operatorname{sgn}(\pi)
\prod_{j=1}^n(P_{\sigma})_{\pi(j),j}.
\]
The product is nonzero only when $\pi(j)=\sigma(j)$ for every $j$;
there is exactly one such term, namely $\pi=\sigma$. Therefore
$\det(P_{\sigma})=\operatorname{sgn}(\sigma)$. \hfill$\square$
\end{claim_box}

\begin{problem_box}
\textbf{Proposition:}
For an $m\times n$ matrix $A$, let $\rho$ permute its $m$ row indices
and $\sigma$ permute its $n$ column indices. Left multiplication by
$P_{\rho}$ reorders rows, with row $i$ of $P_{\rho}A$ equal to row
$\rho^{-1}(i)$ of $A$. Right multiplication by $P_{\sigma}$ reorders
columns, with column $j$ of $AP_{\sigma}$ equal to column $\sigma(j)$
of $A$. Products satisfy
\[
P_{\sigma}P_{\tau}=P_{\sigma\circ\tau}.
\]
\end{problem_box}

\begin{claim_box}
\textit{Proof:} In row $i$ of $P_{\rho}A$, the only nonzero coefficient
selects row $\rho^{-1}(i)$ of $A$. For right multiplication,
\[
(AP_{\sigma})e_j=A(P_{\sigma}e_j)=Ae_{\sigma(j)},
\]
which is column $\sigma(j)$ of $A$. Finally, applying a product to a
standard basis vector gives
\[
P_{\sigma}P_{\tau}e_j
=P_{\sigma}e_{\tau(j)}
=e_{\sigma(\tau(j))}
=P_{\sigma\circ\tau}e_j.
\]
The two matrices agree on every standard basis vector, and therefore
are equal. \hfill$\square$
\end{claim_box}

\begin{example_box}
\textbf{Example:} For the cycle $\sigma=(1\ 2\ 3)$, so that
$\sigma(1)=2$, $\sigma(2)=3$, and $\sigma(3)=1$, we have
\[
P_{\sigma}=\begin{pmatrix}0&0&1\\1&0&0\\0&1&0\end{pmatrix},
\qquad
P_{\sigma}\begin{pmatrix}x_1\\x_2\\x_3\end{pmatrix}
=\begin{pmatrix}x_3\\x_1\\x_2\end{pmatrix}.
\]
For the transposition exchanging $1$ and $2$, the matrix is
\[
P_{(1\ 2)}=\begin{pmatrix}0&1\\1&0\end{pmatrix};
\]
left multiplication by it swaps the first two rows, as in an elementary
row operation.
\end{example_box}
\subsection{Change of basis}

\begin{definition_box}
\textbf{Definition:} Let $V$ be a
finite-dimensional vector space over $\mathbb{F}$ with ordered basis
$\mathcal{B}=(v_1,\ldots,v_n)$. Each $v\in V$ has a unique expansion
$v=x_1 v_1+\cdots+x_n v_n$. The \emph{matrix of $v$ relative to
$\mathcal{B}$}, also called its \emph{coordinate column}, is
\[
{[v]}_{\mathcal{B}}=
\begin{pmatrix}x_1\\\vdots\\x_n\end{pmatrix}
\in\mathbb{F}^{n\times1}.
\]
The order of the entries matches the order of the basis vectors.
We identify this column matrix with a vector of $\mathbb{F}^n$.
It records the coordinates of $v$, not necessarily the vector $v$
itself; changing the basis can change the column.
\end{definition_box}

\begin{example_box}
\textbf{Example:} In
$\mathbb{R}^2$, let $v=(3,1)$, let
$\mathcal{E}=((1,0),(0,1))$ be the standard basis, and choose
$\mathcal{B}=((1,1),(1,-1))$. Since
\[
(3,1)=2(1,1)+(1,-1),
\]
we have ${[v]}_{\mathcal{B}}=\begin{pmatrix}2\\1\end{pmatrix}$,
whereas ${[v]}_{\mathcal{E}}=\begin{pmatrix}3\\1\end{pmatrix}$.
The change-of-coordinates matrix from $\mathcal{E}$ to $\mathcal{B}$ is
\[
P_{\mathcal{B}\leftarrow\mathcal{E}}
=\begin{pmatrix}\frac12&\frac12\\[2pt]\frac12&-\frac12\end{pmatrix},
\qquad
{[v]}_{\mathcal{B}}
=P_{\mathcal{B}\leftarrow\mathcal{E}}{[v]}_{\mathcal{E}}
=\begin{pmatrix}2\\1\end{pmatrix}.
\]
\end{example_box}

\begin{problem_box}
\textbf{Proposition:} Let
$\mathcal{B}=(v_1,\ldots,v_n)$ and $\mathcal{C}=(w_1,\ldots,w_n)$ be
ordered bases of the same finite-dimensional vector space $V$. Define
\[
P_{\mathcal{C}\leftarrow\mathcal{B}}
={[I_V]}_{\mathcal{C}\leftarrow\mathcal{B}}
=\begin{pmatrix}
|& &|\\
{[v_1]}_{\mathcal{C}}&\cdots&{[v_n]}_{\mathcal{C}}\\
|& &|
\end{pmatrix}.
\]
Then for every $v\in V$,
\[
{[v]}_{\mathcal{C}}
=P_{\mathcal{C}\leftarrow\mathcal{B}}{[v]}_{\mathcal{B}},
\qquad
P_{\mathcal{B}\leftarrow\mathcal{C}}
=P_{\mathcal{C}\leftarrow\mathcal{B}}^{-1}.
\]
Thus the matrix converts coordinate columns in basis $\mathcal{B}$ into
coordinate columns in basis $\mathcal{C}$; it changes the coordinates,
not the vector itself.
\end{problem_box}

\begin{claim_box}
\textit{Proof:} Write $v=\sum_{j=1}^n x_jv_j$, so
${[v]}_{\mathcal{B}}=(x_1,\ldots,x_n)^{\mathsf T}$. By linearity of
coordinates,
\[
{[v]}_{\mathcal{C}}
=\sum_{j=1}^n x_j{[v_j]}_{\mathcal{C}}
=P_{\mathcal{C}\leftarrow\mathcal{B}}{[v]}_{\mathcal{B}}.
\]
Applying this formula in both directions shows that the two change-of-
coordinates matrices compose to the identity, so they are inverses.
\hfill$\square$
\end{claim_box}

\begin{definition_box}
\textbf{Definition:} Let
$T:V\longrightarrow W$ be linear, where $V,W$ are finite-dimensional
vector spaces over $\mathbb{F}$. Choose ordered bases
\[
\mathcal{B}=(v_1,\ldots,v_n)\quad\text{of }V,
\qquad \mathcal{C}=(w_1,\ldots,w_m)\quad\text{of }W.
\]
For each $j$, express $T(v_j)$ uniquely in the basis $\mathcal{C}$:
\[
T(v_j)=\sum_{i=1}^m a_{ij}w_i.
\]
The \emph{matrix of $T$ relative to $\mathcal{B}$ and $\mathcal{C}$} is
the $m\times n$ matrix
\[
{[T]}_{\mathcal{C}\leftarrow\mathcal{B}}=(a_{ij}).
\]
Its $j$th column is the coordinate column of $T(v_j)$ in the target
basis $\mathcal{C}$. Thus the number of columns is $\dim V$, and the
number of rows is $\dim W$. The arrow in the notation records that
input coordinates use $\mathcal{B}$ and output coordinates use
$\mathcal{C}$.
\end{definition_box}

\begin{problem_box}
\textbf{Proposition:}
With the ordered bases above fixed, put
$A={[T]}_{\mathcal{C}\leftarrow\mathcal{B}}$. For every $v\in V$,
\[
{[T(v)]}_{\mathcal{C}}
=A{[v]}_{\mathcal{B}}.
\]
Thus applying $T$ becomes multiplication by $A$ when inputs and outputs
are expressed in their respective bases. Abbreviating these coordinate
matrices by $M(T(v))$, $M(T)$, and $M(v)$ gives
\[
M(T(v))=M(T)M(v).
\]
\end{problem_box}

\begin{claim_box}
\textit{Proof:} Write $v=\sum_{j=1}^n x_j v_j$ and $A=(a_{ij})$.
By linearity and the definition of $A$,
\[
T(v)=\sum_{j=1}^n x_j T(v_j)
=\sum_{i=1}^m\left(\sum_{j=1}^n a_{ij}x_j\right)w_i,
\]
so uniqueness of coordinates gives
\[
{[T(v)]}_{\mathcal{C}}
=\begin{pmatrix}
\sum_{j=1}^n a_{1j}x_j\\
\vdots\\
\sum_{j=1}^n a_{mj}x_j
\end{pmatrix}
=A\begin{pmatrix}x_1\\\vdots\\x_n\end{pmatrix}
=A{[v]}_{\mathcal{B}}.
\]
Here an $m\times n$ matrix multiplies an $n\times1$ column to give
an $m\times1$ column. \hfill$\square$
\end{claim_box}

\begin{problem_box}
\textbf{Proposition:}
Fix ordered bases $\mathcal{B}=(v_1,\ldots,v_n)$ of $V$ and
$\mathcal{C}=(w_1,\ldots,w_m)$ of $W$. For
$T\in\mathcal{L}(V,W)$, abbreviate
\[
M(T)={[T]}_{\mathcal{C}\leftarrow\mathcal{B}}.
\]
Then for all $S,T\in\mathcal{L}(V,W)$ and $c\in\mathbb{F}$,
\[
M(S+T)=M(S)+M(T),\qquad M(cT)=cM(T).
\]

\end{problem_box}

\begin{claim_box}
\textit{Proof:} Write $M(S)=(s_{ij})$ and $M(T)=(t_{ij})$. By the
definition of a representing matrix, for each $j$,
\[
S(v_j)=\sum_{i=1}^m s_{ij}w_i,
\qquad T(v_j)=\sum_{i=1}^m t_{ij}w_i.
\]
Using pointwise addition of maps gives
\[
(S+T)(v_j)=S(v_j)+T(v_j)
=\sum_{i=1}^m(s_{ij}+t_{ij})w_i.
\]
Uniqueness of coordinates in $\mathcal{C}$ shows that the entry in row
$i$, column $j$ of $M(S+T)$ is $s_{ij}+t_{ij}$. These are exactly the
entries of $M(S)+M(T)$, proving the first identity.

Similarly, pointwise scalar multiplication gives
\[
(cT)(v_j)=cT(v_j)=\sum_{i=1}^m(ct_{ij})w_i.
\]
Thus the entry in row $i$, column $j$ of $M(cT)$ is $ct_{ij}$, which
is the corresponding entry of $cM(T)$. This proves the second identity.
\hfill$\square$
\end{claim_box}

\begin{claim_box}
\textbf{Consequence:} With these bases fixed, the assignment
\[
M:\mathcal{L}(V,W)\longrightarrow M_{m\times n}(\mathbb{F}),
\qquad T\longmapsto M(T),
\]
is itself a linear map. In particular,
$M(aS+bT)=aM(S)+bM(T)$ for all $a,b\in\mathbb{F}$.
All matrices in these identities must use the same ordered domain and
codomain bases.
\end{claim_box}

\begin{problem_box}
\textbf{Theorem:}
Let
$T:U\longrightarrow V$ and $S:V\longrightarrow W$ be linear maps
between finite-dimensional vector spaces over $\mathbb{F}$. Fix ordered
bases
\[
\mathcal{B}=(u_1,\ldots,u_p),\qquad
\mathcal{C}=(v_1,\ldots,v_n),\qquad
\mathcal{D}=(w_1,\ldots,w_m)
\]
of $U,V,W$, respectively. Then
\[
{[ST]}_{\mathcal{D}\leftarrow\mathcal{B}}
={[S]}_{\mathcal{D}\leftarrow\mathcal{C}}
{[T]}_{\mathcal{C}\leftarrow\mathcal{B}}.
\]
Abbreviating these matrices by $M(ST),M(S),M(T)$, this says
\[
M(ST)=M(S)M(T).
\]
Here $M(S)$ is $m\times n$, $M(T)$ is $n\times p$, and their product
is $m\times p$, as required for the matrix of $ST:U\longrightarrow W$.
\end{problem_box}

\begin{claim_box}
\textit{Proof:} Write $M(S)=(a_{ik})$ and $M(T)=(b_{kj})$.
By the definition of these matrices,
\[
T(u_j)=\sum_{k=1}^n b_{kj}v_k,
\qquad S(v_k)=\sum_{i=1}^m a_{ik}w_i.
\]
For each domain basis vector $u_j$, linearity of $S$ gives
\[
\begin{aligned}
(ST)(u_j)
&=S\left(\sum_{k=1}^n b_{kj}v_k\right)\\
&=\sum_{k=1}^n b_{kj}S(v_k)\\
&=\sum_{k=1}^n\sum_{i=1}^m b_{kj}a_{ik}w_i\\
&=\sum_{i=1}^m\left(\sum_{k=1}^n a_{ik}b_{kj}\right)w_i.
\end{aligned}
\]
Thus the entry in row $i$, column $j$ of $M(ST)$ is
$\sum_{k=1}^n a_{ik}b_{kj}$. By the definition of matrix multiplication,
this is precisely the corresponding entry of $M(S)M(T)$.
Since all corresponding entries agree, $M(ST)=M(S)M(T)$.
\hfill$\square$
\end{claim_box}

\begin{remark_box}
\textbf{Note:} The same intermediate basis $\mathcal{C}$ must be
used for the output of $T$ and the input of $S$. The order also matches
composition: $T$ acts first, so its matrix appears on the right.
\end{remark_box}

\begin{problem_box}
\textbf{Proposition:} Let
$T:V\longrightarrow W$ be linear, and choose ordered bases
$\mathcal{B},\mathcal{B}'$ of $V$ and $\mathcal{C},\mathcal{C}'$ of $W$.
Then
\[
{[T]}_{\mathcal{C}'\leftarrow\mathcal{B}'}
=P_{\mathcal{C}'\leftarrow\mathcal{C}}
{[T]}_{\mathcal{C}\leftarrow\mathcal{B}}
P_{\mathcal{B}\leftarrow\mathcal{B}'}.
\]
For an operator $T:V\longrightarrow V$, if $A$ and $A'$ are its
matrices in the old basis $\mathcal{B}$ and new basis $\mathcal{B}'$,
respectively, and $P=P_{\mathcal{B}\leftarrow\mathcal{B}'}$, this becomes
\[
A'=P^{-1}AP.
\]
Thus matrices of the same operator in different bases are similar.
\end{problem_box}

\begin{claim_box}
\textit{Proof:} For any $v\in V$, first convert its new input
coordinates to the old ones, apply the matrix of $T$ in the old bases,
and then convert the output coordinates to the new basis:
\[
{[T(v)]}_{\mathcal{C}'}
=P_{\mathcal{C}'\leftarrow\mathcal{C}}
{[T(v)]}_{\mathcal{C}}
=P_{\mathcal{C}'\leftarrow\mathcal{C}}
{[T]}_{\mathcal{C}\leftarrow\mathcal{B}}
{[v]}_{\mathcal{B}}
=P_{\mathcal{C}'\leftarrow\mathcal{C}}
{[T]}_{\mathcal{C}\leftarrow\mathcal{B}}
P_{\mathcal{B}\leftarrow\mathcal{B}'}
{[v]}_{\mathcal{B}'}.
\]
Since this holds for every $v$, the matrix identity follows. In the
operator case, $P_{\mathcal{B}'\leftarrow\mathcal{B}}=P^{-1}$, giving
$A'=P^{-1}AP$. \hfill$\square$
\end{claim_box}

\begin{example_box}
\textbf{Example:} Consider
$D:\mathcal{P}_2(\mathbb{R})\longrightarrow\mathcal{P}_1(\mathbb{R})$,
$D(p)=p'$, with ordered bases
$\mathcal{B}=(1,x,x^2)$ and $\mathcal{C}=(1,x)$. Since
\[
D(1)=0,\qquad D(x)=1,\qquad D(x^2)=2x,
\]
their coordinate columns give
\[
{[D]}_{\mathcal{C}\leftarrow\mathcal{B}}
=\begin{pmatrix}0&1&0\\0&0&2\end{pmatrix}.
\]
For example, the polynomial $p=a+bx+cx^2$ has coordinate column
${(a,b,c)}^{\mathsf{T}}$, and
\[
\begin{pmatrix}0&1&0\\0&0&2\end{pmatrix}
\begin{pmatrix}a\\b\\c\end{pmatrix}
=\begin{pmatrix}b\\2c\end{pmatrix},
\]
which is the coordinate column of $p'=b+2cx$ in $\mathcal{C}$.
\end{example_box}

\begin{remark_box}
\textbf{Note:} A linear map is not itself a
matrix: a matrix represents it after ordered bases have been chosen.
Changing either basis can change the matrix without changing the map.
For a map $T_A:\mathbb{F}^n\longrightarrow\mathbb{F}^m$, $T_A(x)=Ax$,
the representing matrix in the standard bases is exactly $A$.
\end{remark_box}

\begin{problem_box}
\textbf{Proposition:} Let $A\in\mathbb{F}^{m\times n}$ have rank $r$.
Then $A$ is equivalent to the $m\times n$ block matrix
\[
\begin{pmatrix}
I_r & 0\\
0 & 0
\end{pmatrix},
\]
where $0\leq r\leq\min(m,n)$.
\end{problem_box}

\begin{claim_box}
\textit{Proof:} Regard $A$ as the matrix of the linear map
$T_A:\mathbb{F}^n\to\mathbb{F}^m$. Choose a basis
$y_1,\ldots,y_r$ of $\operatorname{range}T_A$, and choose
$x_i\in\mathbb{F}^n$ with $T_A(x_i)=y_i$ for $1\leq i\leq r$.
Choose $x_{r+1},\ldots,x_n$ to be a basis of $\ker T_A$. Then
$x_1,\ldots,x_n$ form a basis of $\mathbb{F}^n$: applying $T_A$
to a linear relation first shows that the coefficients of
$x_1,\ldots,x_r$ vanish, and the remaining vectors are independent;
and any vector can be expressed in this basis by subtracting a suitable
linear combination of $x_1,\ldots,x_r$ to place it in the kernel.

Extend $y_1,\ldots,y_r$ to a basis $y_1,\ldots,y_m$ of
$\mathbb{F}^m$. Let $Q$ be the invertible matrix whose columns are the
coordinates of $x_1,\ldots,x_n$ in the standard basis, and let $P$ be
the inverse of the matrix whose columns are $y_1,\ldots,y_m$ in the
standard basis. The columns of $AQ$ are $y_1,\ldots,y_r$ followed by
zeros, and $P$ sends each $y_i$ to the $i$th standard basis vector.
Therefore $PAQ$ is the stated block matrix. Since $P$ and $Q$ are
invertible, $A$ is equivalent to it. Finally, the rank is at most both
$m$ and $n$, so $r\leq\min(m,n)$. \hfill$\square$
\end{claim_box}

\begin{problem_box}
\textbf{Proposition:} Let $V$ and $W$ be finite-dimensional vector
spaces over $\mathbb{F}$, and let $T:V\longrightarrow W$ be linear.
Writing $n=\dim V$, $m=\dim W$, and
$r=\dim(\operatorname{range}T)$, there are ordered bases of $V$ and
$W$ relative to which the matrix of $T$ is
\[
\begin{pmatrix}
I_r & 0\\
0 & 0
\end{pmatrix},
\qquad 0\leq r\leq\min(m,n).
\]
\end{problem_box}

\begin{claim_box}
\textit{Proof:} Choose a basis $y_1,\ldots,y_r$ of
$\operatorname{range}T$, and for each $i$ choose $x_i\in V$ such that
$T(x_i)=y_i$. By rank-nullity, $\ker T$ has dimension $n-r$; choose a
basis $x_{r+1},\ldots,x_n$ of $\ker T$. Then
$x_1,\ldots,x_n$ is a basis of $V$. Indeed, applying $T$ to a linear
relation among these vectors shows that the coefficients of
$x_1,\ldots,x_r$ vanish, after which independence of the kernel basis
forces the remaining coefficients to vanish. For spanning, given
$x\in V$, express $T(x)$ as a linear combination of
$y_1,\ldots,y_r$; subtracting the corresponding combination of
$x_1,\ldots,x_r$ from $x$ leaves a vector in $\ker T$.

Extend $y_1,\ldots,y_r$ to an ordered basis
$y_1,\ldots,y_m$ of $W$. Use
$\mathcal{B}=(x_1,\ldots,x_n)$ as the basis of $V$ and
$\mathcal{C}=(y_1,\ldots,y_m)$ as the basis of $W$. The first $r$
basis vectors of $V$ map to the first $r$ basis vectors of $W$, and
the remaining basis vectors map to zero. Hence the matrix of $T$ in
these bases has the stated form. Since $r$ is the dimension of a
subspace of both $V$ and $W$, $r\leq\min(m,n)$. \hfill$\square$
\end{claim_box}

\subsection{Homogeneous and nonhomogeneous linear systems}

\begin{definition_box}
\textbf{Definition:} A system of $m$ linear equations in $n$
unknowns over $\mathbb{F}$ can be written as
\[
Ax=b,
\qquad A\in\mathbb{F}^{m\times n},\quad
x\in\mathbb{F}^n,\quad b\in\mathbb{F}^m.
\]
Here $A=(a_{ij})$ is the given coefficient matrix, $b$ is the given
right-hand side, and $x$ is the unknown vector. The $i$th equation is
$\sum_{j=1}^n a_{ij}x_j=b_i$. The system is \emph{homogeneous} if
$b=0$, and \emph{nonhomogeneous} if $b\neq0$. A system is
\emph{consistent} if it has at least one solution, and
\emph{inconsistent} otherwise.
\end{definition_box}

\begin{claim_box}
\textbf{Definition:} The \emph{augmented matrix}
of $Ax=b$ records the coefficients and right-hand side together:
\[
[A\mid b]=\left(\begin{array}{cccc|c}
a_{11}&a_{12}&\cdots&a_{1n}&b_1\\
\vdots&\vdots&\ddots&\vdots&\vdots\\
a_{m1}&a_{m2}&\cdots&a_{mn}&b_m
\end{array}\right).
\]
Each row represents one equation; the vertical line separates the
coefficient columns from the right-hand side.
\end{claim_box}

\begin{claim_box}
\textbf{Elementary row operations:} The following operations on
the \emph{entire} augmented matrix preserve the solution set:
\begin{enumerate}
    \item Interchange two rows: $R_i\leftrightarrow R_j$.
    \item Multiply a row by a nonzero scalar:
    $R_i\leftarrow cR_i$, where $c\neq0$.
    \item Add a scalar multiple of another row to a row:
    $R_i\leftarrow R_i+cR_j$, where $i\neq j$.
\end{enumerate}
\end{claim_box}

\begin{problem_box}
\textbf{Theorem:} Suppose
$[A'\mid b']$ is obtained from $[A\mid b]$ by a finite sequence of
elementary row operations. Then
\[
\{x\in\mathbb{F}^n:Ax=b\}
=\{x\in\mathbb{F}^n:A'x=b'\}.
\]

\end{problem_box}

\begin{claim_box}
\textit{Proof:} It suffices to prove the claim for one operation.
Write $L_i(x)=\sum_{k=1}^n a_{ik}x_k$, so the original $i$th
equation is $L_i(x)=b_i$. Fix any candidate vector $x\in\mathbb{F}^n$.
\begin{enumerate}
    \item \emph{Interchanging rows:} Swapping rows $i$ and $j$ only
    changes the order of the equations. Thus $x$ satisfies all the
    new equations if and only if it satisfies all the original ones.
    \item \emph{Multiplying a row by a nonzero scalar:} The only
    changed equation is $cL_i(x)=cb_i$. Since $c\neq0$ has an inverse
    in $\mathbb{F}$,
    \[
    L_i(x)=b_i\quad\Longleftrightarrow\quad cL_i(x)=cb_i.
    \]
    All other equations are unchanged, so the solution set is unchanged.
    \item \emph{Adding a multiple of another row:} The operation
    $R_i\leftarrow R_i+cR_j$, with $i\neq j$, replaces equation $i$ by
    \[
    L_i(x)+cL_j(x)=b_i+cb_j
    \]
    and leaves equation $j$ and every other equation unchanged.
    If $x$ solves the original system, adding the two corresponding
    equalities shows that it satisfies the replacement equation.
    Conversely, if $x$ solves the new system, it satisfies both the
    replacement equation and $L_j(x)=b_j$. Subtracting $c$ times the
    latter gives
    \[
    L_i(x)=(b_i+cb_j)-cb_j=b_i.
    \]
    Hence $x$ also solves the original system.
\end{enumerate}
Each elementary operation therefore preserves exactly the same
solutions. Applying this conclusion at every step proves the result
for any finite sequence. \hfill$\square$
\end{claim_box}

\begin{remark_box}
\textbf{Caution:} Perform each operation on the entire augmented
row, including the right-hand side. The nonzero condition in row
scaling is essential: multiplying the equation $x=1$ by zero gives
$0=0$, which admits every $x$ and loses the original constraint.
\end{remark_box}

\begin{definition_box}
\textbf{Definition:}
In a row echelon matrix, the leading entry of each nonzero row is
called a \emph{pivot entry}. Its location is a \emph{pivot position},
and a column containing one is a \emph{pivot column}. A pivot can be
any nonzero scalar in REF\@; in RREF every pivot is $1$ and all other
entries in its column are zero. Not every nonzero entry is a pivot.
For an arbitrary coefficient matrix $A$, its pivot columns mean the
column indices containing pivots after row reduction.
\end{definition_box}

\begin{claim_box}
\textbf{Pivot variables and free variables:} For a consistent
system, a variable $x_j$ is a \emph{pivot variable} if column $j$ of
the coefficient matrix has a pivot after row reduction. Otherwise it
is a \emph{free variable}, to which an arbitrary parameter can be
assigned. The final column of $[A\mid b]$ is not a variable column:
a pivot there instead signals an equation $0=d$ with $d\neq0$, hence
inconsistency.

For example, the RREF augmented matrix
\[
\left(\begin{array}{ccc|c}1&0&2&5\\0&1&-1&-1\end{array}\right)
\]
has pivot entries $1$ in positions $(1,1)$ and $(2,2)$. Its pivot
variables are $x_1,x_2$, and $x_3$ is free. The entries $2$ and $-1$
in column $3$ are not pivots. The equations are
$x_1=5-2x_3$ and $x_2=-1+x_3$.
\end{claim_box}

\begin{claim_box}
\textbf{Gaussian elimination:}
\begin{enumerate}
    \item Form $[A\mid b]$. Starting with the leftmost coefficient
    column, find a nonzero entry in or below the current row. Swap
    rows to place it in the current row, then use it to eliminate all
    entries below it. If no such entry exists, skip that column.
    Continue with the next row and the columns to the right until
    the matrix is in row echelon form.
    \item If a row has the form
    $\left(\begin{array}{ccc|c}0&\cdots&0&d\end{array}\right)$
    with $d\neq0$, it represents $0=d$. The system is inconsistent.
    A completely zero row imposes no condition.
    \item Otherwise, the variables whose coefficient columns contain
    pivots are \emph{pivot variables}; the remaining variables are
    \emph{free variables}. Assign arbitrary parameters in $\mathbb{F}$
    to the free variables and solve for the pivot variables from the
    bottom equation upward (\emph{back-substitution}).
    \item Alternatively, divide each pivot row by its pivot and
    eliminate entries above the pivots as well. This gives reduced
    row echelon form (\emph{Gauss--Jordan elimination}), from which
    the pivot variables can be read directly in terms of the free ones.
\end{enumerate}
\end{claim_box}

\begin{problem_box}
\textbf{Theorem:}
Let
$A\in\mathbb{F}^{m\times n}$, $b\in\mathbb{F}^m$, and let
$[R\mid d]$ be the RREF of $[A\mid b]$. Suppose $R$ has $r$ pivot
columns. Then:
\begin{enumerate}
    \item $Ax=b$ is consistent if and only if the augmented column
    contains no pivot, equivalently there is no row $[0\ \cdots\ 0\mid1]$.
    \item If the system is consistent, every assignment of the $n-r$
    free variables extends to exactly one solution.
    \item A consistent system has a unique solution if and only if
    $r=n$, that is, every variable column contains a pivot.
    \item The homogeneous system $Ax=0$ has a nonzero solution if and
    only if $r<n$.
\end{enumerate}
\end{problem_box}

\begin{claim_box}
\textit{Proof:} Row operations preserve solutions, so we may work with
$Rx=d$. A pivot in the augmented column gives the equation $0=1$,
making a solution impossible.

Conversely, suppose there is no such pivot. Every nonzero row then
has a pivot in a coefficient column. Let the pivot columns be
$p_1,\ldots,p_r$ in row order, and let $F$ be the set of the remaining
column indices. Since each pivot column has a single $1$ and zeros
elsewhere, row $i$ gives
\[
x_{p_i}=d_i-\sum_{j\in F}R_{ij}x_j,
\qquad 1\leq i\leq r.
\]
All remaining rows say $0=0$. Assign arbitrary values $x_j=t_j$ for
$j\in F$; these formulas determine the pivot variables uniquely and
satisfy every equation. This proves consistency and the second claim.

If $r=n$, there are no free variables, so the formulas give just one
solution. If $r<n$, assigning all free variables zero or assigning
one of them $1$ gives distinct solutions. Thus a consistent system
is unique exactly when $r=n$.

For $Ax=0$, apply the same sequence of row operations used on $[A\mid b]$.
The coefficient matrix again becomes $R$, while the zero right-hand
side remains zero throughout.
Thus the homogeneous system is always consistent, and if $r=n$ its
only solution is zero. If $r<n$, setting one free variable to $1$
produces a nonzero solution. \hfill$\square$
\end{claim_box}

\begin{problem_box}
\textbf{Proposition:} The number $r$
of pivots in a row echelon form of $A$ equals $\operatorname{rank}A$.
Consequently, a consistent system has $n-\operatorname{rank}A$ free
variables.

\end{problem_box}

\begin{claim_box}
\textit{Proof:} Further reducing an echelon form to RREF does not
change the number of pivots. Let $R$ be the resulting coefficient
matrix. Elementary row operations on vectors in $\mathbb{F}^m$ are
invertible linear maps: their inverses are the reverse row operations.
Their composition is an isomorphism $E:\mathbb{F}^m\longrightarrow
\mathbb{F}^m$, and
\[
Rx=E(Ax)\qquad\text{for all }x\in\mathbb{F}^n.
\]
Hence $\operatorname{range}T_R=E(\operatorname{range}T_A)$, so $R$
and $A$ have the same rank. The $r$ pivot columns of $R$ are the first
$r$ standard coordinate columns of $\mathbb{F}^m$. They are linearly
independent and span all columns of $R$, because every entry below
row $r$ is zero. Thus $\operatorname{rank}R=r$, proving the claim.
\hfill$\square$
\end{claim_box}

\begin{problem_box}
\textbf{Corollary:} We have
\[
Ax=b\text{ is consistent}
\quad\Longleftrightarrow\quad
\operatorname{rank}[A\mid b]=\operatorname{rank}A.
\]
Moreover, every $b\in\mathbb{F}^m$ gives a consistent system if and
only if $A$ has a pivot in every row after row reduction.

\end{problem_box}

\begin{claim_box}
\textit{Proof:} In the reduced augmented matrix, the augmented column
adds a pivot exactly when the system is inconsistent. The rank formula
therefore follows from the pivot-count proposition. Finally, every
right-hand side is attainable exactly when $\operatorname{range}T_A$
is all of $\mathbb{F}^m$, equivalently $\operatorname{rank}A=m$.
By the proposition, this means there are $m$ pivots, one in every row.
\hfill$\square$
\end{claim_box}

\begin{remark_box}
\textbf{Note:} For a consistent system,
the free-variable assignments give a bijection between the solution
set and $\mathbb{F}^{n-r}$. Thus if $r<n$ there are infinitely many
solutions over an infinite field; over a finite field with $q$ elements
there are exactly $q^{n-r}$ solutions.
\end{remark_box}

\begin{example_box}
\textbf{Example:} Consider
\[ 
x+2y=3,\qquad -y+z=1.
\]
The augmented matrix reduces as follows:
\[
\begin{aligned}
\left(\begin{array}{ccc|c}1&2&0&3\\0&-1&1&1\end{array}\right)
&\xrightarrow{R_2\leftarrow -R_2}
\left(\begin{array}{ccc|c}1&2&0&3\\0&1&-1&-1\end{array}\right)\\[4pt]
&\xrightarrow{R_1\leftarrow R_1-2R_2}
\left(\begin{array}{ccc|c}1&0&2&5\\0&1&-1&-1\end{array}\right).
\end{aligned}
\]
The pivot variables are $x,y$, and $z$ is free. Setting $z=t$ gives
\[
x=5-2t,\qquad y=t-1,\qquad z=t,
\]
so all solutions are
$(x,y,z)=(5,-1,0)+t(-2,1,1)$ for $t\in\mathbb{R}$.
With $s=t-1$, this is equivalently
$(3,0,1)+s(-2,1,1)$, the form used below.
\end{example_box}

\begin{example_box}
\textbf{Example:} For
$x+y=1$ and $2x+2y=3$,
\[
\left(\begin{array}{cc|c}1&1&1\\2&2&3\end{array}\right)
\xrightarrow{R_2\leftarrow R_2-2R_1}
\left(\begin{array}{cc|c}1&1&1\\0&0&1\end{array}\right).
\]
The last row says $0=1$, so there is no solution.
\end{example_box}

\begin{claim_box}
\textbf{Homogeneous systems:} The system $Ax=0$ always has the
\emph{trivial solution} $x=0$. Its complete solution set is
\[
\{x\in\mathbb{F}^n:Ax=0\}=\ker T_A,
\qquad T_A(x)=Ax,
\]
so it is a subspace of $\mathbb{F}^n$. In particular, every linear
combination of homogeneous solutions is another homogeneous solution.
Writing $r=\operatorname{rank}A=\dim(\operatorname{range}T_A)$,
rank-nullity gives
\[
\dim(\ker T_A)=n-r.
\]
Thus a nontrivial solution exists exactly when $r<n$.
\end{claim_box}

\begin{claim_box}
\textbf{Nonhomogeneous systems and consistency:} The system
$Ax=b$ has a solution exactly when
\[
b\in\operatorname{range}T_A.
\]
Equivalently, $b$ must be a linear combination of the columns of $A$,
because $Ax=x_1 A_1+\cdots+x_n A_n$, where $A_j$ is the $j$th column.
Unlike a homogeneous system, a nonhomogeneous system need not be
consistent.
\end{claim_box}

\begin{problem_box}
\textbf{Proposition:}
Suppose $x_p$
is one particular solution of $Ax=b$. Then all solutions are given by
\[
x=x_p+h,\qquad h\in\ker T_A.
\]
In other words, the solution set is $x_p+\ker T_A$.
\end{problem_box}

\begin{claim_box}
\textit{Proof:} Since $Ax_p=b$, for any $x\in\mathbb{F}^n$,
\[
Ax=b\quad\Longleftrightarrow\quad A(x-x_p)=0
\quad\Longleftrightarrow\quad x-x_p\in\ker T_A.
\]
This proves both that every solution has the stated form and that every
vector of that form is a solution. \hfill$\square$
\end{claim_box}

\begin{claim_box}
\textbf{Interpretation:} A consistent solution set is a translate
of the null space, also called an \emph{affine subspace}. If $b\neq0$,
it is not a vector subspace: it does not contain $0$, since $A0=0\neq b$.
The difference of any two solutions lies in the null space, but the sum
of two solutions need not be a solution of the same system.
\end{claim_box}

\begin{claim_box}
\textbf{Consequences of rank-nullity:}
\begin{enumerate}
    \item A consistent system has a unique solution if and only if
    $\ker T_A=\{0\}$, equivalently $r=n$. Otherwise its solutions have
    $n-r$ free parameters: if $h_1,\ldots,h_{n-r}$ is a basis of the
    kernel, every solution is uniquely expressible as
    \[
    x=x_p+c_1 h_1+\cdots+c_{n-r}h_{n-r},\qquad c_i\in\mathbb{F}.
    \]
    Over an infinite field such as $\mathbb{R}$ or $\mathbb{C}$, a
    consistent system with $r<n$ therefore has infinitely many solutions.
    \item If $m<n$, then $r\leq m<n$, so the homogeneous system has
    nonzero solutions. A nonhomogeneous system may be inconsistent,
    but if it is consistent, its solution cannot be unique.
    \item Every right-hand side $b\in\mathbb{F}^m$ gives a consistent
    system exactly when $T_A$ is surjective, equivalently $r=m$.
    In particular, if $m>n$, some right-hand sides have no solution;
    this does not mean that every right-hand side is inconsistent.
\end{enumerate}
\end{claim_box}

\begin{example_box}
\textbf{Example:} Over $\mathbb{R}$,
consider the homogeneous system
\[
x+2y=0,\qquad -y+z=0.
\]
Its solutions are $(x,y,z)=t(-2,1,1)$, where $t\in\mathbb{R}$.
Now change the right-hand side:
\[
x+2y=3,\qquad -y+z=1.
\]
A particular solution is $(3,0,1)$. Thus all solutions are
\[
(x,y,z)=(3,0,1)+t(-2,1,1),\qquad t\in\mathbb{R}.
\]
The homogeneous solution set is a line through the origin; the
nonhomogeneous solution set is a parallel line that does not pass
through the origin. Both have one free parameter.
\end{example_box}

\begin{example_box}
\textbf{Examples:} Over
$\mathbb{R}$, the system
\[
x+y=1,\qquad 2x+2y=3
\]
is inconsistent: doubling the first equation would give $2x+2y=2$,
contradicting the second. In contrast,
\[
x+y=3,\qquad x-y=1
\]
has the unique solution $(x,y)=(2,1)$. Its associated homogeneous system
has only the zero solution.
\end{example_box}

\subsection{Codimension}

\begin{definition_box}
\textbf{Definition:} Let $U$ be a subspace of a vector space $V$ over
$\mathbb{F}$. The \emph{codimension} of $U$ in $V$ is
\[
\operatorname{codim}_V U=\dim(V/U).
\]
If $V$ is finite-dimensional, the quotient-space dimension theorem gives
\[
\operatorname{codim}_V U=\dim V-\dim U.
\]
Thus, in finite dimensions, codimension counts how many dimensions of
$V$ lie outside $U$. Equivalently, it is the number of additional basis
vectors needed to extend a basis of $U$ to a basis of $V$.
\end{definition_box}

\begin{example_box}
\textbf{Examples:} In $\mathbb{F}^3$, the coordinate plane
$U=\{(x,y,0):x,y\in\mathbb{F}\}$ has dimension $2$, so
$\operatorname{codim}_{\mathbb{F}^3}U=1$. The coordinate axis
$L=\{(x,0,0):x\in\mathbb{F}\}$ has dimension $1$, so
$\operatorname{codim}_{\mathbb{F}^3}L=2$.
\end{example_box}

\begin{problem_box}
\textbf{Proposition:} Let $A\in\mathbb{F}^{m\times n}$ and let
$T_A:\mathbb{F}^n\longrightarrow\mathbb{F}^m$ be the linear map
$T_A(x)=Ax$. Then
\[
\operatorname{codim}_{\mathbb{F}^n}(\ker T_A)=\operatorname{rank}A,
\qquad
\operatorname{codim}_{\mathbb{F}^m}(\operatorname{range}T_A)
=m-\operatorname{rank}A.
\]
In particular, $T_A$ is surjective exactly when its range has
codimension zero, and it is injective exactly when
$\operatorname{rank}A=n$.
\end{problem_box}

\begin{claim_box}
\textit{Proof:} By rank-nullity,
\[
\dim(\ker T_A)=n-\operatorname{rank}A,
\qquad
\dim(\operatorname{range}T_A)=\operatorname{rank}A.
\]
Subtracting these dimensions from the dimensions of their ambient spaces
$\mathbb{F}^n$ and $\mathbb{F}^m$ gives the two codimension formulas.
The surjectivity statement follows because a subspace of $\mathbb{F}^m$
has codimension zero exactly when it is all of $\mathbb{F}^m$. By the
kernel criterion for injectivity, $T_A$ is injective exactly when
$\dim(\ker T_A)=0$, equivalently $\operatorname{rank}A=n$.
\hfill$\square$
\end{claim_box}

\subsection{Rank factorization}

\begin{definition_box}
\textbf{Definition:} A \emph{rank factorization} of a matrix
$A\in\mathbb{F}^{m\times n}$ of rank $r$ is a factorization
\[
A=CR,
\]
where $C\in\mathbb{F}^{m\times r}$ has rank $r$ and
$R\in\mathbb{F}^{r\times n}$ has rank $r$. Thus, $C$ has linearly
independent columns and $R$ has linearly independent rows.
\end{definition_box}

\begin{problem_box}
\textbf{Proposition:} Let $A\in\mathbb{F}^{m\times n}$ and
$0\leq r\leq\min(m,n)$. Then $\operatorname{rank}A=r$ if and only if
there exist $C\in\mathbb{F}^{m\times r}$ and
$R\in\mathbb{F}^{r\times n}$ such that
\[
A=CR,\qquad \operatorname{rank}C=\operatorname{rank}R=r.
\]
\end{problem_box}

\begin{claim_box}
\textit{Proof:} ($\Rightarrow$) Suppose $\operatorname{rank}A=r$.
If $r=0$, then $A=0$ and the factorization with
$C\in\mathbb{F}^{m\times0}$ and $R\in\mathbb{F}^{0\times n}$ is
immediate. Suppose $r>0$. Choose $r$ columns of $A$ that form a basis
of its column space, and let $C$ be the matrix with these columns.
Each column of $A$ has a unique expression as a linear combination of
the columns of $C$. Put the corresponding coordinate vectors into the
columns of a matrix $R\in\mathbb{F}^{r\times n}$. By construction,
$A=CR$ and $C$ has rank $r$. The columns of $R$ corresponding to the
chosen columns of $C$ are the standard basis vectors of $\mathbb{F}^r$,
so $R$ has rank $r$ as well. Thus $A=CR$ is a rank factorization.

($\Leftarrow$) Conversely, suppose $A=CR$ with the stated dimensions
and ranks. The map represented by $R$ is surjective onto $\mathbb{F}^r$,
and the map represented by $C$ is injective on $\mathbb{F}^r$. Hence
\[
\operatorname{range}T_A
=C(\operatorname{range}T_R)
=C(\mathbb{F}^r)
=\operatorname{range}T_C,
\]
so $\operatorname{rank}A=\dim(\operatorname{range}T_C)=r$.
\hfill$\square$
\end{claim_box}

\begin{example_box}
\textbf{Example:} The matrix
\[
A=\begin{pmatrix}1&0&2\\0&1&3\\1&1&5\end{pmatrix}
=\begin{pmatrix}1&0\\0&1\\1&1\end{pmatrix}
\begin{pmatrix}1&0&2\\0&1&3\end{pmatrix}=CR
\]
has rank $2$. The two columns of $C$ are independent, and the two rows
of $R$ are independent, so this is a rank factorization.
\end{example_box}

\begin{problem_box}
\textbf{Proposition:} Let $A,B\in\mathbb{F}^{m\times n}$. Then
\[
\operatorname{rank}(A+B)\leq\operatorname{rank}A+\operatorname{rank}B,
\]
with equality if and only if
\[
\operatorname{Col}(A)\cap\operatorname{Col}(B)=\{0\},
\qquad
\operatorname{Row}(A)\cap\operatorname{Row}(B)=\{0\}.
\]
\end{problem_box}

\begin{claim_box}
\textit{Proof:} Write rank factorizations
$A=C_A R_A$ and $B=C_B R_B$, where
$r=\operatorname{rank}A$ and $s=\operatorname{rank}B$. Then
\[
A+B=
\begin{pmatrix}C_A&C_B\end{pmatrix}
\begin{pmatrix}R_A\\R_B\end{pmatrix}.
\]
The first factor has rank at most $r+s$, and so does the second.
Since the rank of a product is at most the rank of either factor,
\[
\operatorname{rank}(A+B)\leq r+s.
\]
Equality holds exactly when both factors have rank $r+s$: if they do,
 the second factor maps onto $\mathbb{F}^{r+s}$ and the first is injective,
 so their product has rank $r+s$. The columns
of $C_A$ and $C_B$ are bases for $\operatorname{Col}(A)$ and
$\operatorname{Col}(B)$, respectively, so the first factor has rank
$r+s$ exactly when these column spaces intersect only in $0$. The rows
of $R_A$ and $R_B$ are bases for $\operatorname{Row}(A)$ and
$\operatorname{Row}(B)$, respectively, so the second factor has rank
$r+s$ exactly when these row spaces intersect only in $0$. This proves
both directions. \hfill$\square$

\end{claim_box}

\subsection{Generalized inverses}

\begin{definition_box}
\textbf{Definition:} Let $A\in\mathbb{F}^{m\times n}$. A matrix
$G\in\mathbb{F}^{n\times m}$ is a \emph{generalized inverse}, or
\emph{$g$-inverse}, of $A$ if
\[
AGA=A.
\]
This convention is also called an inner inverse (or a $\{1\}$-inverse).
\end{definition_box}
\begin{problem_box}
\textbf{Proposition:} Let $A\in\mathbb{F}^{m\times n}$ and
$G\in\mathbb{F}^{n\times m}$. If
\[
(GA)^2=GA\qquad\text{and}\qquad\operatorname{rank}(GA)=\operatorname{rank}A,
\]
then $G$ is a generalized inverse of $A$.
\end{problem_box}

\begin{claim_box}
\textit{Proof:} Since $\ker A\subseteq\ker(GA)$ and both maps have
domain $\mathbb{F}^n$, the rank equality and rank-nullity theorem imply
$\ker A=\ker(GA)$. For every $x\in\mathbb{F}^n$, idempotence gives
\[
GA(I_n-GA)x=(GA-(GA)^2)x=0.
\]
Thus $(I_n-GA)x\in\ker(GA)=\ker A$, so $A(I_n-GA)x=0$ for every
$x$. Hence $A=AGA$, and $G$ is a generalized inverse of $A$.
\hfill$\square$
\end{claim_box}

\begin{problem_box}
\textbf{Theorem:} Every matrix over a field has a generalized inverse.
\end{problem_box}

\begin{claim_box}
\textit{Proof:} Let $r=\operatorname{rank}A$. Choose a basis
$y_1,\ldots,y_r$ of $\operatorname{range}A$, and choose vectors
$x_1,\ldots,x_r\in\mathbb{F}^n$ such that $Ax_i=y_i$. Extend
$y_1,\ldots,y_r$ to a basis $y_1,\ldots,y_m$ of $\mathbb{F}^m$.
Define a linear map $G:\mathbb{F}^m\to\mathbb{F}^n$ on this basis by
\[
G(y_i)=x_i\quad(1\leq i\leq r),
\qquad
G(y_i)=0\quad(r<i\leq m).
\]
For any $x\in\mathbb{F}^n$, write
$Ax=\sum_{i=1}^r c_i y_i$. Then
\[
AGA x=A\left(\sum_{i=1}^r c_i x_i\right)
=\sum_{i=1}^r c_i y_i=Ax.
\]
Thus the matrix of $G$ is a generalized inverse of $A$.
\hfill$\square$
\end{claim_box}
\begin{example_box}
\textbf{Case-wise method for finding a $g$-inverse:} Let
$r=\operatorname{rank}A$. For index sets $I$ and $J$, let $S_I$
extract rows $I$ and let $T_J$ insert coordinates into positions $J$
(with zeros elsewhere).
\begin{enumerate}
    \item If $r=0$, then $A=0$; $G$ can be any $n$x$m$ matrix.
    \item If $r=n$ (full column rank), choose $n$ independent rows $I$.
    Set $A_I=S_IA$ and $G=A_I^{-1}S_I$. Then $GA=I_n$.
    \item If $r=m<n$ (full row rank), choose $m$ independent columns $J$.
    Set $A_J=AT_J$ and $G=T_JA_J^{-1}$. Then $AG=I_m$.
    \item If $0<r<\min(m,n)$, choose $r$ rows $I$ and $r$ columns
    $J$ such that $M=A_{I,J}$ is invertible. Set
    $G=T_JM^{-1}S_I$. Put $C=AT_J$ and $R=S_IA$. The columns of $C$
    form a basis for the column space of $A$. Since $S_IC=M$, the
    coordinates of each column in this basis are given by applying
    $M^{-1}S_I$. Hence $A=CM^{-1}R$, and
    \[
    AGA=AT_JM^{-1}S_IA=CM^{-1}R=A.
    \]
\end{enumerate}
In every case, $G$ is a generalized inverse of $A$.
\end{example_box}
\begin{claim_box}
\textit{Proof:} Let $G$ be chosen as in the corresponding case. If
$r=0$, then $A=0$, so $AGA=0=A$. If $r=n$, then
\[
GA=A_I^{-1}S_IA=A_I^{-1}A_I=I_n,
\]
and hence $AGA=A$. If $r=m<n$, then
\[
AG=AT_JA_J^{-1}=A_JA_J^{-1}=I_m,
\]
so again $AGA=A$.

Finally, suppose $0<r<\min(m,n)$. Put $C=AT_J$ and $R=S_IA$.
The columns of $C$ form a basis for the column space of $A$. Restricting
each column to rows $I$ shows that its coordinate vector in this basis
is $M^{-1}$ times the corresponding column of $R$. Hence
$A=CM^{-1}R$. Since $AT_J=C$ and $S_IA=R$, we get
\[
AGA=AT_JM^{-1}S_IA=CM^{-1}R=A.
\]
\hfill$\square$
\end{claim_box}

\begin{claim_box}
\textbf{Properties:} If $G$ is a generalized inverse of $A$, then
\[
(AG)^2=A(GA)G=(AGA)G=AG,\qquad
(GA)^2=G(AG)A=G(AGA)=GA.
\]
Moreover, $AG$ is the identity on $\operatorname{range}A$, and
$\operatorname{range}(AG)=\operatorname{range}A$. Hence $AG$ is a
projector onto $\operatorname{range}A$ (not necessarily an orthogonal
projector). The operator $GA$ is also a projector, onto its own range.
If $A$ is square and invertible, its only generalized inverse is
$A^{-1}$, since multiplying $AGA=A$ on the left and right by $A^{-1}$
gives $G=A^{-1}$. A left inverse $L$ or a right inverse $R$ of $A$ is a
generalized inverse, since $LA=I$ implies $ALA=A$ and $AR=I$
implies $ARA=A$.
\end{claim_box}

\begin{problem_box}
\textbf{Proposition:} Let $A\in\mathbb{F}^{m\times n}$ and
$G\in\mathbb{F}^{n\times m}$. Let $X\in\mathbb{F}^{m\times m}$ and
$Y\in\mathbb{F}^{n\times n}$ be invertible. Then $G$ is a
generalized inverse of $A$ if and only if
\[
Y^{-1}GX^{-1}
\]
is a generalized inverse of $XAY$.
\end{problem_box}

\begin{claim_box}
\textit{Proof:} By definition, the required condition is
\[
(XAY)(Y^{-1}GX^{-1})(XAY)=XAY.
\]
Using associativity and $X^{-1}X=I_m$, $YY^{-1}=I_n$, the left-hand
side simplifies to $X(AGA)Y$. Since $X$ and $Y$ are invertible,
\[
X(AGA)Y=XAY
\quad\Longleftrightarrow\quad
AGA=A,
\]
which is exactly the condition that $G$ be a generalized inverse of
$A$. \hfill$\square$
\end{claim_box}

\begin{example_box}
\textbf{Example (nonuniqueness):} For
\[
A=\begin{pmatrix}1&0\\0&0\end{pmatrix},
\]
every matrix
\[
G=\begin{pmatrix}1&a\\b&c\end{pmatrix},\qquad a,b,c\in\mathbb{F},
\]
is a generalized inverse, because $AGA=A$. Thus, unlike an ordinary
inverse, a generalized inverse need not be unique.
\end{example_box}

\subsection{Hermite canonical form}

\begin{definition_box}
\textbf{Definition:} A matrix $H=(h_{ij})\in\mathbb{F}^{n\times n}$
is in \emph{Hermite canonical form} (HCF) if:
\begin{enumerate}
    \item $H$ is upper triangular;
    \item $h_{ii}\in\{0,1\}$ for every $i$;
    \item if $h_{ii}=0$, then row $i$ of $H$ is zero;
    \item if $h_{ii}=1$, then column $i$ of $H$ is the standard basis
    vector $e_i$.
\end{enumerate}
\end{definition_box}

\begin{problem_box}
\textbf{Proposition:} If $H=(h_{ij})\in\mathbb{F}^{n\times n}$ is in
HCF, then
\[
\operatorname{rank}H=\#\{i:h_{ii}\neq0\}.
\]
\end{problem_box}

\begin{claim_box}
\textit{Proof:} Let $I=\{i:h_{ii}=1\}$. By the HCF conditions,
every row indexed outside $I$ is zero. The rows indexed by $I$ are
linearly independent: in any linear relation among them, looking at
column $i$ (which is $e_i$) forces the coefficient of row $i$ to be
zero, for each $i\in I$. Thus the rank is exactly $|I|$, the number
of nonzero diagonal entries. \hfill$\square$
\end{claim_box}

\begin{problem_box}
\textbf{Proposition:} The HCF and reduced row-echelon form are related
by row permutations as follows:
\begin{enumerate}
    \item If $H$ is in HCF, moving all zero rows to the bottom gives a
    matrix in reduced row-echelon form.
    \item If $R$ is in reduced row-echelon form with pivot columns
    $p_1<\cdots<p_r$, moving its $j$th nonzero row to row $p_j$ and
    placing the zero rows in the remaining positions gives a matrix in
    HCF.
\end{enumerate}
\end{problem_box}

\begin{claim_box}
\textit{Proof:}
\begin{enumerate}
\item Write
$I=\{i_1<\cdots<i_r:h_{i_j i_j}=1\}$. The rows outside $I$ are zero.
Move them to the bottom, preserving the order of the other rows. In
each remaining row $i_j$, upper triangularity makes every entry before
column $i_j$ zero, and $h_{i_j i_j}=1$. Moreover, column $i_j$ is
$e_{i_j}$ before reordering, so after reordering it has a $1$ in the
row corresponding to $i_j$ and zeros in every other row. Thus the
leading $1$s occur in increasing columns $i_1,\ldots,i_r$ and are the
only nonzero entries in their columns. This is reduced row-echelon
form.
\hfill$\square$
\item In $R$ the pivot columns are $e_1,\ldots,e_r$. Permute the
rows so that nonzero row $j$ is in position $p_j$, filling the other
positions with the zero rows. Row $p_j$ has zeros in all columns less
than $p_j$ and a $1$ in column $p_j$, so the resulting matrix is upper
triangular with diagonal entries $1$ at the positions $p_j$ and $0$
elsewhere. Its rows at zero-diagonal positions are zero, and column
$p_j$ is $e_{p_j}$ after the same row permutation. Hence the resulting
matrix is in HCF.
\hfill$\square$
\end{enumerate}
\end{claim_box}

\begin{problem_box}
\textbf{Proposition:} If $H\in\mathbb{F}^{n\times n}$ is in HCF and
$r=\operatorname{rank}H$, then there is a permutation matrix $P$ and a
matrix $C\in\mathbb{F}^{r\times(n-r)}$ such that
\[
PHP^{-1}=\begin{pmatrix}I_r&C\\0&0\end{pmatrix}.
\]
In general, the block $C$ need not be zero.
\end{problem_box}

\begin{claim_box}
\textit{Proof:} Let $I=\{i:h_{ii}=1\}$. By the preceding proposition,
$|I|=r$. Choose a permutation that moves the indices in $I$, in their
original order, to the first $r$ positions. The corresponding
permutation matrix conjugates $H$ by applying this same reordering to
its rows and columns. Each column indexed by $I$ is its corresponding
standard basis vector, so the leading $r$ columns of $PHP^{-1}$ form
$\begin{pmatrix}I_r\\0\end{pmatrix}$. Every row indexed outside $I$
is zero, so the remaining rows of $PHP^{-1}$ are zero. These facts give
the stated block form.

The block $C$ cannot in general be removed by permutation similarity.
For example, $\begin{pmatrix}1&1\\0&0\end{pmatrix}$ is in HCF, but
neither it nor its conjugate by the nonidentity permutation matrix is
$\begin{pmatrix}1&0\\0&0\end{pmatrix}$.
\hfill$\square$
\end{claim_box}

\begin{problem_box}
\textbf{Theorem:} For every $A\in\mathbb{F}^{n\times n}$, there
exists $B\in GL_n(\mathbb{F})$ such that $H=BA$ is in HCF, where
$GL_n(\mathbb{F})$ denotes the group of invertible $n\times n$ matrices
over $\mathbb{F}$. Moreover, this HCF is uniquely determined by the
row space of $A$, although $B$ need not be unique.
\end{problem_box}

\begin{claim_box}
\textit{Proof:} Reduce $A$ to its reduced row echelon form $R$. Suppose
$R$ has rank $r$ and its leading $1$s lie in columns
$p_1<\cdots<p_r$. In $R$, these pivot columns are
$e_1,\ldots,e_r$. Permute the rows so that row $i$ of $R$ is moved to
row $p_i$, and the zero rows occupy the remaining positions. The
resulting matrix $H=BA$, for an invertible product $B$ of row-operation
matrices, is upper triangular: row $p_i$ has zeros in every column
before its leading entry $p_i$. Its diagonal entries are $1$ in the
pivot positions and $0$ elsewhere, its zero-diagonal rows are zero,
and each pivot column $p_i$ is $e_{p_i}$. Thus $H$ is in HCF.

For uniqueness, starting with any HCF, move the rows whose diagonal
entries are $1$ to the top, preserving their order, and move its zero
rows to the bottom. The resulting matrix is in RREF. Since the RREF
is uniquely determined by the row space, the pivot positions and rows
of the HCF are uniquely determined as well.
\hfill$\square$
\end{claim_box}

\begin{claim_box}
\textbf{Property:} Every matrix in HCF is idempotent:
\[
H^2=H.
\]
\textit{Proof:} If $h_{ii}=0$, row $i$ of $H$ is zero, so row $i$ of
$H^2$ is also zero. If $h_{ii}=1$, then by upper triangularity
$h_{ik}=0$ for $k<i$. For $k>i$, if $h_{ik}\neq0$, then $h_{kk}$
cannot be $1$, since column $k$ would have to equal $e_k$. Hence row
$k$ is zero. Therefore row $i$ of $H^2$ is just row $i$ of $H$,
proving $H^2=H$. \hfill$\square$
\end{claim_box}

\begin{example_box}
\textbf{Example:} The matrix
\[
H=\begin{pmatrix}
1&2&0&4\\
0&0&0&0\\
0&0&1&-8\\
0&0&0&0
\end{pmatrix}
\]
is in HCF. Its diagonal entries are $1,0,1,0$; rows $2$ and $4$ are
zero, and columns $1$ and $3$ are $e_1$ and $e_3$, respectively.
\end{example_box}

\subsection{LDU decomposition}

\begin{definition_box}
\textbf{Definition:} A matrix $U=(u_{ij})\in\mathbb{F}^{n\times n}$
is \emph{upper unitriangular} if
\[
u_{ij}=0\quad\text{when }i>j,
\qquad
u_{ii}=1\quad\text{for every }i.
\]
It is \emph{lower unitriangular} if
\[
u_{ij}=0\quad\text{when }i<j,
\qquad
u_{ii}=1\quad\text{for every }i.
\]
We denote the sets of these matrices by $\mathcal{U}_n(\mathbb{F})$
and $\mathcal{L}_n(\mathbb{F})$, respectively.
\end{definition_box}

\begin{problem_box}
\textbf{Proposition:} The following properties hold:
\begin{enumerate}
    \item The product of two upper (respectively, lower) unitriangular
    matrices is upper (respectively, lower) unitriangular.
    \item The inverse of an upper (respectively, lower) unitriangular
    matrix is upper (respectively, lower) unitriangular.
    \item Every unitriangular matrix is invertible. In particular,
    $\mathcal{U}_n(\mathbb{F})$ and $\mathcal{L}_n(\mathbb{F})$ are
    subgroups of $GL_n(\mathbb{F})$.
    \item $\mathcal{U}_n(\mathbb{F})\cap\mathcal{L}_n(\mathbb{F})=\{I_n\}$.
\end{enumerate}
\end{problem_box}

\begin{claim_box}
\textit{Proof:}
\begin{enumerate}
\item Let $U,V$ be upper unitriangular. For $i>j$, a
nonzero summand in $(UV)_{ij}=\sum_k u_{ik}v_{kj}$ would require both
$i\leq k$ and $k\leq j$, which is impossible. For $i=j$, the only
possibly nonzero summand is the one with $k=i$, and it equals $1$.
Thus $UV$ is upper unitriangular. The lower triangular case follows
similarly, with the inequalities reversed.
\hfill$\square$
\item If $U=I_n+N$ is upper unitriangular, then $N$ is strictly upper
triangular. A nonzero entry of $N^k$ requires a chain of indices
$i_0<i_1<\cdots<i_k$, so $N^n=0$. Hence
\[
U^{-1}=(I_n+N)^{-1}=I_n-N+N^2-\cdots+(-1)^{n-1}N^{n-1},
\]
which is upper unitriangular. The same argument applies to a lower
unitriangular matrix.
\hfill$\square$
\item The identity, product, and inverse properties show that the
upper and lower unitriangular matrices are subgroups of
$GL_n(\mathbb{F})$.
\hfill$\square$
\item Finally, a matrix that is both upper and lower triangular has all
off-diagonal entries zero. If it is unitriangular, its diagonal entries
are all $1$, so it is $I_n$.
\hfill$\square$
\end{enumerate}
\end{claim_box}

\begin{problem_box}
\textbf{Proposition:} Every invertible upper triangular matrix
$A\in\mathbb{F}^{n\times n}$ has a unique factorization
\[
A=DU,
\]
where $D$ is an invertible diagonal matrix and
$U\in\mathcal{U}_n(\mathbb{F})$.
\end{problem_box}

\begin{claim_box}
\textit{Proof:} Write $A=(a_{ij})$ and set
$D=\operatorname{diag}(a_{11},\ldots,a_{nn})$. Each $a_{ii}$ is
nonzero: otherwise, take the least $k$ with $a_{kk}=0$, set $x_k=1$
and $x_i=0$ for $i>k$, and choose $x_{k-1},\ldots,x_1$ successively
so that the first $k-1$ entries of $Ax$ vanish. The entries from row
$k$ onward also vanish because $A$ is upper triangular, giving a
nonzero vector in the kernel of $A$, a contradiction. Thus $D$ is
invertible. Define $U=D^{-1}A$. Multiplication on the left
by the diagonal matrix $D^{-1}$ rescales the rows of $A$, so $U$ is
upper triangular, and
\[
u_{ii}=a_{ii}^{-1}a_{ii}=1.
\]
Thus $U$ is upper unitriangular and $A=DU$.

For uniqueness, suppose $A=D_1U_1$ with $D_1$ diagonal and $U_1$
upper unitriangular. The diagonal entries of $D_1U_1$ are exactly the
diagonal entries of $D_1$, so necessarily
$D_1=\operatorname{diag}(a_{11},\ldots,a_{nn})=D$. It follows that
$U_1=D^{-1}A=U$. Hence both factors are unique. \hfill$\square$
\end{claim_box}

\begin{definition_box}
\textbf{Definition:} An \emph{LDU decomposition} of
$A\in\mathbb{F}^{n\times n}$ is a factorization
\[
A=LDU,
\]
where $L\in\mathcal{L}_n(\mathbb{F})$, $D$ is diagonal, and
$U\in\mathcal{U}_n(\mathbb{F})$.
\end{definition_box}

\begin{problem_box}
\textbf{Proposition:} A matrix $A\in\mathbb{F}^{n\times n}$ is
upper unitriangular if and only if it is a product of elementary
upper triangular matrices of the form
\[
E_{ij}(c)=I_n+c e_i e_j^{\mathsf T},\qquad i<j,\quad c\in\mathbb{F}.
\]
These are precisely the elementary matrices for row additions
$R_i\leftarrow R_i+cR_j$ with $i<j$.
\end{problem_box}

\begin{claim_box}
\textit{Proof:} Each $E_{ij}(c)$ is upper unitriangular, and products
of upper unitriangular matrices are upper unitriangular. Conversely,
let $A$ be upper unitriangular. Starting with $j=n$ and proceeding
downward to $j=2$, clear the entries above the diagonal in column $j$
by row additions $R_i\leftarrow R_i-a_{ij}R_j$ for $i<j$. At the
stage for column $j$, row $j$ has already been cleared in every column
greater than $j$, so these operations clear column $j$ without
changing columns already treated. The resulting matrix is $I_n$.
Thus a product of elementary upper triangular matrices sends $A$ to
$I_n$, and taking inverses (with $E_{ij}(c)^{-1}=E_{ij}(-c)$) expresses
$A$ as a product of matrices of the stated form. \hfill$\square$
\end{claim_box}

\begin{problem_box}
\textbf{Theorem:} Let $A\in\mathbb{F}^{n\times n}$. For each
$k=1,\ldots,n$, let $A_k$ be its leading $k\times k$ submatrix
(the submatrix in the first $k$ rows and columns). If every $A_k$ is
invertible, then $A$ has a unique LDU decomposition
\[
A=LDU,
\]
with $L\in\mathcal{L}_n(\mathbb{F})$, $U\in\mathcal{U}_n(\mathbb{F})$,
and $D$ an invertible diagonal matrix.
\end{problem_box}

\begin{claim_box}
\textit{Proof:} We use induction on $n$. The assertion is immediate
for $n=1$. Write
\[
A=\begin{pmatrix}a&r\\c&B\end{pmatrix},
\qquad
S=B-ca^{-1}r.
\]
The hypothesis gives $a\neq0$. For every $k<n$, write the leading
$(k+1)\times(k+1)$ submatrix of $A$ in block form
\[
\begin{pmatrix}a&r_k\\c_k&B_k\end{pmatrix}.
\]
Subtracting $c_k/a$ times the first row from the remaining rows is an
invertible row operation and changes this block matrix to
\[
\begin{pmatrix}a&r_k\\0&B_k-c_ka^{-1}r_k\end{pmatrix}.
\]
The resulting matrix is invertible exactly when its lower-right block
is invertible; that block is the leading $k\times k$ submatrix of
$S$. Hence every leading square submatrix of $S$ is invertible. By
induction
$S=L_sD_sU_s$ with unitriangular $L_s,U_s$ and invertible diagonal
$D_s$. Set
\[
L=\begin{pmatrix}1&0\\ca^{-1}&L_s\end{pmatrix},\qquad
D=\begin{pmatrix}a&0\\0&D_s\end{pmatrix},\qquad
U=\begin{pmatrix}1&a^{-1}r\\0&U_s\end{pmatrix}.
\]
These have the required forms, and their product has blocks
$a,r,c,ca^{-1}r+S$, hence equals $A$.

For uniqueness, in any such factorization the $(1,1)$ entry gives
$d_{11}=a$. The first column of $A$ then determines the entries below
$1$ in the first column of $L$, and the first row determines the
entries right of $1$ in the first row of $U$. The remaining factors
give an LDU decomposition of the uniquely determined Schur complement
$S$. Induction proves uniqueness. \hfill$\square$
\end{claim_box}

\begin{problem_box}
\textbf{Theorem (Bruhat factorization):} A matrix
$A\in\mathbb{F}^{n\times n}$ is invertible if and only if it can be
written
\[
A=LPDU,
\]
where $L\in\mathcal{L}_n(\mathbb{F})$, $U\in\mathcal{U}_n(\mathbb{F})$,
$P$ is a permutation matrix, and $D$ is an invertible diagonal
matrix.
\end{problem_box}

\begin{claim_box}
\textit{Proof:} If $A=LPDU$ with the stated factors, then all four
factors are invertible, so $A$ is invertible.

Conversely, use induction by Gaussian elimination with triangular
row and column operations. In the first row of an invertible matrix,
let $q$ be the first column with a nonzero entry. Upper triangular
column additions clear the entries to the right of column $q$ in that
row; lower triangular row additions then clear the entries below the
pivot in column $q$. Scale the pivot to $1$. The pivot row and column
are now those of a permutation matrix, and the remaining
$(n-1)\times(n-1)$ block is invertible. Repeating the process on that
block, using only operations among its rows and its remaining columns,
reduces the matrix to a permutation matrix. Thus there are invertible
lower triangular $B_-$ and upper triangular $B_+$ such that
$A=B_-PB_+$ for some permutation matrix $P$.

Write $B_-=L D_-$ and $B_+=D_+U$, where $L,U$ are unitriangular and
$D_-,D_+$ are invertible diagonal matrices (normalize the diagonal
entries of each triangular factor). Since
$D_-P=P(P^{-1}D_-P)$ and conjugation by a permutation matrix preserves
diagonal matrices,
\[
A=L P\bigl((P^{-1}D_-P)D_+\bigr)U,
\]
which is the required form. \hfill$\square$
\end{claim_box}

\begin{problem_box}
\textbf{Theorem:} Let $A\in\mathbb{F}^{m\times n}$
with $m<n$. Suppose its leading $k\times k$ submatrix is invertible
for every $k=1,\ldots,m$. Then there is a factorization
\[
A=LDU,
\]
where $L\in\mathcal{L}_m(\mathbb{F})$, $D\in\mathbb{F}^{m\times n}$
is rectangular diagonal with nonzero entries in positions
$(1,1),\ldots,(m,m)$, and $U\in\mathcal{U}_n(\mathbb{F})$.
The factors are unique with these requirements. Here “leading”
submatrices means the first $k$ rows and first $k$ columns; this
specifies the condition meant by “primary submatrices” for a
rectangular matrix.
\end{problem_box}

\begin{claim_box}
\textit{Proof:} Write $A=(B\ C)$, where $B$ consists of the first
$m$ columns. The hypothesis and the square theorem give the unique
factorization $B=L D_0 U_0$, with $L,D_0,U_0$ of the corresponding
unit lower triangular, invertible diagonal, and unit upper triangular
forms. Define
\[
D=(D_0\ \ 0),\qquad
U=\begin{pmatrix}U_0&(LD_0)^{-1}C\\0&I_{n-m}\end{pmatrix}.
\]
The matrix $U$ is upper unitriangular, since its upper-right block is
unrestricted. Direct multiplication gives
\[
LDU=(LD_0U_0\ \ LD_0(LD_0)^{-1}C)=(B\ C)=A.
\]
For uniqueness, the first $m$ columns give the unique square
factorization of $B$, determining $L,D_0,U_0$. The remaining columns
then force the upper-right block of $U$ to be $(LD_0)^{-1}C$.
\hfill$\square$
\end{claim_box}

\end{document}
